%---------------------------------------------------------------------------------------------------------------------------- 
\graphicspath{{Parti/P03-Meccanica_SR/C06-Analisi_operativa_dei_SCR/Figure/}}
\chapter{Analisi operativa dei SCR }\label{cap:classificazione_sol_SCR} % \label{cap:analisi_operativa_SCR }
%---------------------------------------------------------------------------------------------------------------------------- 
\begin{sommario}
	Il capitolo ha carattere operativo: obiettivo è fornire
	strategie pratiche di soluzione per sistemi di corpi rigidi
	piani di qualsiasi tipo, senza costruire esplicitamente la
	matrice $\bm{A}$. I problemi cinematico e statico sono
	trattati in parallelo, così da rendere visibile il dualismo
	$\bm{B} = \bm{A}^{\top}$ ad ogni passo; il principio dei lavori virtuali interviene esplicitamente nelle condizioni di compatibilità per i sistemi non determinati. La classificazione
	del sistema --- isostatico, labile o iperstatico, incluso il
	caso degenere --- è oggetto del \CAPref{cap:classificazione_SCR};
	qui si assume che il tipo sia già noto e ci si concentra
	sulle procedure di soluzione.
	
	Per i sistemi determinati si presentano il metodo
	grafo-analitico per il problema cinematico e il metodo dei
	corpi liberi per quello statico. Per i sistemi labili e
	iperstatici si sviluppano le procedure operative duali ---
	vincoli fittizi e vincoli ridondanti --- e le relative
	condizioni di compatibilità, incluso il caso simultaneamente
	labile e iperstatico. Si tratta infine la gerarchia
	strutturale, che consente di decomporre sistemi complessi
	in sottosistemi risolvibili in cascata sfruttando la
	struttura a blocchi triangolare della matrice $\bm{A}$.
\end{sommario}


%---------------------------------------------------------------------------------------------------------------------------- 
\section{I due problemi in forma matriciale}
\label{sec:formulazione}
%---------------------------------------------------------------------------------------------------------------------------- 

Nel \CAPref{cap:problemi_cin_stat_SCR} si è mostrato come i
problemi cinematico e statico di un sistema di $n_c$ corpi
rigidi piani si riducano a due sistemi lineari duali:
\begin{equation}\label{eq:due_problemi}
	\bm{A}\,\bm{u} = \bm{s}
	\qquad\text{(cinematico)},
	\qquad
	\bm{A}^{\top}\bm{r} = -\bm{f}
	\qquad\text{(statico)}.
\end{equation}
La matrice $\bm{A}\in\mathbb{R}^{m\times n}$, con $n=3n_c$ e
$m$ numero di vincoli scalari, è determinata dalla sola
geometria del sistema: ogni riga corrisponde a un vincolo,
ogni blocco di tre colonne a un corpo. Il vettore $\bm{s}$
raccoglie i cedimenti vincolari; $\bm{f}$ raccoglie le forze
generalizzate attive.

Le strategie di soluzione sviluppate nelle sezioni seguenti
operano direttamente sulla geometria del sistema, senza
costruire $\bm{A}$ esplicitamente. Il loro fondamento
algebrico è interamente contenuto
nel \CAPref{cap:problemi_cin_stat_SCR}.

%---------------------------------------------------------------------------------------------------------------------------- 
\section{Sistemi cinematicamente e staticamente determinati}
\label{sec:determinato}
%---------------------------------------------------------------------------------------------------------------------------- 

Un sistema è \emph{isostatico}\index{isostaticita@isostaticità} quando $\bm{A}$ è quadrata e
invertibile ($n = m$, $\rg(\bm{A}) = n$): entrambi i problemi
hanno soluzione unica per ogni dato, senza condizioni di
compatibilità né parametri liberi. La condizione necessaria è
$\nu = n - m = 0$; se i vincoli sono geometricamente non
degeneri, essa è anche sufficiente.

La soluzione unica dei due problemi è:
\begin{equation}
	\bm{u} = \bm{A}^{-1}\bm{s},
	\qquad
	\bm{r} = -(\bm{A}^{\top})^{-1}\bm{f}.
\end{equation}


\subsection{Problema cinematico: metodo grafo-analitico}
\label{sec:grafo_analitico}

Il \emph{metodo grafo-analitico}\index{metodo!grafo-analitico|textbf} sfrutta il fatto che ogni
vincolo semplice restringe il \CRa\ del corpo a giacere su una retta
del piano (\CAPref{cap:cinematica_infinitesima}):
l'intersezione delle rette imposte dai vincoli attivi
individua il \CRa\ del corpo, e da esso si ricava l'ampiezza
della rotazione. Poiché il problema è lineare, cedimenti
diversi producono configurazioni che si sommano per
sovrapposizione.

\paragraph{Procedura.}
Dato il sistema con $n_c$ corpi e $m = 3n_c$ vincoli,
assegnati i cedimenti $\bm{s}$:

\begin{enumerate}[label=\textbf{(\arabic*)}, leftmargin=*,
	itemsep=4pt]
	
	\item \textbf{Un cedimento alla volta.}
	Si risolve il problema per $\bm{s} = s_k\,\bm{e}_k$
	(cedimento al solo vincolo $k$, tutti gli altri nulli) e si
	sommano i contributi. La sovrapposizione è lecita perché il
	problema è lineare.
	
	\item \textbf{Soppressione del vincolo cedente.}
	Si sopprime il vincolo $k$ mantenendo tutti gli altri
	perfetti. Il sistema, che era determinato, diventa labile di
	grado $g_\ell = 1$ e ammette un unico meccanismo a meno di
	un fattore scalare.
	
	\item \textbf{Localizzazione dei centri di rotazione.}
	Si determinano i centri di rotazione assoluti dei corpi dai vincoli rimasti: ogni
	vincolo semplice attivo restringe il \CRa\ del proprio corpo a una
	retta; l'intersezione di due rette non parallele individua il
	centro. Per sistemi a più corpi si applica il teorema di
	allineamento (\CAPref{cap:cinematica_infinitesima}): il \CRr\
	$C_{ij}$ di $\mathcal{C}_j$ rispetto a $\mathcal{C}_i$ giace
	sulla retta congiungente i rispettivi centri di rotazione assoluti $C_i$ e $C_j$.
	
	\item \textbf{Calcolo dell'ampiezza di rotazione.}
	L'ampiezza si ricava imponendo la prestazione cinematica del
	vincolo soppresso. Si distinguono due casi.
	
	\smallskip
	\noindent\textit{Vincolo semplice esterno} (tra il corpo
	$\mathcal{C}_i$ e il suolo, con asse $\ab_k$ nel punto
	$A_k$). La prestazione è $\ub_i(A_k)\cdot\ab_k = s_k$.
	Poiché $A_k$ ruota attorno al \CRa\ $C_i$ del corpo, si ha
	$\ub_i(A_k) = \vartheta_i\,(\ab_z\times\overrightarrow{C_i
		A_k})$, e quindi:
	\begin{equation}\label{eq:ampiezza_esterno}
		\vartheta_i =
		\frac{s_k}
		{\bigl(\overrightarrow{C_i A_k}\times\ab_k\bigr)
			\cdot\ab_z},
	\end{equation}
	dove il prodotto misto al denominatore è il braccio con
	segno di $\ab_k$ rispetto a $C_i$.
	
	\smallskip
	\noindent\textit{Vincolo semplice interno} (tra $\mathcal{C}_i$ e
	$\mathcal{C}_j$, con asse $\ab_k$ nel punto $A_k$). La
	prestazione coinvolge lo spostamento relativo:
	$[\ub_j(A_k)-\ub_i(A_k)]\cdot\ab_k = s_k$.
	Lo spostamento relativo di $A_k$ è quello di un punto che
	ruota attorno al \CRr\ $C_{ij}$ con rotazione relativa
	$\vartheta_{ij} = \vartheta_j - \vartheta_i$; la sua
	ampiezza e il suo segno si ricavano da:
	\begin{equation}\label{eq:ampiezza_interno}
		\vartheta_{ij} =
		\frac{s_k}
		{\bigl(\overrightarrow{C_{ij}A_k}\times\ab_k\bigr)
			\cdot\ab_z}.
	\end{equation}
	Le rotazioni assolute si ricavano da $\vartheta_{ij} =
	\vartheta_j - \vartheta_i$ e dalla relazione
	(\CAPref{cap:cinematica_infinitesima}):
	\begin{equation}\label{eq:rapporto_rotazioni}
		\frac{\vartheta_j}{\vartheta_i} =
		-\frac{\overrightarrow{C_i C_{ij}}}
		{\overrightarrow{C_j C_{ij}}},
	\end{equation}
	dove i vettori sono misurati lungo la retta orientata
	$C_i \to C_j$: il segno del rapporto è negativo se
	$C_{ij}$ è interno al segmento (rotazioni discordi),
	positivo se esterno (rotazioni concordi).
	
	
	\item \textbf{Sovrapposizione.}
	Si ripete il procedimento per ogni cedimento non nullo e si
	sommano le configurazioni ottenute.
	
\end{enumerate}


\paragraph{Metodo delle rette di proiezione.}\index{metodo!delle rette di proiezione|textbf}
Per sistemi a più corpi con cedimenti orizzontali o verticali,
la propagazione delle ampiezze di rotazione si esegue
graficamente senza calcolare esplicitamente bracci e prodotti
misti.

Si tracciano due rette ausiliarie esterne alla figura: una
orizzontale $r_H$ e una verticale $r_V$. Su di esse si
riportano le \emph{tracce} dei centri di rotazione assoluti
$C_i$ e relativi $C_{ij}$:
\begin{itemize}[noitemsep]
	\item su $r_V$ si riporta la proiezione orizzontale di
	ciascun centro (distanza misurata verticalmente sulla
	retta);
	\item su $r_H$ si riporta la proiezione verticale di
	ciascun centro (distanza misurata orizzontalmente sulla
	retta).
\end{itemize}
Per il teorema di allineamento, la traccia di $C_{ij}$ su
ciascuna retta cade automaticamente nella posizione corretta
rispetto alle tracce di $C_i$ e $C_j$.

\smallskip
\noindent\textit{Cedimento verticale $s_k$ al corpo
	$\mathcal{C}_i$ nel punto $A_k$.}
Si usa la retta $r_H$: la componente verticale dello
spostamento di $A_k$ è proporzionale alla distanza
orizzontale tra $C_i$ e $A_k$. L'ampiezza della rotazione
di $\mathcal{C}_i$ è:
\begin{equation}
	\vartheta_i =
	\frac{s_k}{x_{A_k} - x_{C_i}},
\end{equation}
dove $s_k$ è il cedimento con segno (positivo verso l'alto
per $\ab_k = \ab_y$) e $x_{A_k} - x_{C_i}$ è la distanza
orizzontale orientata da $C_i$ ad $A_k$. Il segno di
$\vartheta_i$ --- e dunque il verso di rotazione di
$\mathcal{C}_i$ --- si evince direttamente dal disegno,
senza ricorrere alla formula.

La rotazione di $\mathcal{C}_j$ si ricava dalla posizione
della traccia di $C_{ij}$ su $r_H$:
\begin{equation}
	\vartheta_j = \vartheta_i\,
	\frac{x_{C_{ij}} - x_{C_i}}
	{x_{C_{ij}} - x_{C_j}},
\end{equation}
dove tutte le distanze sono orientate. Anche qui il verso
di rotazione di $\mathcal{C}_j$ si legge direttamente dal
disegno: se $C_{ij}$ è interno al segmento $C_i C_j$ le
rotazioni sono discordi, se esterno sono concordi.

\smallskip
\noindent\textit{Cedimento orizzontale $s_k$ al corpo
	$\mathcal{C}_i$ nel punto $A_k$.}
Si usa la retta $r_V$: il braccio efficace è la distanza
verticale orientata $y_{A_k} - y_{C_i}$ e si procede in
modo analogo al caso verticale.


\smallskip
\noindent\textit{Caso limite: traslazione pura.}
Se il \CRa\ $C_i$ di $\mathcal{C}_i$ è all'infinito in una
direzione, il corpo trasla nella direzione \emph{ortogonale}:
lo spostamento è uniforme su tutto il corpo e la rotazione è
nulla. Analogamente, se il \CRr\ $C_{ij}$ è all'infinito in
una direzione, i due corpi ruotano dello stesso angolo
($\vartheta_i = \vartheta_j$) e il loro moto relativo è una
traslazione pura nella direzione ortogonale. In entrambi i
casi il verso e l'ampiezza della traslazione si leggono
direttamente dal disegno, imponendo la prestazione del
vincolo cedente.


\smallskip
\noindent\textit{Cedimento relativo verticale $s_k$ al vincolo
	interno tra $\mathcal{C}_i$ e $\mathcal{C}_j$ nel punto $A_k$.}
Si usa la retta $r_H$: il braccio efficace è la distanza
orizzontale orientata $x_{A_k} - x_{C_{ij}}$ tra $C_{ij}$
e $A_k$. La rotazione relativa è:
\begin{equation}
	\vartheta_{ij} =
	\frac{s_k}{x_{A_k} - x_{C_{ij}}},
\end{equation}
e il suo verso si evince direttamente dal disegno. Le
rotazioni assolute $\vartheta_i$ e $\vartheta_j$ si ricavano
poi dalla posizione della traccia di $C_{ij}$ tra quelle di
$C_i$ e $C_j$ su $r_H$, con le stesse formule del caso
esterno.

\smallskip
\noindent\textit{Cedimento relativo orizzontale $s_k$.}
Si usa $r_V$ e si procede in modo analogo.

\begin{oss}[Cedimento interno e \CRr]
	La \eqref{eq:ampiezza_interno} ha la stessa struttura
	della \eqref{eq:ampiezza_esterno}: basta sostituire il
	\CRa\ $C_i$ con il \CRr\ $C_{ij}$ e il cedimento
	assoluto con il cedimento relativo. Il caso esterno è
	quindi un caso particolare in cui uno dei due corpi è il
	suolo ($\vartheta_0 = 0$, $C_0$ all'infinito nella
	direzione ortogonale all'asse del vincolo).
\end{oss}

\begin{oss}[Significato del metodo]
	Il metodo grafo-analitico non è solo uno strumento di
	calcolo: è il modo più diretto per visualizzare la
	configurazione variata del sistema rispetto a quella di riferimento. Comprendere come
	ruotano i corpi in risposta a un cedimento è essenziale
	per interpretare i meccanismi nei sistemi labili
	(\PARref{sec:labile}) e per costruire le linee di
	influenza nei sistemi isostatici.
\end{oss}

\begin{esempio}[Arco a tre cerniere --- problema cinematico]
	\label{ex:arco_cedimenti}\index{arco!a tre cerniere}
	Si consideri un arco a profilo circolare di raggio $R$,
	composto da due semiarchi $\mathcal{C}_1$ (sinistro) e
	$\mathcal{C}_2$ (destro), vincolati al suolo da cerniere
	in $A = (-R, 0)$ e $B = (R, 0)$ e collegati fra loro da
	una cerniera interna in chiave $D = (0, R)$ (\FIGref{fig:arco_problema}\,(a)). Si scelgono
	come poli $O_1 \equiv A$ e $O_2 \equiv B$ (\FIGref{fig:arco_problema}\,(b)). Sono assegnati
	i seguenti cedimenti:
	\begin{itemize}[noitemsep]
		\item cerniera in $A$: spostamento orizzontale
		$u(O_1) = s_1 > 0$ (verso destra) e verticale
		$v(O_1) = -s_2 < 0$ (verso il basso);
		\item cerniera interna in $D$: cedimento relativo
		orizzontale $u(D_2) - u(D_1) = s_3 > 0$
		(allontanamento) e verticale
		$v(D_2) - v(D_1) = s_4 > 0$ (innalzamento relativo
		di $D_2$ rispetto a $D_1$).
	\end{itemize}
	Determinare la configurazione del sistema.
	
	% ============================================================
	\begin{figure}[ht]
		\centering
		\begin{tikzpicture}[>=Stealth, line width=0.8pt]
			
			\def\R{2.0}      % raggio dell'arco
			\def\arr{0.55}   % lunghezza frecce di cedimento/spostamento
			\def\eps{0.13}   % offset visivo di separazione D_1, D_2
			\def\rt{0.40}    % raggio degli archi di rotazione
			
			% ============================================================
			% PANNELLO (a) — dati del problema
			% ============================================================
			\begin{scope}[xshift=0cm]
				
				% semiarco C_1 (sinistro): da A=(-R,0) a D=(0,R)
				\draw[line width=1.5pt] (-\R,0)
				arc[start angle=180, end angle=90, radius=\R];
				% semiarco C_2 (destro): da D=(0,R) a B=(R,0)
				\draw[line width=1.5pt] (0,\R)
				arc[start angle=90, end angle=0, radius=\R];
				
				% raggio R dell'arco (dal centro del semicerchio a 45° su C_2)
				\draw[->, thin] (0,0) -- ({\R*cos(45)},{\R*sin(45)})
				node[midway, above left=-2pt, font=\small] {$R$};
				\node[circle, fill=black, inner sep=0pt, minimum size=2pt] at (0,0) {};
				
				
				% cerniera in A (al suolo)
				\draw[thick] (-\R,0) -- (-\R-0.22,-0.38)
				-- (-\R+0.22,-0.38) -- cycle;
				\draw[thick] (-\R-0.32,-0.38) -- (-\R+0.32,-0.38);
				\foreach \x in {-0.24,-0.12,0,0.12,0.24}{
					\draw[thin] (-\R+\x,-0.38) -- (-\R+\x-0.10,-0.52);
				}
				% cerniera in B (al suolo)
				\draw[thick] (\R,0) -- (\R-0.22,-0.38)
				-- (\R+0.22,-0.38) -- cycle;
				\draw[thick] (\R-0.32,-0.38) -- (\R+0.32,-0.38);
				\foreach \x in {-0.24,-0.12,0,0.12,0.24}{
					\draw[thin] (\R+\x,-0.38) -- (\R+\x-0.10,-0.52);
				}
				
				% perni A, B (cerniere a terra)
				\node[circle, fill=white, draw=black, inner sep=0pt,
				minimum size=5pt, line width=0.8pt] at (-\R,0) {};
				\node[circle, fill=white, draw=black, inner sep=0pt,
				minimum size=5pt, line width=0.8pt] at (\R,0) {};
				
				
				% cerniera interna in D — perno unico
				\node[circle, fill=white, draw=black, inner sep=0pt,
				minimum size=5pt, line width=0.8pt] at (0,\R) {};
				
				% coordinate ausiliarie per ancorare le frecce di s_3, s_4
				\coordinate (D1) at (-\eps,\R-\eps*0.3);
				\coordinate (D2) at (\eps,\R+\eps*0.3);
				
				% etichette A, B
				\node[left, font=\scriptsize] at (-\R,+0.2)
				{$A$};
				\node[right, font=\scriptsize] at (\R,+0.2)
				{$B$};
				
				% etichetta D (a sinistra del cluster), D_1, D_2
				\node[font=\scriptsize] at (-0.2,\R+0.30) {$D$};
				\node[font=\scriptsize, below left=-0.5pt and 1pt] at (D1) {$D_1$};
				\node[font=\scriptsize, above right=-1pt and 1pt] at (D2) {$D_2$};
				
				
				% etichette corpi (esterne all'arco, lungo i raggi a 135° e 45°)
				\node[font=\small] at ({(\R+0.30)*cos(135)},{(\R+0.30)*sin(135)})
				{$\mathcal{C}_1$};
				\node[font=\small] at ({(\R+0.30)*cos(45)},{(\R+0.30)*sin(45)})
				{$\mathcal{C}_2$};
				
				% cedimento s_1: orizzontale (verso destra) in A
				\draw[->, thick] (-\R,0.0) -- (-\R+\arr,0.0)
				node[right, font=\scriptsize] {$s_1$};
				% retta ausiliaria da A alla coda di s_2
				\draw[gray!50, thin] (-\R,0) -- (-\R-0.55,0);
				% cedimento s_2: verticale (verso il basso) in A
				% decentrato a sinistra del perno per evitare la cerniera
				\draw[->, thick] (-\R-0.55,0) -- (-\R-0.55,-\arr)
				node[below, font=\scriptsize] {$s_2$};
				
				% cedimento relativo s_3: allontanamento orizzontale in D
				\draw[->, thick] (\eps,\R+0.05) -- ++(\arr-0.05,0)
				node[right, font=\scriptsize] {$s_3$};
				\draw[->, thick] (-\eps,\R+0.05) -- ++(-\arr+0.05,0);
				% cedimento relativo s_4: innalzamento di D_2 rispetto a D_1
				\draw[->, thick] (D2.north) -- ++(0,\arr-0.10)
				node[above, font=\scriptsize] {$s_4$};
				\draw[->, thick] (D1.south) -- ++(0,-\arr+0.10);
				
				% etichetta pannello
				\node[font=\scriptsize] at (0,-1.0) {(a)};
				
			\end{scope}
			
			% ============================================================
			% PANNELLO (b) — parametri di configurazione
			% ============================================================
			\begin{scope}[xshift=6.5cm]
				
				% semiarchi
				\draw[line width=1.5pt] (-\R,0)
				arc[start angle=180, end angle=90, radius=\R];
				\draw[line width=1.5pt] (0,\R)
				arc[start angle=90, end angle=0, radius=\R];
				
				% cerniera in A
				\draw[thick] (-\R,0) -- (-\R-0.22,-0.38)
				-- (-\R+0.22,-0.38) -- cycle;
				\draw[thick] (-\R-0.32,-0.38) -- (-\R+0.32,-0.38);
				\foreach \x in {-0.24,-0.12,0,0.12,0.24}{
					\draw[thin] (-\R+\x,-0.38) -- (-\R+\x-0.10,-0.52);
				}
				% cerniera in B
				\draw[thick] (\R,0) -- (\R-0.22,-0.38)
				-- (\R+0.22,-0.38) -- cycle;
				\draw[thick] (\R-0.32,-0.38) -- (\R+0.32,-0.38);
				\foreach \x in {-0.24,-0.12,0,0.12,0.24}{
					\draw[thin] (\R+\x,-0.38) -- (\R+\x-0.10,-0.52);
				}
				
				% perni A, B, D
				\node[circle, fill=white, draw=black, inner sep=0pt,
				minimum size=5pt, line width=0.8pt] at (-\R,0) {};
				\node[circle, fill=white, draw=black, inner sep=0pt,
				minimum size=5pt, line width=0.8pt] at (\R,0) {};
				\node[circle, fill=white, draw=black, inner sep=0pt,
				minimum size=5pt, line width=0.8pt] at (0,\R) {};
				
				% etichette punti
				\node[left=4pt, font=\scriptsize] at (-\R,-0.)
				{$A \!\equiv\! O_1$};
				\node[left=4pt, font=\scriptsize] at (\R,-0.)
				{$B \!\equiv\! O_2$};
				\node[above, font=\scriptsize] at (0,\R+0.10) {$D$};
				
				% etichette corpi (esterne all'arco, lungo i raggi a 135° e 45°)
				\node[font=\small] at ({(\R+0.30)*cos(135)},{(\R+0.30)*sin(135)})
				{$\mathcal{C}_1$};
				\node[font=\small] at ({(\R+0.30)*cos(45)},{(\R+0.30)*sin(45)})
				{$\mathcal{C}_2$};
				
				
				% parametri di C_1 (in A)
				% u(O_1) orizzontale verso destra
				\draw[->, thick] (-\R,0.0) -- (-\R+\arr+0.15,0.0)
				node[right, font=\scriptsize] {$u(O_1)$};
				% v(O_1) verticale verso l'alto (decentrato a sinistra)
				\draw[->, thick] (-\R-0.05,0) -- (-\R-0.05,\arr+0.15)
				node[left,  yshift=-5pt, font=\scriptsize] {$v(O_1)$};
				% vartheta_1 antiorario attorno a O_1 (traslato verso l'alto)
				\draw[->, thick]
				(-\R+\rt,0.10) arc[start angle=0, end angle=75, radius=\rt];
				\node[font=\scriptsize] at ({-\R+0.1+0.65*cos(45)},{0.06+0.65*sin(45)})
				{$\vartheta_1$};
				
				% parametri di C_2 (in B)
				% u(O_2) orizzontale verso destra
				\draw[->, thick] (\R,0.0) -- (\R+\arr+0.15,0.0)
				node[right, font=\scriptsize] {$u(O_2)$};
				% v(O_2) verticale verso l'alto (decentrato a sinistra)
				\draw[->, thick] (\R+0.05,0) -- (\R+0.05,\arr+0.15)
				node[left, yshift=-5pt, font=\scriptsize] {$v(O_2)$};
				% vartheta_2 antiorario attorno a O_2 (traslato verso l'alto)
				\draw[->, thick]
				(\R+\rt,0.10) arc[start angle=0, end angle=75, radius=\rt];
				\node[font=\scriptsize] at ({\R+0.1+0.65*cos(45)},{0.06+0.65*sin(45)})
				{$\vartheta_2$};
				
				% etichetta pannello
				\node[font=\scriptsize] at (0,-1.0) {(b)};
				
			\end{scope}
			
		\end{tikzpicture}
		\caption{Arco a tre cerniere
			(\ESEref{ex:arco_cedimenti}):
			\emph{(a)}~dati del problema cinematico ---
			cedimenti in $A$ ($s_1$ orizzontale, $s_2$
			verticale) e cedimenti relativi in $D$
			($s_3$ orizzontale, $s_4$ verticale); le
			frecce indicano il verso positivo dei
			cedimenti assegnati.
			\emph{(b)}~parametri di configurazione ---
			spostamenti $u(O_1)$, $v(O_1)$ e rotazione
			$\vartheta_1$ del polo $O_1\equiv A$;
			spostamenti $u(O_2)$, $v(O_2)$ e rotazione
			$\vartheta_2$ del polo $O_2\equiv B$; le
			frecce indicano il verso positivo delle
			incognite.}
		\label{fig:arco_problema}
	\end{figure}
	% ============================================================
	
	\begin{svolgimento}
		
		
		\esottotit{Formulazione analitica.}
		Le incognite sono i parametri di configurazione dei due
		corpi riferiti ai poli scelti (\FIGref{fig:arco_problema}\,b):
		\begin{equation*}
			\bm{u} =
			\bigl\{u(O_1)\; v(O_1)\; \vartheta_1\,\;
			u(O_2)\; v(O_2)\; \vartheta_2
			\bigr\}^\top.
		\end{equation*}
		Le equazioni di vincolo --- cerniera in $A$ (2 equazioni),
		cerniera in $B$ (2 equazioni), cerniera interna in $D$
		(2 equazioni) --- conducono direttamente al sistema
		cinematico \eqref{eq:due_problemi}, $\bm{A}\bm{u} = \bm{s}$:
		\begin{equation*}
			\begin{bmatrix}
				1 & 0 &  0 & 0 & 0 &  0 \\
				0 & 1 &  0 & 0 & 0 &  0 \\
				0 & 0 &  0 & 1 & 0 &  0 \\
				0 & 0 &  0 & 0 & 1 &  0 \\
				-1 & 0 & R & 1 & 0 &  -R \\
				0 & -1 &  R & 0 & 1 & - R
			\end{bmatrix}
			\begin{Bmatrix}
				u(O_1) \\ v(O_1) \\ \vartheta_1 \\
				u(O_2) \\ v(O_2) \\ \vartheta_2
			\end{Bmatrix}
			=
			\begin{Bmatrix}
				s_1 \\ -s_2 \\ 0 \\ 0 \\ s_3 \\ s_4
			\end{Bmatrix}.
		\end{equation*}
		Le prime quattro equazioni danno immediatamente:
		\begin{equation*}
			u(O_1) = s_1, \quad
			v(O_1) = -s_2, \quad
			u(O_2) = 0, \quad
			v(O_2) = 0.
		\end{equation*}
		Le ultime due, sostituite le incognite note, diventano:
		\begin{equation*}
			\begin{bmatrix}
				R & -R \\
				-R & -R
			\end{bmatrix}
			\begin{Bmatrix}
				\vartheta_1 \\ \vartheta_2
			\end{Bmatrix}
			=
			\begin{Bmatrix}
				s_1 + s_3 \\ s_4 - s_2
			\end{Bmatrix},
		\end{equation*}
		da cui:
		\begin{equation*}\label{eq:arco_rotazioni}
			\vartheta_1 = \frac{s_1 + s_2 + s_3 - s_4}{2R},
			\qquad
			\vartheta_2 = \frac{-s_1 + s_2 - s_3 - s_4}{2R}.
		\end{equation*}
		La configurazione è univocamente determinata, come atteso
		per un sistema isostatico.
		
		\noindent	\esottotit{Soluzione grafica per singoli cedimenti.}
		Per il principio di sovrapposizione si considerano i	quattro cedimenti separatamente. Per ciascuno si individuano i centri di rotazione assoluti $C_1$, $C_2$	e relativo $C_{12}$, si tracciano le rette $r_H$ e $r_V$ e si propagano le ampiezze con il metodo delle rette di	proiezione. I risultati sono raccolti nelle	\FIGrefS{fig:arco_cedimenti_esterni}--\ref{fig:arco_cedimenti_interni} e	riepilogati nella	\TABref{tab:arco_cedimenti}.
		
		
		% ============================================================
		\begin{figure}[ht]
			\centering
			\resizebox{\textwidth}{!}{%
				\begin{tikzpicture}[>=Stealth, line width=0.8pt]
					
					\def\R{1.4}
					\def\arr{0.55}
					\def\ss{0.65}
					\def\xV{-2.8}      % retta r_V a sinistra dell'arco
					\def\yH{-1.3}
					\def\rt{0.28}
					\def\dxp{7.0}      % separazione tra i due pannelli (da rifinire)
					\def\rA{1.0}
					
					
					% ============================================================
					% PANNELLO (a) — Cedimento s_1 (orizzontale in A)
					% ============================================================
					\begin{scope}[xshift=0cm]
						
						% --- semiarchi ---
						\draw[line width=1.2pt] (-\R,0)
						arc[start angle=180, end angle=90, radius=\R];
						\draw[line width=1.2pt] (0,\R)
						arc[start angle=90, end angle=0, radius=\R];
						
						% --- cerniera al suolo in A ---
						\draw[thick] (-\R,0) -- (-\R-0.18,-0.30)
						-- (-\R+0.18,-0.30) -- cycle;
						\draw[thick] (-\R-0.28,-0.30) -- (-\R+0.28,-0.30);
						\foreach \x in {-0.20,-0.10,0,0.10,0.20}{
							\draw[thin] (-\R+\x,-0.30) -- (-\R+\x-0.09,-0.42);
						}
						% --- cerniera al suolo in B ---
						\draw[thick] (\R,0) -- (\R-0.18,-0.30)
						-- (\R+0.18,-0.30) -- cycle;
						\draw[thick] (\R-0.28,-0.30) -- (\R+0.28,-0.30);
						\foreach \x in {-0.20,-0.10,0,0.10,0.20}{
							\draw[thin] (\R+\x,-0.30) -- (\R+\x-0.09,-0.42);
						}
						
						% --- perni in A, B (esterni) e D (interno) ---
						\node[circle, fill=white, draw=black, inner sep=0pt,
						minimum size=4pt, line width=0.7pt] at (-\R,0) {};
						\node[circle, fill=white, draw=black, inner sep=0pt,
						minimum size=4pt, line width=0.7pt] at (\R,0) {};
						\node[circle, fill=white, draw=black, inner sep=0pt,
						minimum size=4pt, line width=0.7pt] at (0,\R) {};
						
						% --- costruzione di C_1: verticale per A e retta BD ---
						\draw[gray!55, dashed, thin] (-\R,-0.50) -- (-\R, 2*\R+0.40);
						\draw[gray!55, dashed, thin]
						({\R+0.25},-0.25) -- ({-\R-0.20}, 2*\R+0.20);
						
						% --- centro C_1 (label spostata a destra: la sinistra è
						\fill (-\R, 2*\R) circle (1.8pt);
						\node[font=\scriptsize, anchor=west] at ({-\R+0.10}, 2*\R) {$C_1$};
						
						% --- etichette punti notevoli ---
						\node[font=\scriptsize, anchor=east] at ({-\R-0.10}, 0.18) {$A$};
						\node[font=\scriptsize, anchor=west] at ({\R}, 0.) {$B\!\equiv\!C_2$};
						\node[font=\scriptsize, anchor=south] at (0.0, \R){$D\!\equiv\!C_{12}$};
						
						% --- etichette dei corpi (interno della rispettiva cupola) ---
						\node[font=\scriptsize]	at ({0.95*cos(110)-0.3},{0.95*sin(110)}) {$\mathcal{C}_1$};
						\node[font=\scriptsize]	at ({0.95*cos(70)+1.0},{0.95*sin(70)})	{$\mathcal{C}_2$};
						
						% --- cedimento s_1 in A ---
						\draw[->, thick] (-\R,0) -- ({-\R+\ss},0);
						\node[font=\scriptsize, anchor=south] at ({-\R+\ss/2}, 0.04) {$s_1$};
						
						%			
						% ==========================================================
						% retta r_V (verticale, a sinistra dell'arco) — campo u(y)
						% ==========================================================
						\draw[gray!45, thin] (\xV, -0.30) -- (\xV, 2*\R+0.40)
						node[above, font=\scriptsize] {$r_V$};
						
						% proiezioni dei punti notevoli su r_V
						\draw[gray!40, dashed, thin] (-\R, 0) -- (\xV, 0);
						\draw[gray!40, dashed, thin] (0, \R) -- (\xV, \R);
						\draw[gray!40, dashed, thin] (-\R, 2*\R) -- (\xV, 2*\R);
						
						\fill[gray!75] (\xV, 0) circle (1.3pt);
						\fill[gray!75] (\xV, \R) circle (1.3pt);
						\fill[gray!75] (\xV, 2*\R) circle (1.3pt);
						
						% --- skyline u_1^(1) di C_1 (gray solido) ---
						\draw[gray!75, thin]
						({\xV+\ss}, 0) -- (\xV, 2*\R);
						\foreach \k in {0,1,2,3,4}{
							\pgfmathsetmacro{\yy}{\k*\R/4}
							\pgfmathsetmacro{\uu}{\ss*(1-\yy/(2*\R))}
							\draw[->, gray!80, thin]
							(\xV,\yy) -- ({\xV+\uu},\yy);
						}
						\node[gray!85, font=\scriptsize, anchor=north west]
						at ({\xV+\ss-0.2}, 1.0) {$u_1$};
						
						% --- skyline u_2 di C_2 (nero tratteggiato per distinguerlo) ---
						\draw[black!55, thin, dashed]
						(\xV, 0) -- ({\xV+\ss/2}, \R);
						\foreach \k in {1,2,3,4}{
							\pgfmathsetmacro{\yy}{\k*\R/4}
							\pgfmathsetmacro{\uu}{(\ss/2)*\k/4}
							\draw[->, black!60, thin, dashed]
							(\xV,\yy) -- ({\xV+\uu},\yy);
						}
						\node[black!70, font=\scriptsize, anchor=east] at ({\xV}, \R-0.7) {$u_2$};
						
						
						% --- etichette delle proiezioni dei centri su r_V ---
						\node[font=\scriptsize, anchor=north east]	at ({\xV-0.05}, {2*\R+0.2}) {$C_1$};
						\node[font=\scriptsize, anchor=north east]	at ({\xV-0.05}, {\R+0.3}) {$C_{12}$};
						\node[font=\scriptsize, anchor=north east]	at ({\xV-0.05}, 0.2) {$C_2$};
						
						
						
						% ==========================================================
						% retta r_H (orizzontale) — campo v(x)
						% ==========================================================
						\draw[gray!45, thin] (-\R-0.40, \yH) -- (\R+0.40, \yH)
						node[right, font=\scriptsize] {$r_H$};
						
						\draw[gray!40, dashed, thin] (-\R, -0.45) -- (-\R, \yH);
						\draw[gray!40, dashed, thin] (0, \R-0.10) -- (0, \yH);
						\draw[gray!40, dashed, thin] (\R, -0.45) -- (\R, \yH);
						
						\fill[gray!75] (-\R, \yH) circle (1.3pt);
						\fill[gray!75] (0, \yH) circle (1.3pt);
						\fill[gray!75] (\R, \yH) circle (1.3pt);
						
						% --- skyline v_1 di C_1 ---
						\draw[gray!75, thin]
						(-\R, \yH) -- (0, {\yH+\ss/2});
						\foreach \k in {1,2,3,4}{
							\pgfmathsetmacro{\xx}{-\R+\k*\R/4}
							\pgfmathsetmacro{\vv}{(\ss/2)*\k/4}
							\draw[->, gray!80, thin]
							(\xx,\yH) -- (\xx, {\yH+\vv});
						}
						\node[gray!85, font=\scriptsize] at ({-\R/2-0.05}, {\yH-0.30}) {$v_1$};
						
						\node[gray!85, font=\scriptsize, anchor=south] at ({\xV+\ss/2}, -0.4) {$s_1$};
						
						% --- skyline v_2 di C_2 ---
						\draw[gray!75, thin]
						(0, {\yH+\ss/2}) -- (\R, \yH);
						\foreach \k in {0,1,2,3}{
							\pgfmathsetmacro{\xx}{\k*\R/4}
							\pgfmathsetmacro{\vv}{(\ss/2)*(1-\k/4)}
							\ifnum\k>0
							\draw[->, gray!80, thin]
							(\xx,\yH) -- (\xx, {\yH+\vv});
							\fi
						}
						\node[gray!85, font=\scriptsize]
						at ({\R/2+0.05}, {\yH-0.30}) {$v_2$};
						
						\node[font=\scriptsize, anchor=south] at (0.0, {\yH+\ss/3}) {$s_1/2$};
						
						
						
						% --- etichette delle proiezioni dei centri su r_H ---
						\node[font=\scriptsize, anchor=north] at (-\R, {\yH-0.04}) {$C_1$};
						\node[font=\scriptsize, anchor=north]	at (0, {\yH-0.04}) {$C_{12}$};
						\node[font=\scriptsize, anchor=north] at (\R, {\yH-0.04}) {$C_2$};
						
						
						% Configurazione variata (qualitativa, schematica)
						\draw[gray!65, thin] (-\R, 2*\R) -- ({-\R+\ss}, 0);
						\draw[gray!65, thin] (-\R, 2*\R) -- ({\ss/2}, {\R+\ss/2});
						\draw[gray!65, thin] (\R, 0) -- ({\ss/2}, {\R+\ss/2});
						
						\fill[gray!65] ({\ss/2}, {\R+\ss/2}) circle (1.5pt);
						
						\node[gray!65,font=\scriptsize, anchor=west] at ({-\R+\ss-0.08}, 0.0) {$A'$};
						\node[gray!65,font=\scriptsize, anchor=south west]
						at ({\ss/2+0.05}, {\R+\ss/2+0.05}) {$D'$};
						
						% --- semiarchi deformati (linearizzazione del campo di rotazione) ---
						\draw[gray!65, line width=1pt]
						plot[domain=90:180, samples=30, smooth]
						({\R*cos(\x) + \ss*(1 - sin(\x)/2)},
						{\R*sin(\x) + (\ss/2)*(cos(\x) + 1)});
						\draw[gray!65, line width=1pt]
						plot[domain=0:90, samples=30, smooth]
						({\R*cos(\x) + (\ss/2)*sin(\x)},
						{\R*sin(\x) + (\ss/2)*(1 - cos(\x))});
						
						% --- archi degli angoli di rotazione ---
						\pgfmathsetmacro{\dtheta}{atan(\ss/(2*\R))}
						\def\rA{1.0}
						% vartheta_1 fra verticale per A e C_1A'  (antiorario)
						\draw[->, gray!75, thin]
						({-\R + \rA*cos(270)}, {2*\R + \rA*sin(270)})
						arc[start angle=270, end angle={270+\dtheta}, radius=\rA];
						% vartheta_1 fra C_1D (= retta BD prolungata) e C_1D'  (antiorario)
						\draw[->, gray!75, thin]
						({-\R + \rA*cos(315)}, {2*\R + \rA*sin(315)})
						arc[start angle=315, end angle={315+\dtheta}, radius=\rA];
						% vartheta_2 fra BD e BD'  (orario)
						\draw[->, gray!75, thin]
						({\R + \rA*cos(135)}, {\rA*sin(135)})
						arc[start angle=135, end angle={135-\dtheta}, radius=\rA];
						
						% --- etichette degli angoli, accanto alla punta della freccia ---
						\node[gray!85, font=\scriptsize, anchor=west]
						at ({-\R + \rA*cos(270+\dtheta) - 0.08}, {2*\R + \rA*sin(270+\dtheta)})
						{$\vartheta_1$};
						\node[gray!85, font=\scriptsize, anchor=west]
						at ({-\R + \rA*cos(315+\dtheta) - 0.08}, {2*\R + \rA*sin(315+\dtheta)+0.08})
						{$\vartheta_1$};
						\node[gray!85, font=\scriptsize, anchor=east]
						at ({\R + \rA*cos(135-\dtheta) - 0.06}, {\rA*sin(135-\dtheta)})
						{$\vartheta_2$};
						
						% --- etichetta pannello (centrata sull'estensione effettiva) ---
						\node[font=\scriptsize] at ({(\xV-\R)/2}, \yH-0.55) {(a)};
						
					\end{scope}
					
					% ============================================================
					% PANNELLO (b) — Cedimento s_2 (verticale in A)
					% ============================================================
					\begin{scope}[xshift=\dxp cm]   % posizione del pannello (b)
						
						% --- semiarchi ---
						\draw[line width=1.2pt] (-\R,0)
						arc[start angle=180, end angle=90, radius=\R];
						\draw[line width=1.2pt] (0,\R)
						arc[start angle=90, end angle=0, radius=\R];
						
						% --- cerniera al suolo in A ---
						\draw[thick] (-\R,0) -- (-\R-0.18,-0.30)
						-- (-\R+0.18,-0.30) -- cycle;
						\draw[thick] (-\R-0.28,-0.30) -- (-\R+0.28,-0.30);
						\foreach \x in {-0.20,-0.10,0,0.10,0.20}{
							\draw[thin] (-\R+\x,-0.30) -- (-\R+\x-0.09,-0.42);
						}
						% --- cerniera al suolo in B ---
						\draw[thick] (\R,0) -- (\R-0.18,-0.30)
						-- (\R+0.18,-0.30) -- cycle;
						\draw[thick] (\R-0.28,-0.30) -- (\R+0.28,-0.30);
						\foreach \x in {-0.20,-0.10,0,0.10,0.20}{
							\draw[thin] (\R+\x,-0.30) -- (\R+\x-0.09,-0.42);
						}
						
						% --- perni in A, B (esterni) e D (interno) ---
						\node[circle, fill=white, draw=black, inner sep=0pt,
						minimum size=4pt, line width=0.7pt] at (-\R,0) {};
						\node[circle, fill=white, draw=black, inner sep=0pt,
						minimum size=4pt, line width=0.7pt] at (\R,0) {};
						\node[circle, fill=white, draw=black, inner sep=0pt,
						minimum size=4pt, line width=0.7pt] at (0,\R) {};
						
						
						% --- costruzione di C_1: orizzontale per A e retta BD ---
						\draw[gray!55, dashed, thin] (-\R-0.30, 0) -- (\R+0.30, 0);
						\draw[gray!55, dashed, thin] (\R+0.25, -0.25) -- (-0.40, \R+0.40);
						
						% --- etichette punti notevoli ---
						\node[font=\scriptsize, anchor=east] at ({-\R-0.10}, 0.18) {$A$};
						\node[font=\scriptsize, anchor=west] at ({\R+0.10}, 0.18){$B\!\equiv\!C_1\!\equiv\!C_2$};
						\node[font=\scriptsize, anchor=south] at (0.0, \R){$D\!\equiv\!C_{12}$};
						
						% --- etichette dei corpi (interno della rispettiva cupola) ---
						\node[font=\scriptsize]	at ({0.95*cos(110)-1.2},{0.95*sin(110)}) {$\mathcal{C}_1$};
						\node[font=\scriptsize]	at ({0.95*cos(70)+1.20},{0.95*sin(70)})	{$\mathcal{C}_2$};
						
						% --- cedimento s_2 in A (verticale, verso il basso) ---
						\draw[->, thick] (-\R, 0) -- (-\R, -\ss);
						\node[font=\scriptsize, anchor=east] at ({-\R-0.32}, -\ss/2) {$s_2$};
						
						% --- costruzione di C_1: orizzontale per A e retta BD ---
						\draw[gray!55, dashed, thin] (-\R-0.30, 0) -- (\R+0.30, 0);
						\draw[gray!55, dashed, thin] (\R+0.25, -0.25) -- (-0.40, \R+0.40);
						
						% ============================================================
						% retta r_V (verticale, a sinistra) — campo u(y)
						% ============================================================
						\draw[gray!45, thin] (\xV, -0.30) -- (\xV, \R+0.40)	node[above, font=\scriptsize] {$r_V$};
						
						% proiezioni dei punti notevoli su r_V
						\draw[gray!40, dashed, thin] (\R, 0) -- (\xV, 0);
						\draw[gray!40, dashed, thin] (0, \R) -- (\xV, \R);
						
						\fill[gray!75] (\xV, 0) circle (1.3pt);
						\fill[gray!75] (\xV, \R) circle (1.3pt);
						
						% etichette delle tracce — a destra di r_V (skyline a sx)
						\node[font=\scriptsize, anchor=east]	at ({\xV}, 0.0) {$C_1\!\equiv\!C_2$};
						\node[font=\scriptsize, anchor=west]	at ({\xV+0.05}, {\R-0.04}) {$C_{12}$};
						
						% skyline u(y) — L-shape a sinistra di r_V
						\draw[gray!75, thin]
						(\xV, 0) -- ({\xV-\ss/2}, \R) -- (\xV, \R);
						\foreach \k in {1,2,3,4}{
							\pgfmathsetmacro{\yy}{\k*\R/4}
							\pgfmathsetmacro{\uu}{\k*\ss/8}
							\draw[->, gray!75, thin] (\xV, \yy) -- ({\xV-\uu}, \yy);
						}
						
						% picco label "s_2/2" al midpoint freccia y=R
						\node[font=\scriptsize, anchor=south] at ({\xV-\ss/4}, {\R}) {$s_2/2$};
						
						% nome del campo, sotto la base del r_V
						\node[font=\scriptsize, anchor=north east] at ({\xV-\ss/4}, {\R/2+0.4}) 	{$u_1\!=\!u_2$};
						
						% ============================================================
						% retta r_H (orizzontale, sotto) — campo v(x)
						% ============================================================
						\draw[gray!45, thin] (-\R-0.30, \yH) -- (\R+0.40, \yH)
						node[right, font=\scriptsize] {$r_H$};
						
						% proiezioni dei punti notevoli su r_H
						\draw[gray!40, dashed, thin] (-\R, 0) -- (-\R, \yH);
						\draw[gray!40, dashed, thin] (0, \R) -- (0, \yH);
						\draw[gray!40, dashed, thin] (\R, 0) -- (\R, \yH);
						
						\fill[gray!75] (-\R, \yH) circle (1.3pt);
						\fill[gray!75] (0, \yH) circle (1.3pt);
						\fill[gray!75] (\R, \yH) circle (1.3pt);
						
						% etichette delle tracce — sopra r_H (skyline sotto)
						\node[font=\scriptsize, anchor=south]	at (-\R, {\yH+0.04}) {$A$};
						\node[font=\scriptsize, anchor=south]	at (0, {\yH+0.04}) {$C_{12}$};
						\node[font=\scriptsize, anchor=south]	at (\R, {\yH+0.04}) {$C_1\!\equiv\!C_2$};
						
						% skyline v(x) — L-shape sotto r_H
						\draw[gray!75, thin]
						(-\R, \yH) -- (-\R, {\yH-\ss}) -- (\R, \yH);
						\foreach \k in {0,1,2,3,4,5,6,7}{
							\pgfmathsetmacro{\xx}{-\R + \k*\R/4}
							\pgfmathsetmacro{\vv}{\ss*(1-\k/8)}
							\draw[->, gray!75, thin] (\xx, \yH) -- (\xx, {\yH-\vv});
						}
						
						% picco label "s_2" al midpoint freccia x=-R
						\node[font=\scriptsize, anchor=east] at ({-\R+0.04}, {\yH-\ss/2}) {$s_2$};
						
						% nome del campo, sotto e a sx della picco
						\node[font=\scriptsize, anchor=north east]	at ({\R/2-0.05}, {\yH-\ss+0.2}) {$v_1\!=\!v_2$};
						
						% ============================================================
						% configurazione variata: rotazione rigida attorno a B
						% ============================================================
						
						% raggi della configurazione variata
						\draw[gray!65, thin] (\R, 0) -- (-\R, -\ss);             % BA'
						\draw[gray!65, thin] (\R, 0) -- ({-\ss/2}, {\R-\ss/2});  % BD'
						
						% semiarchi deformati (linearizzazione del campo di rotazione)
						\draw[gray!75, line width=0.9pt]
						plot[domain=90:180, samples=24, smooth, variable=\t]
						({\R*cos(\t) - \ss/2*sin(\t)}, {\R*sin(\t) + \ss/2*(cos(\t)-1)});
						\draw[gray!75, line width=0.9pt]
						plot[domain=0:90, samples=24, smooth, variable=\t]
						({\R*cos(\t) - \ss/2*sin(\t)}, {\R*sin(\t) + \ss/2*(cos(\t)-1)});
						
						% archi degli angoli vartheta_1, vartheta_2 a radius rA (entrambi
						% al medesimo centro B, stessa ampiezza, antiorario)
						\pgfmathsetmacro{\dvartheta}{atan(\ss/(2*\R))}
						
						% vartheta_1: BA → BA' (rotazione del body 1)
						\draw[->, gray!85, thin]
						({\R+\rA*cos(180)}, {\rA*sin(180)})
						arc[start angle=180, end angle={180+\dvartheta}, radius=\rA];
						\node[font=\scriptsize, gray!85, anchor=west]
						at ({\R+\rA*cos(180+\dvartheta)+0.05}, {\rA*sin(180+\dvartheta)})
						{$\vartheta_1$};
						
						% vartheta_2: BD → BD' (rotazione del body 2)
						\draw[->, gray!85, thin]
						({\R+\rA*cos(135)}, {\rA*sin(135)})
						arc[start angle=135, end angle={135+\dvartheta}, radius=\rA];
						\node[font=\scriptsize, gray!85, anchor=east]
						at ({\R+\rA*cos(135+\dvartheta)-0.05}, {\rA*sin(135+\dvartheta)})
						{$\vartheta_2$};
						
						% marker su D' (cerniera interna nella configurazione variata)
						\fill[gray!70] ({-\ss/2}, {\R-\ss/2}) circle (1.8pt);
						
						% --- etichetta pannello, sotto tutto ---
						\node[font=\scriptsize] at ({(\xV-\R)/2}, \yH-0.55) {(b)};
						
					\end{scope}
					
				\end{tikzpicture}
			}
			\caption{Soluzione grafica per i cedimenti esterni in $A$.
				\emph{(a)}~cedimento $s_1$ (orizzontale): tre centri distinti,
				$C_1$ all'intersezione fra la verticale per $A$ e la retta $BD$.
				\emph{(b)}~cedimento $s_2$ (verticale): caso degenere,
				$C_1\equiv C_2\equiv B$, e il sistema ruota rigidamente
				attorno a $B$.}
			\label{fig:arco_cedimenti_esterni}
		\end{figure}
		
		\smallskip
		\noindent\textit{Cedimento $s_1$ (orizzontale in $A$).}
		La cerniera in $B$ fissa il centro assoluto di
		$\mathcal{C}_2$: $C_2 \equiv B$. La cerniera interna in
		$D$ con cedimento nullo impone che $D$ sia punto di uguale
		spostamento per i due semiarchi: $C_{12} \equiv D$.
		Il vincolo verticale in $A$ è attivo, quindi $C_1$ giace
		sulla verticale per $A$; per il teorema di allineamento
		$C_1$ giace inoltre sulla retta per $C_2 \equiv B$ e
		$C_{12} \equiv D$. Il centro $C_1$ è dunque
		l'intersezione della verticale per $A$ con la retta
		$BD$, che cade al di sopra di $A$ a distanza $2R$
		(si veda la \FIGref{fig:arco_cedimenti_esterni}\,(a)). 	
		
		La rotazione di $\mathcal{C}_1$ si ricava dalla
		prestazione del vincolo orizzontale cedente in $A$,
		usando la retta $r_V$: il braccio efficace è la distanza
		verticale tra $C_1$ e $A$, pari a $2R$, quindi:
		\begin{equation*}
			\vartheta_1 = \frac{s_1}{2R} \qquad
			\text{(antiorario, dalla figura)}.
		\end{equation*}
		La traccia di $C_{12} \equiv D$ su $r_V$ è interna al
		segmento tra le tracce di $C_1$ e $C_2$: le rotazioni
		sono discordi e il braccio efficace di $C_2$ rispetto
		a $C_{12}$ è pari a $R$, quindi:
		\begin{equation*}
			\vartheta_2 = -\frac{s_1}{2R} \qquad
			\text{(orario, dalla figura)}.
		\end{equation*}
		
		
		
		
		% ============================================================
		\begin{figure}[ht]
			\centering
			\resizebox{\textwidth}{!}{%
				\begin{tikzpicture}[>=Stealth, line width=0.8pt]
					
					\def\R{1.4}
					\def\arr{0.55}
					\def\ss{0.65}
					\def\xV{-2.8}      % retta r_V a sinistra dell'arco
					\def\yH{-1.3}
					\def\rt{0.28}
					\def\dxp{7.0}      % separazione tra i due pannelli (da rifinire)
					
					% ============================================================
					% PANNELLO (a) — Cedimento s_3 (orizzontale in D)
					% ============================================================
					\begin{scope}[xshift=0cm]
						
						% --- semiarchi ---
						\draw[line width=1.2pt] (-\R,0)
						arc[start angle=180, end angle=90, radius=\R];
						\draw[line width=1.2pt] (0,\R)
						arc[start angle=90, end angle=0, radius=\R];
						
						% --- cerniera al suolo in A ---
						\draw[thick] (-\R,0) -- (-\R-0.18,-0.30)
						-- (-\R+0.18,-0.30) -- cycle;
						\draw[thick] (-\R-0.28,-0.30) -- (-\R+0.28,-0.30);
						\foreach \x in {-0.20,-0.10,0,0.10,0.20}{
							\draw[thin] (-\R+\x,-0.30) -- (-\R+\x-0.09,-0.42);
						}
						% --- cerniera al suolo in B ---
						\draw[thick] (\R,0) -- (\R-0.18,-0.30)
						-- (\R+0.18,-0.30) -- cycle;
						\draw[thick] (\R-0.28,-0.30) -- (\R+0.28,-0.30);
						\foreach \x in {-0.20,-0.10,0,0.10,0.20}{
							\draw[thin] (\R+\x,-0.30) -- (\R+\x-0.09,-0.42);
						}
						
						% --- costruzione di C_{12}: retta AB e verticale per D ---
						\draw[gray!55, dashed, thin] (-\R-0.30, 0) -- (\R+0.30, 0);
						\draw[gray!55, dashed, thin] (0, -0.50) -- (0, \R+0.20);
						
						% --- perni in A, B (esterni) e D (interno) ---
						\node[circle, fill=white, draw=black, inner sep=0pt,
						minimum size=4pt, line width=0.7pt] at (-\R,0) {};
						\node[circle, fill=white, draw=black, inner sep=0pt,
						minimum size=4pt, line width=0.7pt] at (\R,0) {};
						\node[circle, fill=white, draw=black, inner sep=0pt,
						minimum size=4pt, line width=0.7pt] at (0,\R) {};
						
						% --- centro C_{12} (punto medio di AB) ---
						\fill (0, 0) circle (1.8pt);
						
						% --- etichette dei corpi ---
						\node[font=\scriptsize] at ({0.95*cos(110)-1.0},{0.95*sin(110)}) {$\mathcal{C}_1$};
						\node[font=\scriptsize] at ({0.95*cos(70)+1.0},{0.95*sin(70)})  {$\mathcal{C}_2$};
						
						% --- etichette punti notevoli ---
						\node[font=\scriptsize, anchor=east]  at ({-\R-0.10}, 0.18) {$A\!\equiv\!C_1$};
						\node[font=\scriptsize, anchor=west]  at ({\R+0.10}, 0.18)  {$B\!\equiv\!C_2$};
						\node[font=\scriptsize, anchor=north] at (0.0, \R-0.10)     {$D$};
						\node[font=\scriptsize, anchor=south] at (0.05, -0.05) {$C_{12}$};
						
						% --- cedimento relativo s_3 in D (frecce divergenti orizzontali) ---
						\def\eps{0.08}
						\draw[->, thick] (\eps, \R+0.05) -- ++({\ss/2}, 0);
						\draw[->, thick] (-\eps, \R+0.05) -- ++({-\ss/2}, 0);
						\node[above, font=\scriptsize] at (0, \R+0.02) {$s_3$};
						
						% ==========================================================
						% retta r_V (verticale, a sinistra dell'arco) — campo u(y)
						% ==========================================================
						\draw[gray!45, thin] (\xV, -0.30) -- (\xV, \R+0.55)
						node[above, font=\scriptsize] {$r_V$};
						
						% proiezioni dei punti notevoli su r_V (tutti i centri in y=0)
						\draw[gray!40, dashed, thin] (-\R, 0) -- (\xV, 0);
						\draw[gray!40, dashed, thin] (0, \R)  -- (\xV, \R);
						
						\fill[gray!75] (\xV, 0) circle (1.3pt);
						\fill[gray!75] (\xV, \R) circle (1.3pt);
						
						% --- skyline u_1 di C_1 (verso sinistra: u<0) ---
						\draw[gray!75, thin] (\xV, 0) -- ({\xV-\ss/2}, \R);
						\foreach \k in {1,2,3,4}{
							\pgfmathsetmacro{\yy}{\k*\R/4}
							\pgfmathsetmacro{\uu}{(\ss/2)*\k/4}
							\draw[->, gray!80, thin] (\xV, \yy) -- ({\xV-\uu}, \yy);
						}
						\node[gray!85, font=\scriptsize, anchor=east]
						at ({\xV-\ss/2-0.02}, {\R/2+0.05}) {$u_1$};
						
						% --- skyline u_2 di C_2 (verso destra: u>0, dashed) ---
						\draw[black!55, thin, dashed] (\xV, 0) -- ({\xV+\ss/2}, \R);
						\foreach \k in {1,2,3,4}{
							\pgfmathsetmacro{\yy}{\k*\R/4}
							\pgfmathsetmacro{\uu}{(\ss/2)*\k/4}
							\draw[->, black!60, thin, dashed] (\xV, \yy) -- ({\xV+\uu}, \yy);
						}
						\node[black!70, font=\scriptsize, anchor=west]
						at ({\xV+\ss/2+0.02}, {\R/2+0.05}) {$u_2$};
						
						% --- etichette delle proiezioni dei centri/D su r_V ---
						\node[font=\scriptsize, anchor=north]
						at ({\xV-0.05}, 0) {$C_1\!\equiv\!C_2\!\equiv\!C_{12}$};
						%			\node[font=\scriptsize, anchor=east]
						%			at ({\xV-0.05}, \R) {$D$};
						
						% --- etichette dei valori al picco ---
						\node[font=\scriptsize, anchor=south]
						at ({\xV-\ss/4-0.02}, {\R-0.02}) {$\tfrac{s_3}{2}$};
						\node[font=\scriptsize, anchor=south]
						at ({\xV+\ss/4+0.02}, {\R-0.02}) {$\tfrac{s_3}{2}$};
						
						% ==========================================================
						% retta r_H (orizzontale, sotto) — campo v(x)
						% ==========================================================
						\draw[gray!45, thin] (-\R-0.40, \yH) -- (\R+0.40, \yH)
						node[right, font=\scriptsize] {$r_H$};
						
						\draw[gray!40, dashed, thin] (-\R, -0.45) -- (-\R, \yH);
						\draw[gray!40, dashed, thin] (0, -0.10)   -- (0, \yH);
						\draw[gray!40, dashed, thin] (\R, -0.45)  -- (\R, \yH);
						
						\fill[gray!75] (-\R, \yH) circle (1.3pt);
						\fill[gray!75] (0, \yH)   circle (1.3pt);
						\fill[gray!75] (\R, \yH)  circle (1.3pt);
						
						% --- skyline v_1 di C_1 (sopra r_H, sale da A verso D) ---
						\draw[gray!75, thin] (-\R, \yH) -- (0, {\yH+\ss/2});
						\foreach \k in {1,2,3,4}{
							\pgfmathsetmacro{\xx}{-\R + \k*\R/4}
							\pgfmathsetmacro{\vv}{(\ss/2)*\k/4}
							\draw[->, gray!80, thin] (\xx, \yH) -- (\xx, {\yH+\vv});
						}
						\node[gray!85, font=\scriptsize] at ({-\R/2-0.05}, {\yH-0.30}) {$v_1$};
						
						% --- skyline v_2 di C_2 (sopra r_H, scende da D verso B, dashed) ---
						\draw[black!55, thin] (0, {\yH+\ss/2}) -- (\R, \yH);
						\foreach \k in {0,1,2,3}{
							\pgfmathsetmacro{\xx}{\k*\R/4}
							\pgfmathsetmacro{\vv}{(\ss/2)*(1-\k/4)}
							\ifnum\k>0
							\draw[->, black!60, thin] (\xx, \yH) -- (\xx, {\yH+\vv});
							\fi
						}
						\node[black!70, font=\scriptsize] at ({\R/2+0.05}, {\yH-0.30}) {$v_2$};
						
						% --- etichetta picco su r_H ---
						\node[font=\scriptsize, anchor=south] at (0, {\yH+\ss/2+0.00}) {$\tfrac{s_3}{2}$};
						
						% --- etichette delle proiezioni dei centri su r_H ---
						\node[font=\scriptsize, anchor=north] at (-\R, {\yH-0.04}) {$C_1$};
						\node[font=\scriptsize, anchor=north] at (0, {\yH-0.04})   {$C_{12}$};
						\node[font=\scriptsize, anchor=north] at (\R, {\yH-0.04})  {$C_2$};
						
						
						% ==========================================================
						% Configurazione variata
						% ==========================================================
						
						% diagonali A-D_2' e B-D_1' (parallelogramma A D_1' D_2' B)
						\draw[gray!65, thin] (-\R, 0) -- ({\ss/2},  {\R+\ss/2});
						\draw[gray!65, thin] (\R, 0)  -- ({-\ss/2}, {\R+\ss/2});
						
						% raggi AD_1' e BD_2' (chorde dei due semiarchi deformati)
						\draw[gray!65, thin] (-\R, 0) -- ({-\ss/2}, {\R+\ss/2});
						\draw[gray!65, thin] (\R, 0)  -- ({\ss/2},  {\R+\ss/2});
						
						% marker D_1' e D_2'
						\fill[gray!65] ({-\ss/2}, {\R+\ss/2}) circle (1.5pt);
						\fill[gray!65] ({\ss/2},  {\R+\ss/2}) circle (1.5pt);
						
						\node[gray!65, font=\scriptsize, anchor=east]at ({-\ss/4}, {\R+\ss/2+0.2}) {$D_1'$};
						\node[gray!65, font=\scriptsize, anchor=west] at ({\ss/4},  {\R+\ss/2+0.2}) {$D_2'$};
						
						
						% semiarchi deformati (linearizzazione del campo di rotazione)
						\draw[gray!75, line width=0.9pt]
						plot[domain=90:180, samples=24, smooth, variable=\t]
						({\R*cos(\t) - (\ss/2)*sin(\t)},
						{\R*sin(\t) + (\ss/2)*(cos(\t)+1)});
						% C_2: rotazione oraria di ss/(2R) attorno a B=(R,0)
						\draw[gray!75, line width=0.9pt]
						plot[domain=0:90, samples=24, smooth, variable=\t]
						({\R*cos(\t) + (\ss/2)*sin(\t)},
						{\R*sin(\t) + (\ss/2)*(1-cos(\t))});
						
						% --- archi degli angoli di rotazione ---
						\pgfmathsetmacro{\dtheta}{atan(\ss/(2*\R))}
						\def\rB{0.7}
						
						
						\def\rB{0.90}                                  % era 0.7
						\pgfmathsetmacro{\dtheta}{atan(\ss/(2*\R))}
						
						% vartheta_1 in A: coda a 45°, punta a 45°+dtheta, antiorario
						\draw[->, gray!85, thin]
						({-\R + \rB*cos(45)}, {\rB*sin(45)})
						arc[start angle=45, end angle={45+\dtheta}, radius=\rB];
						\node[gray!85, font=\scriptsize, anchor=south west]	at ({-\R + \rB*cos(45)-0.04}, {\rB*sin(45)-0.30})
						{$\vartheta_1$};
						
						% vartheta_2 in B: coda a 135°, punta a 135°-dtheta, orario
						\draw[->, gray!85, thin]
						({\R + \rB*cos(135)}, {\rB*sin(135)})
						arc[start angle=135, end angle={135-\dtheta}, radius=\rB];
						\node[gray!85, font=\scriptsize, anchor=south east]	at ({\R + \rB*cos(135)+0.10}, {\rB*sin(135)-0.30})
						{$\vartheta_2$};
						
						% --- etichetta pannello ---
						\node[font=\scriptsize] at ({(\xV-\R)/2}, \yH-0.55) {(a)};
						
					\end{scope}
					
					% ============================================================
					% PANNELLO (b) — Cedimento s_4 (verticale in D)
					% ============================================================
					\begin{scope}[xshift=\dxp cm]
						
						% --- semiarchi ---
						\draw[line width=1.2pt] (-\R,0)
						arc[start angle=180, end angle=90, radius=\R];
						\draw[line width=1.2pt] (0,\R)
						arc[start angle=90, end angle=0, radius=\R];
						
						% --- cerniera al suolo in A ---
						\draw[thick] (-\R,0) -- (-\R-0.18,-0.30)
						-- (-\R+0.18,-0.30) -- cycle;
						\draw[thick] (-\R-0.28,-0.30) -- (-\R+0.28,-0.30);
						\foreach \x in {-0.20,-0.10,0,0.10,0.20}{
							\draw[thin] (-\R+\x,-0.30) -- (-\R+\x-0.09,-0.42);
						}
						% --- cerniera al suolo in B ---
						\draw[thick] (\R,0) -- (\R-0.18,-0.30)
						-- (\R+0.18,-0.30) -- cycle;
						\draw[thick] (\R-0.28,-0.30) -- (\R+0.28,-0.30);
						\foreach \x in {-0.20,-0.10,0,0.10,0.20}{
							\draw[thin] (\R+\x,-0.30) -- (\R+\x-0.09,-0.42);
						}
						
						% --- costruzione: retta AB estesa (luogo dei tre centri allineati) ---
						\draw[gray!55, dashed, thin] (-\R-0.30, 0) -- (\R+1.20, 0);
						% indicazione di C_{12} all'infinito orizzontale, lungo la retta AB
						\draw[->, gray!55, thin] (\R+0.40, 0) -- (\R+1.10, 0);
						\node[gray!75, font=\scriptsize, anchor=north]	at ({\R+1.10}, 0) {$C_{12}\to\infty$};
						
						% --- perni in A, B (esterni) e D (interno) ---
						\node[circle, fill=white, draw=black, inner sep=0pt,
						minimum size=4pt, line width=0.7pt] at (-\R,0) {};
						\node[circle, fill=white, draw=black, inner sep=0pt,
						minimum size=4pt, line width=0.7pt] at (\R,0) {};
						\node[circle, fill=white, draw=black, inner sep=0pt,
						minimum size=4pt, line width=0.7pt] at (0,\R) {};
						
						% --- etichette dei corpi ---
						\node[font=\scriptsize] at ({0.95*cos(110)-0.6},{0.95*sin(110)+0.50}) {$\mathcal{C}_1$};
						\node[font=\scriptsize] at ({0.95*cos(70)+1.2},{0.95*sin(70)})  {$\mathcal{C}_2$};
						
						% --- etichette punti notevoli ---
						\node[font=\scriptsize, anchor=east]  at ({-\R-0.10}, 0.18) {$A\!\equiv\!C_1$};
						\node[font=\scriptsize, anchor=west]  at ({\R+0.10}, 0.18)  {$B\!\equiv\!C_2$};
						\node[font=\scriptsize, anchor=south] at (-0.20, \R)     {$D$};
						
						% --- cedimento relativo s_4 in D: D_2 sale, D_1 scende ---
						\def\eps{0.08}
						\draw[->, thick] ({\eps}, \R+0.05)  -- ++(0, {\ss/2});
						\draw[->, thick] ({-\eps}, \R-0.05) -- ++(0, {-\ss/2});
						\node[font=\scriptsize, anchor=west] at ({\eps+0.10}, {\R+\ss/4}) {$s_4$};
						
						% ==========================================================
						% retta r_V (verticale, a sinistra dell'arco) — campo u(y)
						% ==========================================================
						\draw[gray!45, thin] (\xV, -0.30) -- (\xV, \R+0.55)
						node[above, font=\scriptsize] {$r_V$};
						
						\draw[gray!40, dashed, thin] (-\R, 0) -- (\xV, 0);
						\draw[gray!40, dashed, thin] (0, \R)  -- (\xV, \R);
						
						\fill[gray!75] (\xV, 0) circle (1.3pt);
						\fill[gray!75] (\xV, \R) circle (1.3pt);
						
						% skyline u_1 = u_2 (coincidenti)
						\draw[gray!75, thin] (\xV, 0) -- ({\xV+\ss/2}, \R);
						\foreach \k in {1,2,3,4}{
							\pgfmathsetmacro{\yy}{\k*\R/4}
							\pgfmathsetmacro{\uu}{(\ss/2)*\k/4}
							\draw[->, gray!80, thin] (\xV, \yy) -- ({\xV+\uu}, \yy);
						}
						\node[gray!85, font=\scriptsize, anchor=west]
						at ({\xV+\ss/3}, {\R/2+0.10}) {$u_1\!=\!u_2$};
						
						% --- etichette delle proiezioni dei centri/D su r_V ---
						\node[font=\scriptsize, anchor=north]
						at ({\xV-0.05}, 0) {$C_1\!\equiv\!C_2\!\equiv\!C_{12}$};
						\node[font=\scriptsize, anchor=east]
						at ({\xV-0.05}, \R) {$D$};
						
						% etichetta del valore al picco
						\node[font=\scriptsize, anchor=south]	at ({\xV+\ss/4+0.02}, {\R-0.02}) {$\tfrac{s_4}{2}$};
						
						% ==========================================================
						% retta r_H (orizzontale, sotto) — campo v(x)
						% ==========================================================
						\draw[gray!45, thin] (-\R-0.40, \yH) -- (\R+0.40, \yH)
						node[right, font=\scriptsize] {$r_H$};
						
						\draw[gray!40, dashed, thin] (-\R, -0.45) -- (-\R, \yH);
						\draw[gray!40, dashed, thin] (0, -0.10)   -- (0, \yH);
						\draw[gray!40, dashed, thin] (\R, -0.45)  -- (\R, \yH);
						
						\fill[gray!75] (-\R, \yH) circle (1.3pt);
						\fill[gray!75] (0, \yH)   circle (1.3pt);
						\fill[gray!75] (\R, \yH)  circle (1.3pt);
						
						% --- skyline v_1 di C_1 (sotto r_H, scende da A a D_1) ---
						\draw[gray!75, thin] (-\R, \yH) -- (0, {\yH-\ss/2});
						\foreach \k in {1,2,3,4}{
							\pgfmathsetmacro{\xx}{-\R + \k*\R/4}
							\pgfmathsetmacro{\vv}{(\ss/2)*\k/4}
							\draw[->, gray!80, thin] (\xx, \yH) -- (\xx, {\yH-\vv});
						}
						\node[gray!85, font=\scriptsize] at ({-\R/2-0.05}, {\yH-\ss/2-0.25}) {$v_1$};
						
						% --- skyline v_2 di C_2 (sopra r_H, scende da D_2 a B) ---
						\draw[gray!75, thin] (0, {\yH+\ss/2}) -- (\R, \yH);
						\foreach \k in {0,1,2,3}{
							\pgfmathsetmacro{\xx}{\k*\R/4}
							\pgfmathsetmacro{\vv}{(\ss/2)*(1-\k/4)}
							\ifnum\k>0
							\draw[->, gray!80, thin] (\xx, \yH) -- (\xx, {\yH+\vv});
							\fi
						}
						\node[gray!85, font=\scriptsize] at ({\R/2+0.05}, {\yH+\ss/2+0.20}) {$v_2$};
						\node[black!70, font=\scriptsize] at ({\R/2+0.05}, {\yH+\ss/2+0.20}) {$v_2$};
						
						% etichette del salto in D
						\node[font=\scriptsize, anchor=east] at (0.04, {\yH+\ss/4}) {$\tfrac{s_4}{2}$};
						\node[font=\scriptsize, anchor=west] at (0.02, {\yH-\ss/4-0.04}) {$\tfrac{s_4}{2}$};
						
						% --- etichette delle proiezioni dei centri su r_H ---
						\node[font=\scriptsize, anchor=north] at (-\R, {\yH-0.04}) {$C_1$};
						\node[font=\scriptsize, anchor=north] at (\R, {\yH-0.04})  {$C_2$};
						
						% ==========================================================
						% Configurazione variata
						% ==========================================================
						
						% diagonali AD_2' e BD_1' (parallelogramma A D_1' D_2' B)
						\draw[gray!65, thin] (-\R, 0) -- ({\ss/2}, {\R+\ss/2});
						\draw[gray!65, thin] (\R, 0)  -- ({\ss/2}, {\R-\ss/2});
						
						% corde AD_1' e BD_2'
						\draw[gray!65, thin] (-\R, 0) -- ({\ss/2}, {\R-\ss/2});
						\draw[gray!65, thin] (\R, 0)  -- ({\ss/2}, {\R+\ss/2});
						
						% marker D_1' e D_2'
						\fill[gray!65] ({\ss/2}, {\R-\ss/2}) circle (1.5pt);
						\fill[gray!65] ({\ss/2}, {\R+\ss/2}) circle (1.5pt);
						
						\node[gray!65, font=\scriptsize, anchor=west]  at ({0}, {\R-\ss/2-0.25}) {$D_1'$};
						\node[gray!65, font=\scriptsize, anchor=west]  at ({0}, {\R+\ss/2+0.20}) {$D_2'$};
						
						% semiarchi deformati (linearizzazione del campo di rotazione)
						% C_1: rotazione oraria di ss/(2R) attorno ad A=(-R,0)
						\draw[gray!75, line width=0.9pt]
						plot[domain=90:180, samples=24, smooth, variable=\t]
						({\R*cos(\t) + (\ss/2)*sin(\t)},
						{\R*sin(\t) - (\ss/2)*(cos(\t)+1)});
						% C_2: rotazione oraria di ss/(2R) attorno a B=(R,0)
						\draw[gray!75, line width=0.9pt]
						plot[domain=0:90, samples=24, smooth, variable=\t]
						({\R*cos(\t) + (\ss/2)*sin(\t)},
						{\R*sin(\t) + (\ss/2)*(1-cos(\t))});
						
						% --- archi degli angoli di rotazione (entrambi orari) ---
						\pgfmathsetmacro{\dtheta}{atan(\ss/(2*\R))}
						\def\rB{1.0}
						
						% vartheta_1 in A: coda a 45°, punta a 45°-dtheta, ORARIO
						\draw[->, gray!85, thin]
						({-\R + \rB*cos(45)}, {\rB*sin(45)})
						arc[start angle=45, end angle={45-\dtheta}, radius=\rB];
						\node[gray!85, font=\scriptsize, anchor=south west]	at ({-\R + \rB*cos(45)-0.02}, {\rB*sin(45)-0.15})
						{$\vartheta_1$};
						
						% vartheta_2 in B: coda a 135°, punta a 135°-dtheta, ORARIO
						\draw[->, gray!85, thin]
						({\R + \rB*cos(135)}, {\rB*sin(135)})
						arc[start angle=135, end angle={135-\dtheta}, radius=\rB];
						\node[gray!85, font=\scriptsize, anchor=south east]	at ({\R + \rB*cos(135)+0.60}, {\rB*sin(135)-0.40})
						{$\vartheta_2$};
						
						% --- etichetta pannello ---
						\node[font=\scriptsize] at ({(\xV-\R)/2}, \yH-0.55) {(b)};
						
					\end{scope}
					
				\end{tikzpicture}
			}
			\caption{Soluzione grafica per i cedimenti interni in $D$.
				\emph{(a)}~cedimento $s_3$ (orizzontale): tre centri
				distinti, $C_1\equiv A$, $C_2\equiv B$ fissati dai
				vincoli a cedimento nullo, $C_{12}$ punto medio di $AB$.
				\emph{(b)}~cedimento $s_4$ (verticale): caso degenere,
				$C_{12}$ punto improprio della direzione orizzontale,
				e i due semiarchi ruotano della stessa quantità attorno
				ai rispettivi centri $C_1\equiv A$ e $C_2\equiv B$.}
			\label{fig:arco_cedimenti_interni}
		\end{figure}
		
		\noindent\textit{Cedimento $s_2$ (verticale in $A$).}
		La cerniera in $B$ fissa il centro assoluto di
		$\mathcal{C}_2$: $C_2 \equiv B$. La cerniera interna in
		$D$ con cedimento nullo impone $C_{12} \equiv D$.
		Il vincolo orizzontale in $A$ è attivo, quindi $C_1$
		giace sull'orizzontale per $A$; per il teorema di
		allineamento $C_1$ giace inoltre sulla retta per
		$C_2 \equiv B$ e $C_{12} \equiv D$. L'intersezione
		della retta $BD$ con l'orizzontale per $A$ cade
		proprio in $B$: dunque $C_1 \equiv C_2 \equiv B$
		(si veda la \FIGref{fig:arco_cedimenti_esterni}\,(b))).
		
		I due centri assoluti coincidono: i due semiarchi
		ruotano dello stesso angolo. La rotazione si ricava
		dalla prestazione del vincolo verticale cedente in $A$,
		usando la retta $r_H$: il braccio efficace è la distanza
		orizzontale tra $C_1 \equiv B$ e $A$, pari a $2R$,
		quindi:
		\begin{equation*}
			\vartheta_1 = \vartheta_2 = \frac{s_2}{2R}
			\qquad \text{(antiorario, dalla figura)}.
		\end{equation*}
		
		
		\smallskip
		\noindent\textit{Cedimento relativo $s_3$
			(orizzontale in $D$).}
		La cerniera in $A$ con cedimento nullo fissa
		$C_1 \equiv A$; la cerniera in $B$ fissa $C_2 \equiv B$.
		Il \CRr\ $C_{12}$ giace sulla retta $AB$; la condizione
		di cedimento relativo verticale nullo in $D$ impone che
		$C_{12}$ coincida con la proiezione di $D$ sull'asse
		$AB$, cioè con il punto medio di $AB$
		(si veda la \FIGref{fig:arco_cedimenti_interni}\,(a)). 
		Si usa $r_V$: il braccio efficace è la distanza verticale
		tra $C_{12}$ e $D$, pari a $R$, quindi:
		\begin{equation*}
			\vartheta_{12} = -\frac{s_3}{R}, \qquad
			\vartheta_1 = \frac{s_3}{2R}, \qquad
			\vartheta_2 = -\frac{s_3}{2R},
		\end{equation*}
		
		\noindent\textit{Cedimento relativo $s_4$
			(verticale in $D$).}
		La cerniera in $A$ con cedimento nullo fissa
		$C_1 \equiv A$; la cerniera in $B$ fissa $C_2 \equiv B$.
		La condizione di cedimento relativo orizzontale nullo
		in $D$ impone che il moto relativo dei due semiarchi
		in $D$ sia puramente verticale: il \CRr\ $C_{12}$ è
		il punto improprio della direzione orizzontale, ossia
		i due semiarchi traslano relativamente in direzione
		verticale con la stessa rotazione
		$\vartheta_1 = \vartheta_2$ (si veda la \FIGref{fig:arco_cedimenti_interni}\,(b)). Si usa $r_H$: il braccio
		efficace è la distanza verticale tra $C_1 \equiv A$ e
		$D$, pari a $R$, quindi:
		\begin{equation*}
			\vartheta_{12} = 0, \qquad
			\vartheta_1 = \vartheta_2 = -\frac{s_4}{2R}.
		\end{equation*}
		La configurazione complessiva si ottiene sovrapponendo
		i quattro contributi: i risultati sono raccolti nella
		\TABref{tab:arco_cedimenti}, la cui riga
		``Totale'' coincide con la soluzione analitica	precedentemente ottenuta.
		
		
		
		\begin{table}[H]
			\small 
			\centering
			\renewcommand{\arraystretch}{1.8}
			\begin{tabular}{cccccc}
				\toprule
				Cedimento & $u(O_1)$ & $v(O_1)$ &
				$\vartheta_1$ & $\vartheta_2$ \\
				\midrule
				$s_1$ & $s_1$ & $0$ &
				$\dfrac{s_1}{2R}$ & $-\dfrac{s_1}{2R}$ \\[6pt]
				$s_2$ & $0$ & $-s_2$ &
				$\dfrac{s_2}{2R}$ & $\dfrac{s_2}{2R}$ \\[6pt]
				$s_3$ & $0$ & $0$ &
				$\dfrac{s_3}{2R}$ & $-\dfrac{s_3}{2R}$ \\[6pt]
				$s_4$ & $0$ & $0$ &
				$-\dfrac{s_4}{2R}$ & $-\dfrac{s_4}{2R}$ \\[6pt]
				\midrule
				Totale & $s_1$ & $-s_2$ &
				$\dfrac{s_1+s_2+s_3-s_4}{2R}$ &
				$\dfrac{-s_1+s_2-s_3-s_4}{2R}$ \\
				\bottomrule
			\end{tabular}
			\caption{Arco a tre cerniere: contributi dei singoli
				cedimenti. Le colonne $u(O_2)=0$ e $v(O_2)=0$
				sono omesse. La riga ``Totale'' coincide con la
				soluzione analitica.}
			\label{tab:arco_cedimenti}
		\end{table}
		
	\end{svolgimento}
\end{esempio}

\subsection{Problema statico: metodo dei corpi liberi}
\label{sec:corpi_liberi}



Il metodo\index{metodo!dei corpi liberi|textbf} sfrutta la possibilità di isolare ciascun corpo
dal sistema e scriverne le equazioni di equilibrio. La scelta
del polo di momento è lo strumento chiave per estrarre
un'incognita alla volta. L'equilibrio deve essere garantito
per ogni corpo singolo, per ogni sottoinsieme di corpi e per
l'intero sistema: queste equazioni non sono indipendenti, ma
scelte opportunamente consentono di disaccoppiare le incognite.

\paragraph{Procedura.}
\begin{enumerate}[label=\textbf{(\arabic*)}, leftmargin=*,
	itemsep=4pt]
	
	\item \textbf{Schema del corpo libero.}\index{schema del corpo libero}
	Si isola ciascun corpo, sostituendo ogni vincolo semplice
	esterno con la corrispondente reazione $R_k\,\ab_k$ e ogni
	vincolo interno con una coppia di reazioni uguali e
	opposte sui due corpi connessi.
	
	\item \textbf{Scelta del polo di momento.}
	Per determinare la reazione $R_k$:
	\begin{enumerate}[label=\textbf{(\alph*)},
		leftmargin=*, itemsep=2pt]
		\item si individua il sottosistema (corpo singolo
		o insieme di corpi) su cui $R_k$ è l'unica reazione
		incognita, o una delle poche;
		\item si scrive l'equazione di momento rispetto al
		polo $P$ che annulla il momento di tutte le altre
		reazioni incognite, ossia il punto per cui passano
		(o verso cui convergono) le loro linee d'azione;
		\item l'equazione risultante contiene la sola $R_k$
		e si risolve immediatamente.
	\end{enumerate}
	
	\item \textbf{Equilibrio globale.}
	Le equazioni di equilibrio dell'intero sistema si ottengono
	sommando quelle dei singoli corpi: le reazioni interne si
	cancellano a coppie e rimangono le sole reazioni esterne.
	Queste equazioni, combinate opportunamente con quelle dei
	singoli corpi, consentono di determinare tutte le incognite.
	
\end{enumerate}

\paragraph{Indipendenza delle equazioni di momento.}\index{equilibrio!equazioni di momento}
Poiché si dispone di tre equazioni di equilibrio per ogni
corpo rigido piano, si pone la questione se sia possibile
sostituire le equazioni di forza con ulteriori equazioni di
momento, e in quali condizioni le equazioni così ottenute
siano indipendenti.

\begin{proposizione}[Indipendenza delle equazioni di momento]
	\label{prop:indip_momenti}
	Le equazioni di momento scritte rispetto a tre poli $O$,
	$P$, $Q$ per un corpo rigido piano sono tre equazioni
	indipendenti --- equivalenti alle tre equazioni di
	equilibrio standard --- se e solo se i tre poli non sono
	collineari.
\end{proposizione}

\begin{proof}
	Il momento rispetto a un polo $P$ si lega a quello
	rispetto a un polo $O$ dalla formula di trasporto:
	\[
	\mu(P) = \mu(O) -
	\bigl(\overrightarrow{OP}\times \fb\bigr)\cdot\ab_z,
	\]
	dove $\fb$ è la risultante delle forze. Se i tre poli
	$O$, $P$, $Q$ sono collineari, esiste
	$\lambda\in\mathbb{R}$ tale che
	$\overrightarrow{OQ} = \lambda\,\overrightarrow{OP}$,
	e quindi:
	\[
	\mu(Q)
	= \mu(O) -
	\bigl(\overrightarrow{OQ}\times\fb\bigr)\cdot\ab_z
	= \mu(O) - \lambda\,
	\bigl(\overrightarrow{OP}\times\fb\bigr)\cdot\ab_z
	= (1-\lambda)\,\mu(O) + \lambda\,\mu(P).
	\]
	L'equazione $\mu(Q) = 0$ è dunque combinazione lineare
	di $\mu(O) = 0$ e $\mu(P) = 0$: le tre equazioni sono
	dipendenti.
	
	Viceversa, supponiamo $\mu(O) = \mu(P) = \mu(Q) = 0$
	con $O$, $P$, $Q$ non collineari. Dalle prime due
	condizioni:
	\[
	\mu(P) - \mu(O)
	= -\bigl(\overrightarrow{OP}\times\fb\bigr)
	\cdot\ab_z = 0,
	\]
	ossia $\bm{f}$ è parallela a $\overrightarrow{OP}$.
	Analogamente, da $\mu(Q) = \mu(O) = 0$, si ottiene
	che $\fb$ è parallela anche a $\overrightarrow{OQ}$.
	Poiché $O$, $P$, $Q$ non sono collineari,
	$\overrightarrow{OP}$ e $\overrightarrow{OQ}$ non sono
	paralleli: l'unico vettore piano parallelo a entrambi è
	$\fb = \zerovb$. Sostituendo in $\mu(O) = 0$ si
	ottiene infine $\mu(O) = 0$: le tre equazioni di momento
	implicano dunque $\fb = \zerovb$ e $\mu(O) = 0$,
	ovvero le tre equazioni di equilibrio standard.
\end{proof}

\begin{oss}[Conseguenza pratica]
	Se si vogliono sostituire le equazioni di forza con
	equazioni di momento, i poli scelti devono formare un
	triangolo non degenere. Il caso limite di tre poli
	allineati produce tre equazioni dipendenti: le due
	equazioni di forza nella direzione della retta di
	allineamento non vengono implicate e l'equilibrio non
	è completamente determinato.
\end{oss}


\begin{esempio}[Arco a tre cerniere --- problema statico]
	\label{ex:arco_statico}
	Si consideri lo stesso arco a tre cerniere
	dell'\ESEref{ex:arco_cedimenti}. Sono applicati
	una forza verticale $F$ (verso il basso) nel punto $P_2$,
	punto medio dell'asse di $\mathcal{C}_2$, e una coppia
	$\mu$ (antioraria) nel punto $P_1$, punto medio dell'asse
	di $\mathcal{C}_1$ (\FIGref{fig:arco_stat_dati}). Determinare le reazioni vincolari.
	
	% ============================================================
	% FIGURA: Arco a tre cerniere — problema statico
	% ============================================================
	\begin{figure}[ht]
		\centering
		\begin{tikzpicture}[>=Stealth, line width=0.8pt]
			
			\def\R{2.0}
			\def\arr{0.55}
			\def\dx{1.4}    % separazione tra C_1 e C_2 nei corpi liberi
			
			% ============================================================
			% PANNELLO (a) — dati del problema
			% ============================================================
			\begin{scope}[xshift=0cm]
				
				% semiarchi
				\draw[line width=1.5pt] (-\R,0)
				arc[start angle=180, end angle=90, radius=\R];
				\draw[line width=1.5pt] (0,\R)
				arc[start angle=90, end angle=0, radius=\R];
				
				% cerniera in A
				\draw[thick] (-\R,0) -- (-\R-0.22,-0.38)
				-- (-\R+0.22,-0.38) -- cycle;
				\draw[thick] (-\R-0.32,-0.38) -- (-\R+0.32,-0.38);
				\foreach \x in {-0.24,-0.12,0,0.12,0.24}{
					\draw[thin] (-\R+\x,-0.38) -- (-\R+\x-0.10,-0.52);
				}
				% cerniera in B
				\draw[thick] (\R,0) -- (\R-0.22,-0.38)
				-- (\R+0.22,-0.38) -- cycle;
				\draw[thick] (\R-0.32,-0.38) -- (\R+0.32,-0.38);
				\foreach \x in {-0.24,-0.12,0,0.12,0.24}{
					\draw[thin] (\R+\x,-0.38) -- (\R+\x-0.10,-0.52);
				}
				
				% perni A, B, D
				\node[circle, fill=white, draw=black, inner sep=0pt,
				minimum size=5pt, line width=0.8pt] at (-\R,0) {};
				\node[circle, fill=white, draw=black, inner sep=0pt,
				minimum size=5pt, line width=0.8pt] at (\R,0) {};
				\node[circle, fill=white, draw=black, inner sep=0pt,
				minimum size=5pt, line width=0.8pt] at (0,\R) {};
				
				% etichette punti
				\node[left, font=\scriptsize] at (-\R,+0.2) {$A$};
				\node[right, font=\scriptsize] at (\R,+0.2) {$B$};
				\node[font=\scriptsize] at (-0.,\R+0.30) {$D$};
				
				% etichette corpi (esterne)
				\node[font=\small] at ({(\R+0.35)*cos(150)},{(\R+0.35)*sin(150)})	{$\mathcal{C}_1$};
				\node[font=\small] at ({(\R+0.35)*cos(35)},{(\R+0.35)*sin(35)})		{$\mathcal{C}_2$};
				
				% punto P_1 (centro asse C_1, a 135°)
				\node[circle, fill=black, inner sep=0pt, minimum size=2pt]
				at ({\R*cos(135)},{\R*sin(135)}) {};
				\node[font=\scriptsize] at ({(\R-0.30)*cos(135)},{(\R-0.30)*sin(135)+0.07}) {$P_1$};
				% coppia mu antioraria in P_1
				\draw[->, thick] ({\R*cos(135)+0.30},{\R*sin(135)+0.05}) arc[start angle=10, end angle=310, radius=0.30];
				\node[font=\scriptsize]	at ({\R*cos(135)+0.55},{\R*sin(135)+0.10}) {$\mu$};
				
				% punto P_2 (centro asse C_2, a 45°)
				\node[circle, fill=black, inner sep=0pt, minimum size=2pt]
				at ({\R*cos(45)},{\R*sin(45)}) {};
				\node[font=\scriptsize]	at ({(\R-0.30)*cos(45)},{(\R-0.30)*sin(45)}) {$P_2$};
				% forza F verticale verso il basso in P_2
				\draw[->, thick] ({\R*cos(45)},{\R*sin(45)+\arr}) -- ({\R*cos(45)},{\R*sin(45)});
				\node[font=\scriptsize]	at ({\R*cos(45)+0.25},{\R*sin(45)+\arr-0.10}) {$F$};
				
				
				
				% etichetta pannello
				\node[font=\scriptsize] at (0,-0.50) {(a)};
				
			\end{scope}
			
			% ============================================================
			% PANNELLO (b) — schemi dei corpi liberi
			% ============================================================
			\begin{scope}[xshift=7.0cm]
				
				% ----- C_1 libero (decentrato a sinistra) -----
				\begin{scope}[xshift=-\dx cm]
					
					% semiarco C_1
					\draw[line width=1.5pt] (-\R,0)
					arc[start angle=180, end angle=90, radius=\R];
					
					% cerniera in A (grigia, vincolo rotto)
					\draw[thick, gray!50] (-\R,0) -- (-\R-0.22,-0.38)
					-- (-\R+0.22,-0.38) -- cycle;
					\draw[thick, gray!50] (-\R-0.32,-0.38) -- (-\R+0.32,-0.38);
					\foreach \x in {-0.24,-0.12,0,0.12,0.24}{
						\draw[thin, gray!50] (-\R+\x,-0.38) -- (-\R+\x-0.10,-0.52);
					}
					
					% perni A, D
					\node[circle, fill=white, draw=black, inner sep=0pt,
					minimum size=5pt, line width=0.8pt] at (-\R,0) {};
					\node[circle, fill=white, draw=black, inner sep=0pt,
					minimum size=5pt, line width=0.8pt] at (0,\R) {};
					
					% etichette
					\node[left, font=\scriptsize] at (-\R,-0.) {$A$};
					\node[below=2pt, font=\scriptsize] at (0,\R+0.05) {$D$};
					\node[font=\small] at ({(\R+0.35)*cos(150)},{(\R+0.35)*sin(150)})	{$\mathcal{C}_1$};
					
					
					
					% reazioni in A
					\draw[->, thick] (-\R,0.0) -- (-\R+\arr,0.0)
					node[above, font=\scriptsize] {$X_A$};
					\draw[->, thick] (-\R-0.05,0) -- (-\R-0.05,\arr)
					node[left, yshift=-3pt, font=\scriptsize] {$Y_A$};
					
					% reazioni interne in D (su C_1): X_D verso destra, Y_D verso l'alto
					\draw[->, thick] (0,\R) -- (\arr,\R)
					node[below, font=\scriptsize] {$X_D$};
					\draw[->, thick] (-0.05,\R) -- (-0.05,\R+\arr)
					node[left, yshift=-3pt, font=\scriptsize] {$Y_D$};
					
					% coppia mu in P_1
					\node[circle, fill=black, inner sep=0pt, minimum size=2pt]
					at ({\R*cos(135)},{\R*sin(135)}) {};
					\draw[->, thick]
					({\R*cos(135)+0.25},{\R*sin(135)+0.05})
					arc[start angle=10, end angle=310, radius=0.25];
					\node[font=\scriptsize]	at ({\R*cos(135)+0.50},{\R*sin(135)+0.05}) {$\mu$};
					\node[font=\scriptsize] at ({(\R-0.30)*cos(135)},{(\R-0.30)*sin(135)+0.10}) {$P_1$};
					
				\end{scope}
				
				% ----- C_2 libero (decentrato a destra) -----
				\begin{scope}[xshift=\dx-0.10 cm]
					
					% semiarco C_2
					\draw[line width=1.5pt] (0,\R)
					arc[start angle=90, end angle=0, radius=\R];
					
					% cerniera in B (grigia)
					\draw[thick, gray!50] (\R,0) -- (\R-0.22,-0.38)
					-- (\R+0.22,-0.38) -- cycle;
					\draw[thick, gray!50] (\R-0.32,-0.38) -- (\R+0.32,-0.38);
					\foreach \x in {-0.24,-0.12,0,0.12,0.24}{
						\draw[thin, gray!50] (\R+\x,-0.38) -- (\R+\x-0.10,-0.52);
					}
					
					% perni B, D
					\node[circle, fill=white, draw=black, inner sep=0pt,
					minimum size=5pt, line width=0.8pt] at (\R,0) {};
					\node[circle, fill=white, draw=black, inner sep=0pt,
					minimum size=5pt, line width=0.8pt] at (0,\R) {};
					
					% etichette
					\node[left, font=\scriptsize] at (\R,+0.) {$B$};
					\node[above, font=\scriptsize] at (0,\R) {$D$};
					\node[font=\small] at ({(\R+0.35)*cos(35)},{(\R+0.35)*sin(35)})	{$\mathcal{C}_2$};
					
					
					% reazioni in B
					\draw[->, thick] (\R,0.0) -- (\R+\arr,0.0)
					node[above, font=\scriptsize] {$X_B$};
					\draw[->, thick] (\R+0.05,0) -- (\R+0.05,\arr)
					node[left, yshift=-3pt, font=\scriptsize] {$Y_B$};
					
					% reazioni interne in D (su C_2): opposte (verso sinistra e basso)
					\draw[->, thick] (0,\R) -- (-\arr,\R)
					node[above, font=\scriptsize] {$X_D$};
					\draw[->, thick] (0.05,\R) -- (0.05,\R-\arr)
					node[right, yshift=3pt, font=\scriptsize] {$Y_D$};
					
					% forza F in P_2
					\node[circle, fill=black, inner sep=0pt, minimum size=2pt]
					at ({\R*cos(45)},{\R*sin(45)}) {};
					\draw[->, thick] ({\R*cos(45)},{\R*sin(45)+\arr}) -- ({\R*cos(45)},{\R*sin(45)});
					\node[font=\scriptsize]	at ({\R*cos(45)+0.25},{\R*sin(45)+\arr-0.10}) {$F$};
					\node[font=\scriptsize]	at ({(\R-0.30)*cos(45)},{(\R-0.30)*sin(45)}) {$P_2$};
					
				\end{scope}
				
				% etichetta pannello
				\node[font=\scriptsize] at (-0.60,-0.50) {(b)};
				
			\end{scope}
			
		\end{tikzpicture}
		\caption{Arco a tre cerniere
			(\ESEref{ex:arco_statico}):
			\emph{(a)}~dati del problema statico ---
			carico verticale $F$ in $P_2$ (punto medio
			dell'asse di $\mathcal{C}_2$) e coppia $\mu$
			antioraria in $P_1$ (punto medio dell'asse
			di $\mathcal{C}_1$); le frecce indicano il
			verso positivo dei carichi assegnati.
			\emph{(b)}~schemi dei corpi liberi di
			$\mathcal{C}_1$ (sinistra) e $\mathcal{C}_2$
			(destra): reazioni esterne $X_A, Y_A$ in $A$
			e $X_B, Y_B$ in $B$; reazioni interne $X_D, Y_D$
			in $D$, uguali e opposte sui due corpi.}
		\label{fig:arco_statico}
	\end{figure}
	
	\begin{svolgimento}
		A differenza del problema cinematico, per il problema
		statico non si costruisce esplicitamente il sistema
		$\bm{A}^{\top}\bm{r} = -\bm{f}$: la strategia dei
		corpi liberi con scelta opportuna del polo di momento
		consente di estrarre le reazioni una alla volta,
		rendendo superflua la formulazione matriciale.
		
		Le reazioni incognite sono: $X_A$, $Y_A$ in $A$;
		$X_B$, $Y_B$ in $B$; $X_D$, $Y_D$ in $D$ (uguali e
		opposte sui due corpi). I versi positivi assunti per
		ciascuna reazione sono indicati nella
		\FIGref{fig:arco_stat_cl}.
		
		\smallskip
		\noindent\textit{Equilibrio globale $(\mathcal{C}_1 +
			\mathcal{C}_2)$ ai momenti rispetto ad $A$.}
		Le reazioni interne in $D$ si cancellano:
		\begin{equation*}
			\sum_{(1+2)}\mu(A)= 0:\quad
			Y_B \cdot 2R
			- \frac{F}{2} \cdot \frac{3}{2}R
			+ \mu = 0
			\;\Longrightarrow\;
			Y_B = \frac{\frac{3}{4}FR - FR}{2R}
			= -\frac{F}{8}.
		\end{equation*}
		La reazione $Y_B$ è diretta verso il basso.
		
		\smallskip
		\noindent\textit{Equilibrio di $\mathcal{C}_2$ ai momenti
			rispetto a $D$.}
		Le reazioni interne in $D$ si eliminano; $P_2$ dista
		$R/2$ orizzontalmente da $D$:
		\begin{equation*}
			\sum_{(2)}\mu(D) = 0:\quad
			X_B \cdot R
			- \frac{F}{2} \cdot \frac{R}{2}
			- \frac{F}{8} \cdot R = 0
			\;\Longrightarrow\;
			X_B = \frac{F}{4} + \frac{F}{8} = \frac{3F}{8}.
		\end{equation*}
		
		\smallskip
		\noindent\textit{Equilibrio globale alla traslazione.}
		\begin{equation*}
			\sum_{(1+2)} X = 0:\quad
			X_A + X_B = 0
			\;\Longrightarrow\;
			X_A = -\frac{3F}{8}
			\quad\text{(verso sinistra)},
		\end{equation*}
		\begin{equation*}
			\sum_{(1+2)} Y = 0:\quad
			Y_A + Y_B - \frac{F}{2} = 0
			\;\Longrightarrow\;
			Y_A = \frac{F}{2} + \frac{F}{8} = \frac{5F}{8}
			\quad\text{(verso l'alto)}.
		\end{equation*}
		
		\smallskip
		\noindent\textit{Reazioni interne in $D$.}
		Dall'equilibrio alla traslazione di $\mathcal{C}_2$:
		\begin{equation*}
			\sum_{(2)} X = 0:\quad
			X_D = X_B = \frac{3F}{8}
			\quad\text{(verso sinistra su $\mathcal{C}_2$)},
		\end{equation*}
		\begin{equation*}
			\sum_{(2)} Y = 0:\quad
			Y_D = \frac{F}{2} - Y_B =
			\frac{F}{2} + \frac{F}{8} = \frac{5F}{8}
			\quad\text{(verso l'alto su $\mathcal{C}_2$)}.
		\end{equation*}
		
		\noindent\textit{Verifica: equilibrio di $\mathcal{C}_1$.}
		\begin{equation*}
			\sum_{(1)} X = 0:\quad
			X_A + (-X_D) = -\frac{3F}{8} + \left(-\frac{3F}{8}\right)
			= 0, \checkmark
		\end{equation*}
		\begin{equation*}
			\sum_{(1)} Y = 0:\quad
			Y_A + (-Y_D) = \frac{5F}{8} - \frac{5F}{8} = 0,
			\checkmark
		\end{equation*}
		\begin{equation*}
			\sum_{(1)}\mu(A) = 0:\quad
			\mu + (-Y_D)\cdot R - (-X_D)\cdot R =
			FR - \frac{5F}{8}R - \frac{3F}{8}R =
			FR - FR = 0. \checkmark
		\end{equation*}
	\end{svolgimento}
\end{esempio}



%---------------------------------------------------------------------------------------------------------------------------- 
\section{Sistemi labili}\label{sec:labile}
%---------------------------------------------------------------------------------------------------------------------------- 
Un sistema è \emph{labile}\index{labilita@labilità} di grado $g_\ell$ quando
$\dim\Ker(\Abb) = g_\ell > 0$: esistono $g_\ell$ configurazioni
linearmente indipendenti compatibili con tutti i vincoli perfetti,
dette \emph{meccanismi}. La condizione necessaria è $\nu = n - m > 0$
(vincoli in numero insufficiente); tuttavia la labilità può presentarsi
anche con $\nu = 0$, in concorrenza con un'iperstaticità di pari grado
per effetto della degenerazione geometrica: è il caso labile-iperstatico,
le cui condizioni operative sono trattate in
\PARref{sec:labile_iperstatico}.

I due problemi si comportano in modo duale: il problema cinematico è
\emph{sottodeterminato} (infinite soluzioni), mentre quello statico è
\emph{sovradeterminato} (soluzione possibile solo per certi carichi).


\subsection{Problema cinematico: vincoli fittizi e	coordinate lagrangiane}
\label{sec:labile_cin}

Poiché il sistema ha $g_\ell$ gradi di labilità, la
soluzione generale del problema cinematico
\eqref{eq:due_problemi}, $\bm{A}\bm{u} = \bm{s}$, è:
\begin{equation}\label{eq:sol_labile_cin}
	\bm{u} = \bm{u}_0 + \bm{U}_q\,\bm{q},
	\qquad \bm{q} \in \mathbb{R}^{g_\ell},
\end{equation}
dove $\bm{u}_0$ è una soluzione particolare e le
colonne di $\bm{U}_q$ formano una base di
$\Ker(\bm{A})$. Le componenti di $\bm{q}$ sono le
\emph{coordinate lagrangiane} del meccanismo,
parametri liberi che descrivono la parte di
configurazione non determinata dai vincoli reali.

Il metodo operativo per costruire questa soluzione
procede aggiungendo al sistema $g_\ell$
\emph{vincoli fittizi}\index{vincolo!fittizio|textbf}, scelti in modo da rendere
il sistema isostatico, e poi sfruttando la linearità
per separare i contributi dei cedimenti reali da
quelli fittizi.

\paragraph{Procedura.}
\begin{enumerate}[label=\textbf{(\arabic*)},
	leftmargin=*, itemsep=4pt]
	
	\item \textbf{Aggiunta di vincoli fittizi.}
	Si aggiungono $g_\ell$ vincoli fittizi, scelti
	in posizioni e direzioni tali che il sistema
	risultante, anche detto \textit{sistema principale}\index{sistema principale},  sia isostatico ($\nu = 0$,
	$\rg(\bm{A}_{\rm est}) = n$). La scelta è
	arbitraria purché i vincoli aggiunti siano
	geometricamente non degeneri rispetto a quelli
	reali.
	
	\item \textbf{Soluzione particolare: cedimenti
		reali, cedimenti fittizi nulli.}
	Si assegnano ai vincoli fittizi cedimento zero
	e si risolve il problema cinematico del sistema
	isostatico esteso con i soli cedimenti reali
	$\bm{s}$, applicando il metodo
	grafo-analitico (\PARref{sec:grafo_analitico}).
	La soluzione $\bm{u}_0$ così ottenuta
	è una soluzione particolare di
	$\bm{A}_{\rm est}\bm{u} = \bm{s}$: soddisfa tutti i
	vincoli reali e impone che i vincoli fittizi
	non cedano.
	
	\item \textbf{Meccanismi: cedimenti reali nulli, cedimenti fittizi unitari.}
	Si pongono tutti i cedimenti reali a zero e si
	assegna valore unitario a ciascun cedimento
	fittizio alla volta: $q_k = 1$, $q_j = 0$ per
	$j \neq k$. Ogni problema così ottenuto è un
	problema cinematico del sistema isostatico esteso,
	risolto con il metodo grafo-analitico
	(\PARref{sec:grafo_analitico}); la soluzione
	$\bm{u}_k$ è un meccanismo del sistema originale
	(soddisfa $\bm{A}\bm{u}_k = \bm{0}$) e
	costituisce la $k$-esima colonna di $\bm{U}_q$.
	
	\item \textbf{Soluzione generale.}
	Per linearità, la soluzione generale si ottiene
	sovrapponendo i contributi dei passi 2 e 3:
	\begin{equation}
		\bm{u} = \bm{u}_0 +
		\sum_{k=1}^{g_\ell} q_k\,\bm{u}_k,
	\end{equation}
	con i parametri $q_k$ liberi: ciascuno
	rappresenta l'ampiezza con cui il $k$-esimo
	meccanismo si sovrappone alla soluzione
	particolare.
	
\end{enumerate}

\begin{oss}[Arbitrarietà della scelta dei vincoli fittizi]
	La scelta dei vincoli fittizi non è unica: qualsiasi
	insieme di $g_\ell$ vincoli geometricamente non
	degeneri rispetto a quelli reali rende il sistema
	isostatico. Nella formulazione algebrica, i vincoli
	fittizi corrispondono alla scelta delle variabili
	fuori base del sistema $\bm{A}\bm{u} = \bm{s}$:
	scelte diverse corrispondono a basi diverse per
	$\Ker(\bm{A})$, cioè a coordinate lagrangiane che
	sono combinazioni lineari l'una dell'altra.
	L'insieme di tutte le configurazioni ammissibili
	$\bm{u}_0 + \Ker(\bm{A})$ è però invariante:
	soluzioni particolari e meccanismi cambiano forma
	ma descrivono sempre lo stesso insieme di
	configurazioni. In pratica: si scelgano i vincoli
	fittizi in modo da rendere i calcoli il più
	semplice possibile.
\end{oss}


\begin{esempio}[Sistema labile di grado 1 --- problema cinematico]
	\label{ex:labile_L}
	Si consideri un corpo rigido piano a forma di L,
	composto da un tratto orizzontale di lunghezza $\ell$
	che va da $A$ a $D$ e da un tratto verticale di
	lunghezza $\ell$ che va da $D$ a $B$. Il corpo è
	vincolato da un carrello verticale in $A$ (asse $\ab_y$)
	e da un carrello orizzontale in $B$ (asse $\ab_x$).
	Sono assegnati i cedimenti $s_A$ (verticale in $A$) e
	$s_B$ (orizzontale in $B$)
	(\FIGref{fig:labile_L}\,(a)). Determinare la
	configurazione del sistema.
	
	\begin{figure}[ht]
		\centering
		\begin{tikzpicture}[>=Stealth, line width=0.8pt]
			
			\def\L{2.0}
			\def\sA{0.3}
			\def\sB{0.3}
			\def\q{0.3}
			
			% ============================================================
			% PANNELLO SINISTRO — dati del problema
			% ============================================================
			\begin{scope}[xshift=0cm]
				
				% corpo a L
				\draw[line width=1.5pt] (0,0) -- (\L,0) -- (\L,\L);
				
				% rette tratteggiate per C
				\draw[gray!50, dashed, thin]
				(0,-0.3) -- (0,\L+0.2);
				\draw[gray!50, dashed, thin]
				(-0.3,\L) -- (\L+0.2,\L);
				
				% carrello verticale in A
				\draw[thick] (0,0) -- (-0.28,-0.48)
				-- (0.28,-0.48) -- cycle;
				\draw[thick] (-0.18,-0.55) circle (0.08);
				\draw[thick] ( 0.18,-0.55) circle (0.08);
				\draw[thick] (-0.38,-0.65) -- (0.38,-0.65);
				\foreach \x in {-0.30,-0.15,0,0.15,0.30}{
					\draw[thin] (\x,-0.65) -- (\x-0.12,-0.83);
				}
				
				% carrello orizzontale in B
				\draw[thick] (\L,\L) -- (\L+0.48,\L+0.28)
				-- (\L+0.48,\L-0.28) -- cycle;
				\draw[thick] (\L+0.55,\L+0.18) circle (0.08);
				\draw[thick] (\L+0.55,\L-0.18) circle (0.08);
				\draw[thick] (\L+0.65,\L+0.38) -- (\L+0.65,\L-0.38);
				\foreach \y in {-0.30,-0.15,0,0.15,0.30}{
					\draw[thin] (\L+0.65,\L+\y)
					-- (\L+0.83,\L+\y-0.12);
				}
				
				
				% punti
				\draw[fill=white, thick] (0,0) circle (3pt)
				node[above right, font=\scriptsize] {$A$};
				\draw[fill=white, thick] (\L,\L) circle (3pt)
				node[above left, font=\scriptsize] {$B$};
				\node[left, font=\scriptsize] at (0,\L) {$C$};
				
				% cedimento s_A (verticale, a sinistra di A)
				\draw[thick] (-0.35,0) -- (-0.15,0);
				\draw[->, thick] (-0.25,0) -- (-0.25,\sA)
				node[left, font=\scriptsize] {$s_A$};
				
				% cedimento s_B (orizzontale, sopra B)
				\draw[thick] (\L,\L+0.15) -- (\L,\L+0.35);
				\draw[->, thick] (\L,\L+0.25) -- ({\L+\sB},\L+0.25)
				node[above, font=\scriptsize] {$s_B$};
				
				% quota tratto orizzontale (sotto)
				\draw[<->, thin] (0,-1.0) -- (\L,-1.0)
				node[midway, above=2pt, font=\scriptsize] {$\ell$};
				
				% quota tratto verticale (a destra)
				\draw[<->, thin] (\L+1.2,0) -- (\L+1.2,\L)
				node[midway, right=2pt, font=\scriptsize] {$\ell$};
				
				% etichetta pannello
				\node[font=\scriptsize] at (\L/2,-1.4) {(a)};
				
			\end{scope}
			
			% ============================================================
			% PANNELLO DESTRO — sistema principale
			% ============================================================
			\begin{scope}[xshift=5.5cm]
				
				
				% carrello verticale fittizio in D (grigio)
				\draw[gray!50, thick] (\L,0) -- (\L-0.28,-0.48)
				-- (\L+0.28,-0.48) -- cycle;
				\draw[gray!50, thick] (\L-0.18,-0.55) circle (0.08);
				\draw[gray!50, thick] (\L+0.18,-0.55) circle (0.08);
				\draw[gray!50, thick] (\L-0.38,-0.65)
				-- (\L+0.38,-0.65);
				\foreach \x in {-0.30,-0.15,0,0.15,0.30}{
					\draw[gray!50, thin] (\L+\x,-0.65)
					-- (\L+\x-0.12,-0.83);
				}
				
				\fill[white]            (\L,0.10) arc(90:-180:0.10) -- cycle;
				\draw[gray!50, thick]   (\L,0.10) arc(90:-180:0.10);
				
				% corpo a L
				\draw[line width=1.5pt] (0,0) -- (\L,0) -- (\L,\L);
				
				
				
				% carrello verticale in A
				\draw[thick] (0,0) -- (-0.28,-0.48)
				-- (0.28,-0.48) -- cycle;
				\draw[thick] (-0.18,-0.55) circle (0.08);
				\draw[thick] ( 0.18,-0.55) circle (0.08);
				\draw[thick] (-0.38,-0.65) -- (0.38,-0.65);
				\foreach \x in {-0.30,-0.15,0,0.15,0.30}{
					\draw[thin] (\x,-0.65) -- (\x-0.12,-0.83);
				}
				
				% carrello orizzontale in B
				\draw[thick] (\L,\L) -- (\L+0.48,\L+0.28)
				-- (\L+0.48,\L-0.28) -- cycle;
				\draw[thick] (\L+0.55,\L+0.18) circle (0.08);
				\draw[thick] (\L+0.55,\L-0.18) circle (0.08);
				\draw[thick] (\L+0.65,\L+0.38) -- (\L+0.65,\L-0.38);
				\foreach \y in {-0.30,-0.15,0,0.15,0.30}{
					\draw[thin] (\L+0.65,\L+\y)
					-- (\L+0.83,\L+\y-0.12);
				}
				
				% cedimento q al carrello fittizio in D
				\draw[gray!60, thick] (\L+0.35,0) -- (\L+0.15,0);
				\draw[->, gray!60, thick] (\L+0.25,0) -- (\L+0.25,\q)
				node[right, font=\scriptsize] {$q$};
				
				% punti
				\draw[fill=white, thick] (0,0) circle (3pt)
				node[above right, font=\scriptsize] {$A$};
				\draw[fill=white, thick] (\L,\L) circle (3pt)
				node[above left, font=\scriptsize] {$B$};
				\node[above left, font=\scriptsize] at (\L,0) {$D$};
				
				
				
				% etichetta pannello
				\node[font=\scriptsize] at (\L/2,-1.4) {(b)};
				
			\end{scope}
			
		\end{tikzpicture}
		\caption{Corpo a L
			(\ESEref{ex:labile_L}): dati del problema
			cinematico \emph{(a)} e sistema principale con
			carrello fittizio in $D$ (grigio) e cedimento
			$q$ \emph{(b)}.}
		\label{fig:labile_L}
	\end{figure}
	
	\begin{svolgimento}
		
		\noindent\textbf{Classificazione.}
		Il corpo ha $n = 3$ gradi di libertà e $m = 2$ vincoli
		scalari (un carrello verticale in $A$ e un carrello
		orizzontale in $B$): $\nu = 3 - 2 = 1$. Il sistema è
		labile di grado $g_\ell = 1$. 		
		I due carrelli hanno assi ortogonali tra loro: la loro
		equivalenza cinematica è quella di una cerniera nel
		punto $C$, intersezione della verticale per $A$ e
		dell'orizzontale per $B$. Il meccanismo è dunque una
		rotazione attorno a $C$.
		
		\noindent\textbf{Sistema principale.}
		Si aggiunge un carrello verticale in $D$ con cedimento
		$q$ arbitrario (coordinata lagrangiana). Il sistema
		esteso ha $m = 3$ vincoli scalari ed è isostatico.
		
		\noindent\textbf{Formulazione analitica.}
		Le incognite sono $\bm{u} = \{u_O,\, v_O,\,
		\vartheta\}^\top$ con polo $O \equiv A$. Le equazioni
		di vincolo del sistema principale sono:
		\begin{equation*}
			\begin{bmatrix}
				0 & 1 &  0 \\
				1 & 0 & -\ell \\
				0 & 1 &  \ell
			\end{bmatrix}
			\begin{Bmatrix}
				u_O \\ v_O \\ \vartheta
			\end{Bmatrix}
			=
			\begin{Bmatrix}
				s_A \\ s_B \\ q
			\end{Bmatrix},
		\end{equation*}
		dove la prima riga è il carrello verticale in $A$,
		la seconda il carrello orizzontale in $B$, la terza
		il carrello verticale in $D$.
		
		Le prime due equazioni danno:
		\begin{equation*}
			v_O = s_A, \qquad u_O = s_B + \ell\,\vartheta.
		\end{equation*}
		La terza equazione dà:
		\begin{equation*}
			v_O + \ell\,\vartheta = q
			\;\Longrightarrow\;
			\vartheta = \frac{q - s_A}{\ell}.
		\end{equation*}
		Sostituendo nella seconda:
		\begin{equation*}
			u_O = s_B + q - s_A.
		\end{equation*}
		La soluzione generale è:
		\begin{equation*}\label{eq:labile_L_sol}
			\begin{Bmatrix}
				u_O \\ v_O \\ \vartheta
			\end{Bmatrix}
			=
			\underbrace{
				\begin{Bmatrix}
					s_B - s_A \\ s_A \\ -s_A/\ell
				\end{Bmatrix}
			}_{\bm{u}_0}
			+\, q\,
			\underbrace{
				\begin{Bmatrix}
					1 \\ 0 \\ 1/\ell
				\end{Bmatrix}
			}_{\bm{u}_1}.
		\end{equation*}
		
		\noindent\textbf{Interpretazione grafica per sovrapposizione.}
		Per il principio di sovrapposizione si considerano i
		tre cedimenti separatamente.
		
		\smallskip
		\noindent\textit{Cedimento $s_A$ (verticale in $A$).}
		I vincoli attivi sono il carrello orizzontale in $B$
		e il carrello verticale fittizio in $D$: il \CRa\ del
		corpo è l'intersezione dei loro assi, cioè il punto
		$B$. Il cedimento verticale in $A$ produce una
		rotazione attorno a $B$ di ampiezza:
		\begin{equation*}
			\vartheta = \frac{-s_A}{\ell} \qquad
			\text{(oraria, dalla figura)}.
		\end{equation*}
		Il contributo alla soluzione è
		$\{-s_A,\, s_A,\, -s_A/\ell\}^\top$.
		
		\smallskip
		\noindent\textit{Cedimento $s_B$ (orizzontale in $B$).}
		I vincoli attivi sono il carrello verticale in $A$
		e il carrello verticale fittizio in $D$: i loro assi
		sono paralleli, quindi il \CRa\ è all'infinito in
		direzione verticale. Il corpo trasla orizzontalmente
		di $s_B$:
		\begin{equation*}
			u_O = s_B, \quad v_O = 0, \quad \vartheta = 0.
		\end{equation*}
		
		\smallskip
		\noindent\textit{Cedimento $q$ (verticale in $D$).}
		Entrambi i carrelli reali sono perfetti ($s_A = s_B =
		0$): il \CRa\ è il punto $C$, intersezione degli assi
		dei due carrelli. Il cedimento verticale in $D$ produce
		una rotazione attorno a $C$ di ampiezza:
		\begin{equation*}
			\vartheta = \frac{q}{\ell} \qquad
			\text{(antioraria, dalla figura)}.
		\end{equation*}
		Il contributo alla soluzione è
		$\{q,\, 0,\, q/\ell\}^\top$.
		
		
		
		\begin{figure} [ht]
			\centering
			\begin{tikzpicture}[>=Stealth, line width=0.8pt]
				
				\def\L{1.5}
				\def\sA{0.3}
				\def\sB{0.3}
				\def\q{0.3}
				\def\yH{1.0}
				
				% ============================================================
				% PANNELLO (a) — cedimento s_A
				% ============================================================
				\begin{scope}[xshift=0cm]
					
					% carrello verticale fittizio in D (grigio) — gomito della L
					\draw[gray!50, thick] (\L,0) -- (\L-0.21,-0.36)
					-- (\L+0.21,-0.36) -- cycle;
					\draw[gray!50, thick] (\L-0.14,-0.41) circle (0.06);
					\draw[gray!50, thick] (\L+0.14,-0.41) circle (0.06);
					\draw[gray!50, thick] (\L-0.29,-0.49) -- (\L+0.29,-0.49);
					\foreach \x in {-0.22,-0.11,0,0.11,0.22}{
						\draw[gray!50, thin] (\L+\x,-0.49) -- (\L+\x-0.09,-0.62);
					}
					% tre-quarti di luna grigio attorno al gomito (passaggio libero in alto-sinistra)
					\fill[white]          (\L,0.075) arc(90:-180:0.075) -- cycle;
					\draw[gray!50, thick] (\L,0.075) arc(90:-180:0.075);
					
					
					\draw[gray!40, thin] (-1.0,-0.3) -- (-1.0,2.0)
					node[above, font=\scriptsize] {$r_V$};
					\node[gray!60, left=2pt, font=\scriptsize]
					at (-1.0,1.0) {$u$};
					\draw[gray!40, thin] (-0.3,-\yH) -- (\L+0.3,-\yH)
					node[right, font=\scriptsize] {$r_H$};
					\node[gray!60, above=3pt, font=\scriptsize]
					at (0.6,-\yH) {$v$};
					
					\draw[gray!50, dashed, thin]
					(0,{-\yH-0.1}) -- (0,\L+0.3);
					
					\draw[gray!60, thick] (0,0) -- (\L,0) -- (\L,\L);
					\draw[line width=1.5pt]
					({-\sA},{\sA}) -- ({\L-\sA},0) -- (\L,\L);
					
					\fill (0,0)  %circle (2pt)
					node[left, font=\scriptsize] {$A$};
					\fill (\L,0) %circle (2pt)
					node[right, font=\scriptsize] {$D$};
					\fill (\L,\L) %circle (2pt)
					node[left, font=\scriptsize] {$C_1\equiv B$};
					
					
					% carrello orizzontale in B (in grigio: spettatore in questo pannello)
					\draw[gray!50, thick] (\L,\L) -- (\L+0.29,\L+0.17)
					-- (\L+0.29,\L-0.17) -- cycle;
					\draw[gray!50, thick] (\L+0.29,\L-0.22) -- (\L+0.29,\L+0.22);
					\draw[gray!50, thick] (\L+0.36,\L-0.10) circle (0.05);
					\draw[gray!50, thick] (\L+0.36,\L+0.10) circle (0.05);
					\draw[gray!50, thick] (\L+0.43,\L-0.22) -- (\L+0.43,\L+0.22);
					\foreach \y in {-0.18,-0.09,0,0.09,0.18}{
						\draw[gray!50, thin] (\L+0.43,\L+\y) -- (\L+0.55,\L+\y-0.08);
					}
					% perno grigio vuoto sopra il vertice del triangolo
					\draw[gray!50, fill=white, thick] (\L,\L) circle (2.5pt);
					
					% campo u su r_V
					\draw[gray!60, thin] ({-1.0-\sA},0) -- (-1.0,\L);
					\foreach \k in {0,1,2,3,4}{
						\pgfmathsetmacro{\yy}{\k*\L/4}
						\pgfmathsetmacro{\uu}{-\sA*(1-\k/4)}
						\fill[gray!60] (-1.0,\yy) circle (1pt);
						\ifnum\k<4
						\draw[->, gray!70, thin]
						(-1.0,\yy) -- ({-1.0+\uu},\yy);
						\fi
					}
					\draw[->, gray!70, thick] (-1.0,0) -- ({-1.0-\sA},0)
					node[below, font=\scriptsize] {$s_A$};
					
					% campo v su r_H
					\draw[gray!60, thin]
					(0,{-\yH+\sA}) -- (\L,{-\yH});
					\foreach \k in {0,1,2,3,4}{
						\pgfmathsetmacro{\xx}{\k*\L/4}
						\pgfmathsetmacro{\vv}{\sA*(1-\k/4)}
						\fill[gray!60] (\xx,-\yH) circle (1pt);
						\ifnum\k<4
						\draw[->, gray!70, thin]
						(\xx,-\yH) -- (\xx,{-\yH+\vv});
						\fi
					}
					\draw[->, gray!70, thick] (0,-\yH) -- (0,{-\yH+\sA})
					node[left, font=\scriptsize] {$s_A$};
					
					\node[font=\scriptsize] at (-\L/2,-1.4) {(a)};
					
				\end{scope}
				
				% ============================================================
				% PANNELLO (b) — cedimento s_B
				% ============================================================
				\begin{scope}[xshift=4.0cm]
					
					% carrello verticale fittizio in D (grigio) — gomito della L
					\draw[gray!50, thick] (\L,0) -- (\L-0.21,-0.36)
					-- (\L+0.21,-0.36) -- cycle;
					\draw[gray!50, thick] (\L-0.14,-0.41) circle (0.06);
					\draw[gray!50, thick] (\L+0.14,-0.41) circle (0.06);
					\draw[gray!50, thick] (\L-0.29,-0.49) -- (\L+0.29,-0.49);
					\foreach \x in {-0.22,-0.11,0,0.11,0.22}{
						\draw[gray!50, thin] (\L+\x,-0.49) -- (\L+\x-0.09,-0.62);
					}
					% tre-quarti di luna grigio attorno al gomito (passaggio libero in alto-sinistra)
					\fill[white]          (\L,0.075) arc(90:-180:0.075) -- cycle;
					\draw[gray!50, thick] (\L,0.075) arc(90:-180:0.075);
					
					\draw[gray!40, thin] (-1.0,-0.3) -- (-1.0,2.0)
					node[above, font=\scriptsize] {$r_V$};
					\node[gray!60, left=2pt, font=\scriptsize]
					at (-1.0,1.0) {$u$};
					
					\draw[gray!60, thick] (0,0) -- (\L,0) -- (\L,\L);
					\draw[line width=1.5pt]
					({\sB},0) -- ({\L+\sB},0) -- ({\L+\sB},\L);
					
					\fill (0,0)  node[left,  font=\scriptsize] {$A$};
					\fill (\L,0) node[above left, font=\scriptsize] {$D$};
					\fill (\L,\L) node[left, font=\scriptsize] {$B$};
					%\node[right, font=\scriptsize] at ({\L+\sB},\L) {$s_B$};
					
					
					% carrello verticale rimosso in A (tutto grigio)
					\draw[gray!50, thick] (0,0) -- (-0.17,-0.29)
					-- (0.17,-0.29) -- cycle;
					\draw[gray!50, thick] (-0.22,-0.29) -- (0.22,-0.29);
					\draw[gray!50, thick] (-0.10,-0.36) circle (0.05);
					\draw[gray!50, thick] ( 0.10,-0.36) circle (0.05);
					\draw[gray!50, thick] (-0.22,-0.43) -- (0.22,-0.43);
					\foreach \x in {-0.18,-0.09,0,0.09,0.18}{
						\draw[gray!50, thin] (\x,-0.43) -- (\x-0.08,-0.55);
					}
					% perno grigio vuoto sopra il vertice del triangolo
					\draw[gray!50, fill=white, thick] (0,0) circle (2.5pt);
					
					% centro di rotazione all'infinito
					\draw[->, gray!60, thick] (-1.0,1.7) -- (-1.0,2.0)
					node[left, font=\scriptsize] {$C_1^\infty$};
					
					% campo u: retta parallela a r_V, limitata a [0,L]
					\draw[gray!60, thin]
					({-1.0+\sB},0) -- ({-1.0+\sB},\L);
					\foreach \k in {0,1,2,3,4}{
						\pgfmathsetmacro{\yy}{\k*\L/4}
						\draw[->, gray!70, thin]
						(-1.0,\yy) -- ({-1.0+\sB},\yy);
					}
					% quota s_B
					%\draw[<->, gray!70, thin]	(-1.0,{\L+0.2}) -- ({-1.0+\sB},{\L+0.2});
					\node[gray!70, above, font=\scriptsize]
					at ({-0.85+\sB/2},{\L}) {$s_B$};
					
					\node[font=\scriptsize] at (-\L/2,-1.4) {(b)};
					
				\end{scope}
				
				% ============================================================
				% PANNELLO (c) — cedimento q
				% ============================================================
				\begin{scope}[xshift=8.0cm]
					
					
					\draw[gray!40, thin] (-1.0,-0.3) -- (-1.0,2.0)
					node[above, font=\scriptsize] {$r_V$};
					\node[gray!60, left=2pt, font=\scriptsize]
					at (-1.0,1.0) {$u$};
					\draw[gray!40, thin] (-0.3,-\yH) -- (\L+0.3,-\yH)
					node[right, font=\scriptsize] {$r_H$};
					\node[gray!60, above=2pt, font=\scriptsize]
					at (0.6,-\yH) {$v$};
					
					\draw[gray!50, dashed, thin]
					(0,{-\yH-0.1}) -- (0,\L+0.3);
					\draw[gray!50, dashed, thin]
					(-0.3,\L) -- (\L+0.3,\L);
					
					
					\fill (0,0)  %circle (2pt)
					node[left,  font=\scriptsize] {$A$};
					\fill (\L,0) %circle (2pt)
					node[right, font=\scriptsize] {$D$};
					\fill (\L,\L) %circle (2pt)
					node[below right, font=\scriptsize] {$B$};
					\fill (0,\L)  circle (2pt)
					node[left, font=\scriptsize] {$C$};
					
					
					% carrello verticale rimosso in A (tutto grigio)
					\draw[gray!50, thick] (0,0) -- (-0.17,-0.29)
					-- (0.17,-0.29) -- cycle;
					\draw[gray!50, thick] (-0.22,-0.29) -- (0.22,-0.29);
					\draw[gray!50, thick] (-0.10,-0.36) circle (0.05);
					\draw[gray!50, thick] ( 0.10,-0.36) circle (0.05);
					\draw[gray!50, thick] (-0.22,-0.43) -- (0.22,-0.43);
					\foreach \x in {-0.18,-0.09,0,0.09,0.18}{
						\draw[gray!50, thin] (\x,-0.43) -- (\x-0.08,-0.55);
					}
					% perno grigio vuoto sopra il vertice del triangolo
					\draw[gray!50, fill=white, thick] (0,0) circle (2.5pt);
					
					
					% carrello orizzontale in B (in grigio: spettatore in questo pannello)
					\draw[gray!50, thick] (\L,\L) -- (\L+0.29,\L+0.17)
					-- (\L+0.29,\L-0.17) -- cycle;
					\draw[gray!50, thick] (\L+0.29,\L-0.22) -- (\L+0.29,\L+0.22);
					\draw[gray!50, thick] (\L+0.36,\L-0.10) circle (0.05);
					\draw[gray!50, thick] (\L+0.36,\L+0.10) circle (0.05);
					\draw[gray!50, thick] (\L+0.43,\L-0.22) -- (\L+0.43,\L+0.22);
					\foreach \y in {-0.18,-0.09,0,0.09,0.18}{
						\draw[gray!50, thin] (\L+0.43,\L+\y) -- (\L+0.55,\L+\y-0.08);
					}
					% perno grigio vuoto sopra il vertice del triangolo
					\draw[gray!50, fill=white, thick] (\L,\L) circle (2.5pt);
					
					% configurazioni iniizale e deformata
					\draw[gray!60, thick] (0,0) -- (\L,0) -- (\L,\L);
					\draw[line width=1.5pt]
					({\q},0) -- ({\L+\q},{\q}) -- (\L,{\L+\q});
					
					% campo u su r_V
					\draw[gray!60, thin] ({-1.0+\q},0) -- (-1.0,\L);
					\foreach \k in {0,1,2,3,4}{
						\pgfmathsetmacro{\yy}{\k*\L/4}
						\pgfmathsetmacro{\uu}{\q*(1-\k/4)}
						\fill[gray!60] (-1.0,\yy) circle (1pt);
						\ifnum\k<4
						\draw[->, gray!70, thin]
						(-1.0,\yy) -- ({-1.0+\uu},\yy);
						\fi
					}
					\draw[->, gray!70, thick] (-1.0,0) -- ({-1.0+\q},0)
					node[below, font=\scriptsize] {$q$};
					
					% campo v su r_H
					\draw[gray!60, thin] (0,-\yH) -- (\L,{-\yH+\q});
					\foreach \k in {0,1,2,3,4}{
						\pgfmathsetmacro{\xx}{\k*\L/4}
						\pgfmathsetmacro{\vv}{\q*\k/4}
						\fill[gray!60] (\xx,-\yH) circle (1pt);
						\ifnum\k>0
						\draw[->, gray!70, thin]
						(\xx,-\yH) -- (\xx,{-\yH+\vv});
						\fi
					}
					\draw[->, gray!70, thick] (\L,-\yH) -- (\L,{-\yH+\q})
					node[right, font=\scriptsize] {$q$};
					
					\node[font=\scriptsize] at (-\L/2,-1.4) {(c)};
					
				\end{scope}
				
			\end{tikzpicture}
			\caption{Corpo a L (\ESEref{ex:labile_L}):
				contributi separati dei tre cedimenti alla
				configurazione variata, ottenuti per
				sovrapposizione.
				\emph{(a)}~Cedimento $s_A \neq 0$: rotazione
				oraria attorno a $B$, $\vartheta = -s_A/\ell$.
				\emph{(b)}~Cedimento $s_B \neq 0$: traslazione
				orizzontale di ampiezza $s_B$, centro di
				rotazione all'infinito.
				\emph{(c)}~Cedimento $q \neq 0$: rotazione
				antioraria attorno a $C$, $\vartheta = q/\ell$
				--- questo è il meccanismo del sistema labile.
				In ciascun pannello: vincoli attivi in nero,
				fittizi in grigio; rette ausiliarie $r_V$ e
				$r_H$ con i campi $u$ e $v$; configurazione
				variata in nero continuo.}
			\label{fig:labile_L_contributi}
		\end{figure}
		
		\smallskip
		La soluzione generale per sovrapposizione coincide con quella ottenuta analiticamente in precedenza, come atteso.
		
	\end{svolgimento}
\end{esempio}




\begin{esempio}[Tre scelte di vincoli fittizi per lo
	stesso sistema labile]
	\label{ex:scelte_fittizi}
	Si considerino due aste verticali di lunghezza $\ell$:
	$\mathcal{C}_1$ con estremi $A = (0,0)$ (cerniera al
	suolo) e $B = (0,\ell)$, e $\mathcal{C}_2$ con estremi
	$B$ (cerniera interna) e $C = (0, 2\ell)$ (estremità
	libera). Il sistema ha $n_c = 2$ corpi, $m = 4$
	vincoli scalari (2 dalla cerniera in $A$, 2 dalla
	cerniera interna in $B$) e $n = 6$ gradi di libertà:
	$\nu = 2$. Determinare tre scelte di vincoli fittizi
	che rendono il sistema isostatico e i corrispondenti
	meccanismi (\FIGref{fig:scelte_fittizi}).
	
	\begin{figure}[!ht]
		\centering
		\begin{tikzpicture}[>=Stealth, line width=0.8pt]
			
			\def\dx{4.2}  % distanza orizzontale tra pannelli
			\def\dy{3.8}  % distanza verticale tra righe
			
			% ============================================================
			% RIGA 1 — sistemi con vincoli fittizi
			% ============================================================
			\begin{scope}[yshift=0cm]
				% pannello (a)
				% ============================================================
				% PANNELLO (a) — Scelta 1: pendoli orizzontali in B e in C
				% ============================================================
				\begin{scope}[xshift=0cm]
					
					%% --- Parametri ---
					\def\ell{1.2}      % unità grafica (lunghezza dell'asta)
					\def\xast{-0.7}    % ascissa dell'asta verticale
					\def\Lpend{1.4}    % lunghezza orizzontale del pendolo
					\def\Lq{0.55}      % lunghezza della freccia di cedimento
					\def\dsuolo{0.18}
					\def\hsuolo{0.30}
					\def\hatch{0.13}
					
					%% --- Coordinate ---
					\pgfmathsetmacro{\yA}{-\ell}
					\pgfmathsetmacro{\yB}{0}
					\pgfmathsetmacro{\yC}{\ell}
					\pgfmathsetmacro{\xPend}{\xast+\Lpend}
					
					%% --- Aste verticali (in nero, due segmenti distinti) ---
					\draw[line width=1.2pt] (\xast,\yA) -- (\xast,\yB);  % C_1
					\draw[line width=1.2pt] (\xast,\yB) -- (\xast,\yC);  % C_2
					
					%% --- Pendoli orizzontali in B e in C (vincoli fittizi, grigio) ---
					\draw[gray!75, line width=0.8pt] (\xast,\yB) -- (\xPend,\yB);
					\draw[gray!75, line width=0.8pt] (\xast,\yC) -- (\xPend,\yC);
					
					%% --- Cerniera al suolo in A (nera, pavimento sotto) ---
					\draw[thick]
					(\xast,\yA) -- (\xast-\dsuolo,\yA-\hsuolo)
					-- (\xast+\dsuolo,\yA-\hsuolo) -- cycle;
					\draw[thick] (\xast-0.28,\yA-\hsuolo) -- (\xast+0.28,\yA-\hsuolo);
					\foreach \x in {-0.20,-0.10,0,0.10,0.20}{
						\draw[thin] (\xast+\x,\yA-\hsuolo)
						-- (\xast+\x-0.09,\yA-\hsuolo-\hatch);
					}
					%% --- Anellino in A (perno della cerniera al suolo) ---
					\node[circle, fill=white, draw=black, inner sep=0pt,
					minimum size=4pt, line width=0.7pt] at (\xast,\yA) {};
					
					
					%% --- Cerniera interna in B (nera, sopra l'asta) ---
					\node[circle, fill=white, draw=black, inner sep=0pt,
					minimum size=4pt, line width=0.7pt] at (\xast,\yB) {};
					
					%% --- Anellini di estremità dei pendoli ---
					\draw[gray!75, fill=white, line width=0.6pt] (\xast,\yC)  circle (1.6pt);
					\draw[gray!75, fill=white, line width=0.6pt] (\xPend,\yB) circle (1.6pt);
					\draw[gray!75, fill=white, line width=0.6pt] (\xPend,\yC) circle (1.6pt);
					
					%% --- Cerniera a muro per il pendolo in B (grigia, muro a destra) ---
					\draw[gray!75, thick]
					(\xPend,\yB) -- (\xPend+\hsuolo,\yB+\dsuolo)
					-- (\xPend+\hsuolo,\yB-\dsuolo) -- cycle;
					\draw[gray!75, thick] (\xPend+\hsuolo,\yB-0.28)
					-- (\xPend+\hsuolo,\yB+0.28);
					\foreach \y in {-0.20,-0.10,0,0.10,0.20}{
						\draw[gray!75, thin] (\xPend+\hsuolo,\yB+\y)
						-- (\xPend+\hsuolo+\hatch,\yB+\y-0.09);
					}
					
					%% --- Cerniera a muro per il pendolo in C (grigia, muro a destra) ---
					\draw[gray!75, thick]
					(\xPend,\yC) -- (\xPend+\hsuolo,\yC+\dsuolo)
					-- (\xPend+\hsuolo,\yC-\dsuolo) -- cycle;
					\draw[gray!75, thick] (\xPend+\hsuolo,\yC-0.28)
					-- (\xPend+\hsuolo,\yC+0.28);
					\foreach \y in {-0.20,-0.10,0,0.10,0.20}{
						\draw[gray!75, thin] (\xPend+\hsuolo,\yC+\y)
						-- (\xPend+\hsuolo+\hatch,\yC+\y-0.09);
					}
					
					%% --- Frecce di cedimento q_1 (in B) e q_2 (in C), in nero ---
					\draw[->, thick] (\xast,\yB) -- (\xast+\Lq,\yB);
					\draw[->, thick] (\xast,\yC) -- (\xast+\Lq,\yC);
					\node[font=\scriptsize, anchor=south]
					at ({\xast+\Lq/2}, \yB+0.05) {$q_1$};
					\node[font=\scriptsize, anchor=south]
					at ({\xast+\Lq/2}, \yC+0.05) {$q_2$};
					
					%% --- Etichette dei punti notevoli (a sinistra dell'asta) ---
					\node[font=\scriptsize, anchor=east] at (\xast-0.10,\yA) {$A$};
					\node[font=\scriptsize, anchor=east] at (\xast-0.10,\yB) {$B$};
					\node[font=\scriptsize, anchor=east] at (\xast-0.10,\yC) {$C$};
					
					%% --- Etichette dei corpi (a sinistra, mid-asta) ---
					\node[font=\scriptsize, anchor=east]
					at (\xast-0.10, {(\yA+\yB)/2}) {$\mathcal{C}_1$};
					\node[font=\scriptsize, anchor=east]
					at (\xast-0.10, {(\yB+\yC)/2}) {$\mathcal{C}_2$};
					
					%% --- Etichetta del pannello ---
					\node[font=\scriptsize] at (0.8,-1.4) {(a)};
					
				\end{scope}
				% ============================================================
				% PANNELLO (b) — Scelta 2: bipendolo su C_1 e bipendolo su C_2,
				% ============================================================
				\begin{scope}[xshift=\dx cm]
					
					%% --- Parametri ---
					\def\ell{1.2}
					\def\xast{-0.7}
					\def\Lb{0.55}        % lunghezza delle bielle (corte)
					\def\angBip{30}      % inclinazione delle bielle interne/esterne
					\def\dAtt{0.20}      % distanza tra i due anellini sull'asta
					\def\dsuolo{0.18}
					\def\hsuolo{0.30}
					\def\hatch{0.13}
					\def\rn{0.06}        % raggio degli anellini semicircolari
					\def\rArc{0.13}      % raggio degli archi di rotazione q_i
					
					%% --- Coordinate di riferimento ---
					\pgfmathsetmacro{\yA}{-\ell}
					\pgfmathsetmacro{\yB}{0}
					\pgfmathsetmacro{\yC}{\ell}
					\pgfmathsetmacro{\ymidUno}{(\yA+\yB)/2}
					\pgfmathsetmacro{\ymidDue}{(\yB+\yC)/2}
					
					%% --- Costanti geometriche del pantografo ---
					\pgfmathsetmacro{\dxBie}{\Lb*cos(\angBip)}
					\pgfmathsetmacro{\dyBie}{\Lb*sin(\angBip)}
					\pgfmathsetmacro{\xMuro}{\xast+2*\dxBie}
					
					%				%% --- Aste verticali (in nero) ---
					%				\draw[line width=1.2pt] (\xast,\yA) -- (\xast,\yB);
					%				\draw[line width=1.2pt] (\xast,\yB) -- (\xast,\yC);
					
					%% =========================================================
					%% Bipendolo su C_1
					%% =========================================================
					\pgfmathsetmacro{\yAiU}{\ymidUno+\dAtt/2}
					\pgfmathsetmacro{\yAiD}{\ymidUno-\dAtt/2}
					\pgfmathsetmacro{\yBiU}{\yAiU-\dyBie}
					\pgfmathsetmacro{\yBiD}{\yAiD-\dyBie}
					
					% bielle interne (parallelogramma 1)
					\draw[gray!75, line width=0.8pt] (\xast,\yAiU) -- ({\xast+\dxBie},\yBiU);
					\draw[gray!75, line width=0.8pt] (\xast,\yAiD) -- ({\xast+\dxBie},\yBiD);
					% traversa
					\draw[gray!75, line width=0.8pt] ({\xast+\dxBie},\yBiU) -- ({\xast+\dxBie},\yBiD);
					% bielle esterne (parallelogramma 2)
					\draw[gray!75, line width=0.8pt] ({\xast+\dxBie},\yBiU) -- (\xMuro,\yAiU);
					\draw[gray!75, line width=0.8pt] ({\xast+\dxBie},\yBiD) -- (\xMuro,\yAiD);
					% cerniere intermedie
					\draw[gray!75, fill=white, line width=0.6pt] ({\xast+\dxBie},\yBiU) circle (1.5pt);
					\draw[gray!75, fill=white, line width=0.6pt] ({\xast+\dxBie},\yBiD) circle (1.5pt);
					% anellini sul corpo C_1 (semicerchi sporgenti a destra)
					\draw[gray!75, fill=white, line width=0.6pt]
					(\xast,\yAiU-\rn) arc(-90:90:\rn) -- cycle;
					\draw[gray!75, fill=white, line width=0.6pt]
					(\xast,\yAiD-\rn) arc(-90:90:\rn) -- cycle;
					% anellini sul muro (semicerchi sporgenti a sinistra)
					\draw[gray!75, fill=white, line width=0.6pt]
					(\xMuro,\yAiU+\rn) arc(90:270:\rn) -- cycle;
					\draw[gray!75, fill=white, line width=0.6pt]
					(\xMuro,\yAiD+\rn) arc(90:270:\rn) -- cycle;
					% muro verticale con hatching
					\pgfmathsetmacro{\yMurMin}{\yAiD-0.12}
					\pgfmathsetmacro{\yMurMax}{\yAiU+0.12}
					\draw[gray!75, thick] (\xMuro,\yMurMin) -- (\xMuro,\yMurMax);
					\foreach \k in {0,1,2,3,4,5}{
						\pgfmathsetmacro{\yh}{\yMurMin+\k*(\yMurMax-\yMurMin)/5}
						\draw[gray!75, thin] (\xMuro,\yh) -- ({\xMuro+\hatch},\yh-0.09);
					}
					
					%% =========================================================
					%% Bipendolo su C_2  (variabili \yAiU, \yAiD, \yBiU, \yBiD ridefinite)
					%% =========================================================
					\pgfmathsetmacro{\yAiU}{\ymidDue+\dAtt/2}
					\pgfmathsetmacro{\yAiD}{\ymidDue-\dAtt/2}
					\pgfmathsetmacro{\yBiU}{\yAiU-\dyBie}
					\pgfmathsetmacro{\yBiD}{\yAiD-\dyBie}
					
					\draw[gray!75, line width=0.8pt] (\xast,\yAiU) -- ({\xast+\dxBie},\yBiU);
					\draw[gray!75, line width=0.8pt] (\xast,\yAiD) -- ({\xast+\dxBie},\yBiD);
					\draw[gray!75, line width=0.8pt] ({\xast+\dxBie},\yBiU) -- ({\xast+\dxBie},\yBiD);
					\draw[gray!75, line width=0.8pt] ({\xast+\dxBie},\yBiU) -- (\xMuro,\yAiU);
					\draw[gray!75, line width=0.8pt] ({\xast+\dxBie},\yBiD) -- (\xMuro,\yAiD);
					\draw[gray!75, fill=white, line width=0.6pt] ({\xast+\dxBie},\yBiU) circle (1.5pt);
					\draw[gray!75, fill=white, line width=0.6pt] ({\xast+\dxBie},\yBiD) circle (1.5pt);
					\draw[gray!75, fill=white, line width=0.6pt]
					(\xast,\yAiU-\rn) arc(-90:90:\rn) -- cycle;
					\draw[gray!75, fill=white, line width=0.6pt]
					(\xast,\yAiD-\rn) arc(-90:90:\rn) -- cycle;
					\draw[gray!75, fill=white, line width=0.6pt]
					(\xMuro,\yAiU+\rn) arc(90:270:\rn) -- cycle;
					\draw[gray!75, fill=white, line width=0.6pt]
					(\xMuro,\yAiD+\rn) arc(90:270:\rn) -- cycle;
					\pgfmathsetmacro{\yMurMin}{\yAiD-0.12}
					\pgfmathsetmacro{\yMurMax}{\yAiU+0.12}
					\draw[gray!75, thick] (\xMuro,\yMurMin) -- (\xMuro,\yMurMax);
					\foreach \k in {0,1,2,3,4,5}{
						\pgfmathsetmacro{\yh}{\yMurMin+\k*(\yMurMax-\yMurMin)/5}
						\draw[gray!75, thin] (\xMuro,\yh) -- ({\xMuro+\hatch},\yh-0.09);
					}
					
					%% --- Aste verticali (in nero) ---
					\draw[line width=1.2pt] (\xast,\yA) -- (\xast,\yB);
					\draw[line width=1.2pt] (\xast,\yB) -- (\xast,\yC);
					
					
					%% --- Cerniera al suolo in A (nero) ---
					\draw[thick]
					(\xast,\yA) -- (\xast-\dsuolo,\yA-\hsuolo)
					-- (\xast+\dsuolo,\yA-\hsuolo) -- cycle;
					\draw[thick] (\xast-0.28,\yA-\hsuolo) -- (\xast+0.28,\yA-\hsuolo);
					\foreach \x in {-0.20,-0.10,0,0.10,0.20}{
						\draw[thin] (\xast+\x,\yA-\hsuolo)
						-- (\xast+\x-0.09,\yA-\hsuolo-\hatch);
					}
					
					%% --- Anellini in A e B (nero) ---
					\node[circle, fill=white, draw=black, inner sep=0pt,
					minimum size=4pt, line width=0.7pt] at (\xast,\yA) {};
					\node[circle, fill=white, draw=black, inner sep=0pt,
					minimum size=4pt, line width=0.7pt] at (\xast,\yB) {};
					
					%				%% --- Frecce di rotazione q_1 e q_2 (nero), centrate sull'asta ---
					
					\def\qDeg{20}       % apertura grafica della rotazione
					\def\Lref{0.90}     % lunghezza dei due segmenti
					\def\rArcQ{0.70}    % raggio dell'archetto (minore di \Lref)
					\pgfmathsetmacro{\xRot}{\xMuro+0.55}   % a destra del muro
					
					% --- q_1: a fianco del bipendolo su C_1 ---
					\draw[thin] (\xRot,\ymidUno) -- (\xRot,{\ymidUno+\Lref});
					\draw[thin] (\xRot,\ymidUno) -- 
					({\xRot-\Lref*sin(\qDeg)},{\ymidUno+\Lref*cos(\qDeg)});
					\draw[->, thin] 
					(\xRot,{\ymidUno+\rArcQ}) 
					arc[start angle=90, end angle={90+\qDeg}, radius=\rArcQ];
					\node[font=\scriptsize] at ({\xRot+(\rArcQ+0.12)*cos(90+2*\qDeg)+0.1},{\ymidUno+(\rArcQ+0.12)*sin(90+\qDeg/2)-0.2}) {$q_1$};
					
					% --- q_2: a fianco del bipendolo su C_2 ---
					\draw[thin] (\xRot,\ymidDue) -- (\xRot,{\ymidDue+\Lref});
					\draw[thin] (\xRot,\ymidDue) -- 
					({\xRot-\Lref*sin(\qDeg)},{\ymidDue+\Lref*cos(\qDeg)});
					\draw[->, thin] 
					(\xRot,{\ymidDue+\rArcQ}) 
					arc[start angle=90, end angle={90+\qDeg}, radius=\rArcQ];
					\node[font=\scriptsize] at ({\xRot+(\rArcQ+0.12)*cos(90+2*\qDeg)+0.1},{\ymidDue+(\rArcQ+0.12)*sin(90+\qDeg/2)-0.2}) {$q_2$};
					
					%% --- Etichette ---
					\node[font=\scriptsize, anchor=east] at (\xast-0.10,\yA) {$A$};
					\node[font=\scriptsize, anchor=east] at (\xast-0.10,\yB) {$B$};
					\node[font=\scriptsize, anchor=east] at (\xast-0.10,\yC) {$C$};
					\node[font=\scriptsize, anchor=east]
					at (\xast-0.10, \ymidUno) {$\mathcal{C}_1$};
					\node[font=\scriptsize, anchor=east]
					at (\xast-0.10, \ymidDue) {$\mathcal{C}_2$};
					
					%% --- Etichetta del pannello ---
					\node[font=\scriptsize] at (0.8,-1.4) {(b)};
				\end{scope}
				% ============================================================
				% PANNELLO (c) — Scelta 3: bipendolo su C_1 al suolo
				%                          + bipendolo interno tra C_1 e C_2
				%                          (scavalca B, traversa orizzontale)
				% ============================================================
				\begin{scope}[xshift=2*\dx cm]
					
					%% --- Parametri comuni ---
					\def\ell{1.2}
					\def\xast{-0.7}
					\def\dsuolo{0.18}
					\def\hsuolo{0.30}
					\def\hatch{0.13}
					\def\rn{0.06}
					
					%% --- Coordinate dei punti notevoli ---
					\pgfmathsetmacro{\yA}{-\ell}
					\pgfmathsetmacro{\yB}{0}
					\pgfmathsetmacro{\yC}{\ell}
					
					%% =========================================================
					%% Bipendolo su C_1 al suolo (traslato verso il basso di Dymid)
					%% =========================================================
					\def\Lb{0.55}
					\def\angBip{30}
					\def\dAtt{0.20}
					\def\Dymid{-0.20}    % traslazione verticale rispetto al pannello (b)
					\pgfmathsetmacro{\ymidUno}{(\yA+\yB)/2 + \Dymid}
					\pgfmathsetmacro{\dxBie}{\Lb*cos(\angBip)}
					\pgfmathsetmacro{\dyBie}{\Lb*sin(\angBip)}
					\pgfmathsetmacro{\xMuro}{\xast+2*\dxBie}
					
					\pgfmathsetmacro{\yAiU}{\ymidUno+\dAtt/2}
					\pgfmathsetmacro{\yAiD}{\ymidUno-\dAtt/2}
					\pgfmathsetmacro{\yBiU}{\yAiU-\dyBie}
					\pgfmathsetmacro{\yBiD}{\yAiD-\dyBie}
					
					\draw[gray!75, line width=0.8pt] (\xast,\yAiU) -- ({\xast+\dxBie},\yBiU);
					\draw[gray!75, line width=0.8pt] (\xast,\yAiD) -- ({\xast+\dxBie},\yBiD);
					\draw[gray!75, line width=0.8pt] ({\xast+\dxBie},\yBiU) -- ({\xast+\dxBie},\yBiD);
					\draw[gray!75, line width=0.8pt] ({\xast+\dxBie},\yBiU) -- (\xMuro,\yAiU);
					\draw[gray!75, line width=0.8pt] ({\xast+\dxBie},\yBiD) -- (\xMuro,\yAiD);
					\draw[gray!75, fill=white, line width=0.6pt] ({\xast+\dxBie},\yBiU) circle (1.5pt);
					\draw[gray!75, fill=white, line width=0.6pt] ({\xast+\dxBie},\yBiD) circle (1.5pt);
					\draw[gray!75, fill=white, line width=0.6pt]
					(\xast,\yAiU-\rn) arc(-90:90:\rn) -- cycle;
					\draw[gray!75, fill=white, line width=0.6pt]
					(\xast,\yAiD-\rn) arc(-90:90:\rn) -- cycle;
					\draw[gray!75, fill=white, line width=0.6pt]
					(\xMuro,\yAiU+\rn) arc(90:270:\rn) -- cycle;
					\draw[gray!75, fill=white, line width=0.6pt]
					(\xMuro,\yAiD+\rn) arc(90:270:\rn) -- cycle;
					\pgfmathsetmacro{\yMurMin}{\yAiD-0.12}
					\pgfmathsetmacro{\yMurMax}{\yAiU+0.12}
					\draw[gray!75, thick] (\xMuro,\yMurMin) -- (\xMuro,\yMurMax);
					\foreach \k in {0,1,2,3,4,5}{
						\pgfmathsetmacro{\yh}{\yMurMin+\k*(\yMurMax-\yMurMin)/5}
						\draw[gray!75, thin] (\xMuro,\yh) -- ({\xMuro+\hatch},\yh-0.09);
					}
					
					%					%% =========================================================
					%					%% Bipendolo interno tra C_1 e C_2 
					%					%% =========================================================
					\def\dIn{0.25}         % distanza degli anellini interni da B (sull'asta)
					\def\dOut{0.40}        % distanza degli anellini esterni da B
					\def\angInt{30}        % inclinazione delle bielle dal positivo x (CCW)
					\pgfmathsetmacro{\cotInt}{cos(\angInt)/sin(\angInt)}
					
					% Quote dei 4 anellini, tutti a x = \xast
					\pgfmathsetmacro{\yAlowIn} {\yB-\dIn}     % su C_1, vicino a B
					\pgfmathsetmacro{\yAlowOut}{\yB-\dOut}    % su C_1, lontano da B
					\pgfmathsetmacro{\yAupIn}  {\yB+\dIn}     % su C_2, vicino a B
					\pgfmathsetmacro{\yAupOut} {\yB+\dOut}    % su C_2, lontano da B
					
					% Perni della traversa orizzontale (a y = yB)
					\pgfmathsetmacro{\xTin} {\xast+\dIn *\cotInt}
					\pgfmathsetmacro{\xTout}{\xast+\dOut*\cotInt}
					
					% Bielle interne (parallele, dirette in alto-destra) — corta + lunga
					\draw[gray!75, line width=0.8pt] (\xast,\yAlowIn)  -- (\xTin, \yB);
					\draw[gray!75, line width=0.8pt] (\xast,\yAlowOut) -- (\xTout,\yB);
					% Traversa orizzontale a y = yB
					\draw[gray!75, line width=0.8pt] (\xTin,\yB) -- (\xTout,\yB);
					% Bielle esterne (parallele, dirette in alto-sinistra) — corta + lunga
					\draw[gray!75, line width=0.8pt] (\xTin, \yB) -- (\xast,\yAupIn);
					\draw[gray!75, line width=0.8pt] (\xTout,\yB) -- (\xast,\yAupOut);
					% Cerniere sulla traversa
					\draw[gray!75, fill=white, line width=0.6pt] (\xTin, \yB) circle (1.5pt);
					\draw[gray!75, fill=white, line width=0.6pt] (\xTout,\yB) circle (1.5pt);
					
					% 4 anellini sulle aste AB e BC: semicerchi sporgenti a destra,
					% parte piatta appoggiata sull'asta verticale
					\draw[gray!75, fill=white, line width=0.6pt]
					(\xast,\yAlowIn -\rn) arc(-90:90:\rn) -- cycle;
					\draw[gray!75, fill=white, line width=0.6pt]
					(\xast,\yAlowOut-\rn) arc(-90:90:\rn) -- cycle;
					\draw[gray!75, fill=white, line width=0.6pt]
					(\xast,\yAupIn  -\rn) arc(-90:90:\rn) -- cycle;
					\draw[gray!75, fill=white, line width=0.6pt]
					(\xast,\yAupOut -\rn) arc(-90:90:\rn) -- cycle;
					
					%% =========================================================
					%% Aste verticali, cerniere, anellini reali (in nero, sopra)
					%% =========================================================
					\draw[line width=1.2pt] (\xast,\yA) -- (\xast,\yB);
					\draw[line width=1.2pt] (\xast,\yB) -- (\xast,\yC);
					
					\draw[thick]
					(\xast,\yA) -- (\xast-\dsuolo,\yA-\hsuolo)
					-- (\xast+\dsuolo,\yA-\hsuolo) -- cycle;
					\draw[thick] (\xast-0.28,\yA-\hsuolo) -- (\xast+0.28,\yA-\hsuolo);
					\foreach \x in {-0.20,-0.10,0,0.10,0.20}{
						\draw[thin] (\xast+\x,\yA-\hsuolo)
						-- (\xast+\x-0.09,\yA-\hsuolo-\hatch);
					}
					\node[circle, fill=white, draw=black, inner sep=0pt,
					minimum size=4pt, line width=0.7pt] at (\xast,\yA) {};
					\node[circle, fill=white, draw=black, inner sep=0pt,
					minimum size=4pt, line width=0.7pt] at (\xast,\yB) {};
					
					%% =========================================================
					%% Indicazione di q_1 (traslata insieme al bipendolo su C_1)
					%% =========================================================
					\def\qDeg{20}
					\def\Lref{0.90}
					\def\rArcQ{0.70}
					\def\Dyq{-0.30}        % traslazione verticale di q_1 (negativa = verso il basso)
					\pgfmathsetmacro{\xRot}{\xMuro+0.55}
					\pgfmathsetmacro{\yq}{\ymidUno+\Dyq}
					
					\draw[thin] (\xRot,\yq) -- (\xRot,{\yq+\Lref});
					\draw[thin] (\xRot,\yq) --
					({\xRot-\Lref*sin(\qDeg)},{\yq+\Lref*cos(\qDeg)});
					\draw[->, thin]
					(\xRot,{\yq+\rArcQ})
					arc[start angle=90, end angle={90+\qDeg}, radius=\rArcQ];
					\node[font=\scriptsize]
					at ({\xRot+(\rArcQ+0.12)*cos(90+2*\qDeg)+0.1},
					{\yq+(\rArcQ+0.12)*sin(90+\qDeg/2)-0.2}) {$q_1$};
					
					%% =========================================================
					%% Indicazione di q_2 (rotazione relativa C_2/C_1)
					%% =========================================================
					\def\xq{0.40}              % ascissa comune dei due semicerchi
					\def\dq{0.07}              % semi-distanza verticale fra le due frecce
					\def\rq{0.18}              % raggio dei semicerchi
					\pgfmathsetmacro{\yqup}{\yB+\dq}
					\pgfmathsetmacro{\yqdn}{\yB-\dq}
					
					% Semicerchio superiore: antiorario, convessità verso l'alto
					% (da destra, sopra il vertice, fino a sinistra — la freccia punta in basso)
					\draw[->, thin]
					({\xq+\rq},\yqup)
					arc[start angle=0, end angle=180, radius=\rq];
					
					% Semicerchio inferiore: orario, convessità verso il basso
					% (da destra, sotto il vertice, fino a sinistra — la freccia punta in alto)
					\draw[->, thin]
					({\xq+\rq},\yqdn)
					arc[start angle=0, end angle=-180, radius=\rq];
					
					% Etichetta q_2 a destra delle due frecce
					\node[font=\scriptsize] at ({\xq+\rq+0.20},\yB) {$q_2$};
					
					%% =========================================================
					%% Etichette
					%% =========================================================
					\node[font=\scriptsize, anchor=east] at (\xast-0.10,\yA) {$A$};
					\node[font=\scriptsize, anchor=east] at (\xast-0.10,\yB) {$B$};
					\node[font=\scriptsize, anchor=east] at (\xast-0.10,\yC) {$C$};
					\node[font=\scriptsize, anchor=east]
					at (\xast-0.10, {(\yA+\yB)/2}) {$\mathcal{C}_1$};
					\node[font=\scriptsize, anchor=east]
					at (\xast-0.10, {(\yB+\yC)/2}) {$\mathcal{C}_2$};
					
					%% --- Etichetta del pannello ---
					\node[font=\small] at (0.8,-1.4) {(c)};
					
				\end{scope}
			\end{scope}
			
			% ============================================================
			% RIGA 2 — meccanismi q_1
			% ============================================================
			\begin{scope}[yshift=-\dy cm]
				% pannello (d)
				\begin{scope}[xshift=0cm]
					
					%% --- Parametri ---
					\def\ell{1.2}      % unità grafica (lunghezza dell'asta)
					\def\xast{-0.7}    % ascissa dell'asta verticale
					\def\Lpend{1.4}    % lunghezza orizzontale del pendolo
					\def\Lq{0.55}      % lunghezza della freccia di cedimento
					\def\dsuolo{0.18}
					\def\hsuolo{0.30}
					\def\hatch{0.13}
					
					%% --- Coordinate ---
					\pgfmathsetmacro{\yA}{-\ell}
					\pgfmathsetmacro{\yB}{0}
					\pgfmathsetmacro{\yC}{\ell}
					\pgfmathsetmacro{\xPend}{\xast+\Lpend}
					
					%% --- Aste verticali (in nero, due segmenti distinti) ---
					\draw[line width=1.2pt] (\xast,\yA) -- (\xast,\yB);  % C_1
					\draw[line width=1.2pt] (\xast,\yB) -- (\xast,\yC);  % C_2
					
					%% --- Pendoli orizzontali in B e in C (vincoli fittizi, grigio) ---
					%\draw[gray!75, line width=0.8pt] (\xast,\yB) -- (\xPend,\yB);
					\draw[gray!75, line width=0.8pt] (\xast,\yC) -- (\xPend,\yC);
					
					%% --- Cerniera al suolo in A (nera, pavimento sotto) ---
					\draw[thick]
					(\xast,\yA) -- (\xast-\dsuolo,\yA-\hsuolo)
					-- (\xast+\dsuolo,\yA-\hsuolo) -- cycle;
					\draw[thick] (\xast-0.28,\yA-\hsuolo) -- (\xast+0.28,\yA-\hsuolo);
					\foreach \x in {-0.20,-0.10,0,0.10,0.20}{
						\draw[thin] (\xast+\x,\yA-\hsuolo)
						-- (\xast+\x-0.09,\yA-\hsuolo-\hatch);
					}
					%% --- Anellino in A (perno della cerniera al suolo) ---
					\node[circle, fill=white, draw=black, inner sep=0pt,
					minimum size=4pt, line width=0.7pt] at (\xast,\yA) {};
					
					
					%% --- Cerniera interna in B (nera, sopra l'asta) ---
					\node[circle, fill=white, draw=black, inner sep=0pt,
					minimum size=4pt, line width=0.7pt] at (\xast,\yB) {};
					
					%% --- Anellini di estremità dei pendoli ---
					\draw[gray!75, fill=white, line width=0.6pt] (\xast,\yC)  circle (1.6pt);
					%\draw[gray!75, fill=white, line width=0.6pt] (\xPend,\yB) circle (1.6pt);
					\draw[gray!75, fill=white, line width=0.6pt] (\xPend,\yC) circle (1.6pt);
					
					%% --- Cerniera a muro per il pendolo in C (grigia, muro a destra) ---
					\draw[gray!75, thick]
					(\xPend,\yC) -- (\xPend+\hsuolo,\yC+\dsuolo)
					-- (\xPend+\hsuolo,\yC-\dsuolo) -- cycle;
					\draw[gray!75, thick] (\xPend+\hsuolo,\yC-0.28)
					-- (\xPend+\hsuolo,\yC+0.28);
					\foreach \y in {-0.20,-0.10,0,0.10,0.20}{
						\draw[gray!75, thin] (\xPend+\hsuolo,\yC+\y)
						-- (\xPend+\hsuolo+\hatch,\yC+\y-0.09);
					}
					
					%% --- Frecce di cedimento q_1 (in B) e q_2 (in C), in nero ---
					\draw[->, thick] (\xast,\yB) -- (\xast+\Lq,\yB);
					%\draw[->, thick] (\xast,\yC) -- (\xast+\Lq,\yC);
					\node[font=\scriptsize, anchor=south]	at ({\xast+\Lq/2}, \yB+0.05) {$q_1$};
					%\node[font=\scriptsize, anchor=south]	at ({\xast+\Lq/2}, \yC+0.05) {$q_2$};
					
					%% --- Etichette dei punti notevoli (a sinistra dell'asta) ---
					\node[font=\scriptsize, anchor=east] at (\xast-0.10,\yA) {$A$};
					\node[font=\scriptsize, anchor=east] at (\xast-0.10,\yB) {$B$};
					\node[font=\scriptsize, anchor=east] at (\xast-0.10,\yC) {$C$};
					
					%% --- Etichette dei corpi (a sinistra, mid-asta) ---
					\node[font=\scriptsize, anchor=east]
					at (\xast-0.10, {(\yA+\yB)/2}) {$\mathcal{C}_1$};
					\node[font=\scriptsize, anchor=east]
					at (\xast-0.10, {(\yB+\yC)/2}) {$\mathcal{C}_2$};
					
					%% --- Deformata in grigio: A--B' e C--B' ---
					%% B' è l'estremo della freccia q_1 in B
					\pgfmathsetmacro{\xBp}{\xast+\Lq}
					\pgfmathsetmacro{\yBp}{\yB}
					
					% Segmento A -- B' (deformata di C_1)
					\draw[gray!75, line width=1.0pt] (\xast,\yA) -- (\xBp,\yBp);
					
					% Segmento C -- B' (deformata di C_2)
					\draw[gray!75, line width=1.0pt] (\xast,\yC) -- (\xBp,\yBp);
					
					% Etichetta B' a destra del punto
					\node[gray!75, font=\scriptsize, anchor=west] at (\xBp+0.05,\yBp) {$B'$};
					
					%% --- Anellino grigio in B' ---
					\draw[gray!75, fill=white, line width=0.6pt] (\xBp,\yBp) circle (1.6pt);
					
					%% --- Etichetta del pannello ---
					\node[font=\scriptsize] at (0.8,-1.4) {(d)};
					
				\end{scope}
				% pannello (e)
				\begin{scope}[xshift=\dx cm]
					
					%% --- Parametri ---
					\def\ell{1.2}
					\def\xast{-0.7}
					\def\Lb{0.55}        % lunghezza delle bielle (corte)
					\def\angBip{30}      % inclinazione delle bielle interne/esterne
					\def\dAtt{0.20}      % distanza tra i due anellini sull'asta
					\def\dsuolo{0.18}
					\def\hsuolo{0.30}
					\def\hatch{0.13}
					\def\rn{0.06}        % raggio degli anellini semicircolari
					\def\rArc{0.13}      % raggio degli archi di rotazione q_i
					
					%% --- Coordinate di riferimento ---
					\pgfmathsetmacro{\yA}{-\ell}
					\pgfmathsetmacro{\yB}{0}
					\pgfmathsetmacro{\yC}{\ell}
					\pgfmathsetmacro{\ymidUno}{(\yA+\yB)/2}
					\pgfmathsetmacro{\ymidDue}{(\yB+\yC)/2}
					
					%% --- Costanti geometriche del pantografo ---
					\pgfmathsetmacro{\dxBie}{\Lb*cos(\angBip)}
					\pgfmathsetmacro{\dyBie}{\Lb*sin(\angBip)}
					\pgfmathsetmacro{\xMuro}{\xast+2*\dxBie}
					
					
					%% =========================================================
					%% Bipendolo su C_2  (variabili \yAiU, \yAiD, \yBiU, \yBiD ridefinite)
					%% =========================================================
					\pgfmathsetmacro{\yAiU}{\ymidDue+\dAtt/2}
					\pgfmathsetmacro{\yAiD}{\ymidDue-\dAtt/2}
					\pgfmathsetmacro{\yBiU}{\yAiU-\dyBie}
					\pgfmathsetmacro{\yBiD}{\yAiD-\dyBie}
					
					\draw[gray!75, line width=0.8pt] (\xast,\yAiU) -- ({\xast+\dxBie},\yBiU);
					\draw[gray!75, line width=0.8pt] (\xast,\yAiD) -- ({\xast+\dxBie},\yBiD);
					\draw[gray!75, line width=0.8pt] ({\xast+\dxBie},\yBiU) -- ({\xast+\dxBie},\yBiD);
					\draw[gray!75, line width=0.8pt] ({\xast+\dxBie},\yBiU) -- (\xMuro,\yAiU);
					\draw[gray!75, line width=0.8pt] ({\xast+\dxBie},\yBiD) -- (\xMuro,\yAiD);
					\draw[gray!75, fill=white, line width=0.6pt] ({\xast+\dxBie},\yBiU) circle (1.5pt);
					\draw[gray!75, fill=white, line width=0.6pt] ({\xast+\dxBie},\yBiD) circle (1.5pt);
					\draw[gray!75, fill=white, line width=0.6pt]
					(\xast,\yAiU-\rn) arc(-90:90:\rn) -- cycle;
					\draw[gray!75, fill=white, line width=0.6pt]
					(\xast,\yAiD-\rn) arc(-90:90:\rn) -- cycle;
					\draw[gray!75, fill=white, line width=0.6pt]
					(\xMuro,\yAiU+\rn) arc(90:270:\rn) -- cycle;
					\draw[gray!75, fill=white, line width=0.6pt]
					(\xMuro,\yAiD+\rn) arc(90:270:\rn) -- cycle;
					\pgfmathsetmacro{\yMurMin}{\yAiD-0.12}
					\pgfmathsetmacro{\yMurMax}{\yAiU+0.12}
					\draw[gray!75, thick] (\xMuro,\yMurMin) -- (\xMuro,\yMurMax);
					\foreach \k in {0,1,2,3,4,5}{
						\pgfmathsetmacro{\yh}{\yMurMin+\k*(\yMurMax-\yMurMin)/5}
						\draw[gray!75, thin] (\xMuro,\yh) -- ({\xMuro+\hatch},\yh-0.09);
					}
					
					%% --- Aste verticali (in nero) ---
					\draw[line width=1.2pt] (\xast,\yA) -- (\xast,\yB);
					\draw[line width=1.2pt] (\xast,\yB) -- (\xast,\yC);
					
					
					%% --- Cerniera al suolo in A (nero) ---
					\draw[thick]
					(\xast,\yA) -- (\xast-\dsuolo,\yA-\hsuolo)
					-- (\xast+\dsuolo,\yA-\hsuolo) -- cycle;
					\draw[thick] (\xast-0.28,\yA-\hsuolo) -- (\xast+0.28,\yA-\hsuolo);
					\foreach \x in {-0.20,-0.10,0,0.10,0.20}{
						\draw[thin] (\xast+\x,\yA-\hsuolo)
						-- (\xast+\x-0.09,\yA-\hsuolo-\hatch);
					}
					
					%% --- Anellini in A e B (nero) ---
					\node[circle, fill=white, draw=black, inner sep=0pt,
					minimum size=4pt, line width=0.7pt] at (\xast,\yA) {};
					\node[circle, fill=white, draw=black, inner sep=0pt,
					minimum size=4pt, line width=0.7pt] at (\xast,\yB) {};
					
					%				%% --- Frecce di rotazione q_1 e q_2 (nero), centrate sull'asta ---
					
					\def\qDeg{20}       % apertura grafica della rotazione
					\def\Lref{0.90}     % lunghezza dei due segmenti
					\def\rArcQ{0.70}    % raggio dell'archetto (minore di \Lref)
					\pgfmathsetmacro{\xRot}{\xMuro+0.55}   % a destra del muro
					
					% --- q_1: a fianco del bipendolo su C_1 ---
					\draw[thin] (\xRot,\ymidUno) -- (\xRot,{\ymidUno+\Lref});
					\draw[thin] (\xRot,\ymidUno) -- 
					({\xRot-\Lref*sin(\qDeg)},{\ymidUno+\Lref*cos(\qDeg)});
					\draw[->, thin] 
					(\xRot,{\ymidUno+\rArcQ}) 
					arc[start angle=90, end angle={90+\qDeg}, radius=\rArcQ];
					\node[font=\scriptsize] at ({\xRot+(\rArcQ+0.12)*cos(90+2*\qDeg)+0.1},{\ymidUno+(\rArcQ+0.12)*sin(90+\qDeg/2)-0.2}) {$q_1$};
					
					%% --- Etichette ---
					\node[font=\scriptsize, anchor=east] at (\xast-0.10,\yA) {$A$};
					\node[font=\scriptsize, anchor=east] at (\xast-0.10,\yB) {$B$};
					\node[font=\scriptsize, anchor=east] at (\xast-0.01,\yC) {$C$};
					%\node[font=\scriptsize, anchor=south] at (\xast+0.06,\yC) {$C$};
					\node[font=\scriptsize, anchor=east]
					at (\xast-0.10, \ymidUno) {$\mathcal{C}_1$};
					\node[font=\scriptsize, anchor=east]
					at (\xast-0.10, \ymidDue) {$\mathcal{C}_2$};
					
					%% --- Deformata in grigio: A -- B' (rotazione di C_1) e B' -- C' (traslazione di C_2) ---
					\def\Lq{0.55}
					\pgfmathsetmacro{\xBp}{\xast-\Lq}
					\pgfmathsetmacro{\yBp}{\yB}
					\pgfmathsetmacro{\xCp}{\xast-\Lq}
					\pgfmathsetmacro{\yCp}{\yC}
					
					% Segmento A -- B' (deformata di C_1)
					\draw[gray!75, line width=1.0pt] (\xast,\yA) -- (\xBp,\yBp);
					
					% Segmento B' -- C' (deformata di C_2, verticale traslata)
					\draw[gray!75, line width=1.0pt] (\xBp,\yBp) -- (\xCp,\yCp);
					
					% Anellini grigi in B' e C'
					\draw[gray!75, fill=white, line width=0.6pt] (\xBp,\yBp) circle (1.6pt);
					\draw[gray!75, fill=white, line width=0.6pt] (\xCp,\yCp) circle (1.6pt);
					
					% Etichette B' e C'
					\node[gray!75, font=\scriptsize, anchor=east] at (\xBp-0.05,\yBp) {$B'$};
					\node[gray!75, font=\scriptsize, anchor=east] at (\xCp-0.05,\yCp) {$C'$};
					
					%% --- Archetto della rotazione q_1 (antioraria, attorno ad A) ---
					\def\qDegI{22}        % apertura grafica della rotazione
					\def\rArcI{0.85}      % raggio dell'archetto (< ell)
					
					\draw[->, gray!75, thin]
					(\xast,{\yA+\rArcI})
					arc[start angle=90, end angle={90+\qDegI}, radius=\rArcI];
					
					% Etichetta q_1 al punto medio dell'arco, sul lato esterno
					\node[gray!75, font=\scriptsize, anchor=east] at ({\xast+\rArcI*cos(90+\qDegI/2)-0.03},	{\yA+\rArcI*sin(90+\qDegI/2)})
					{$q_1$};
					
					%% --- Etichetta del pannello ---
					\node[font=\scriptsize] at (0.8,-1.4) {(e)};
				\end{scope}
				% pannello (f)
				\begin{scope}[xshift=2*\dx cm]
					
					%% --- Parametri comuni ---
					\def\ell{1.2}
					\def\xast{-0.7}
					\def\dsuolo{0.18}
					\def\hsuolo{0.30}
					\def\hatch{0.13}
					\def\rn{0.06}
					
					%% --- Coordinate dei punti notevoli ---
					\pgfmathsetmacro{\yA}{-\ell}
					\pgfmathsetmacro{\yB}{0}
					\pgfmathsetmacro{\yC}{\ell}
					
					%% =========================================================
					%% Bipendolo su C_1 al suolo (traslato verso il basso di Dymid)
					%% =========================================================
					\def\Lb{0.55}
					\def\angBip{30}
					\def\dAtt{0.20}
					\def\Dymid{-0.20}    % traslazione verticale rispetto al pannello (b)
					\pgfmathsetmacro{\ymidUno}{(\yA+\yB)/2 + \Dymid}
					\pgfmathsetmacro{\dxBie}{\Lb*cos(\angBip)}
					\pgfmathsetmacro{\dyBie}{\Lb*sin(\angBip)}
					\pgfmathsetmacro{\xMuro}{\xast+2*\dxBie}
					
					\pgfmathsetmacro{\yAiU}{\ymidUno+\dAtt/2}
					\pgfmathsetmacro{\yAiD}{\ymidUno-\dAtt/2}
					\pgfmathsetmacro{\yBiU}{\yAiU-\dyBie}
					\pgfmathsetmacro{\yBiD}{\yAiD-\dyBie}
					
					
					%					%% =========================================================
					%					%% Bipendolo interno tra C_1 e C_2 
					%					%% =========================================================
					\def\dIn{0.25}         % distanza degli anellini interni da B (sull'asta)
					\def\dOut{0.40}        % distanza degli anellini esterni da B
					\def\angInt{30}        % inclinazione delle bielle dal positivo x (CCW)
					\pgfmathsetmacro{\cotInt}{cos(\angInt)/sin(\angInt)}
					
					% Quote dei 4 anellini, tutti a x = \xast
					\pgfmathsetmacro{\yAlowIn} {\yB-\dIn}     % su C_1, vicino a B
					\pgfmathsetmacro{\yAlowOut}{\yB-\dOut}    % su C_1, lontano da B
					\pgfmathsetmacro{\yAupIn}  {\yB+\dIn}     % su C_2, vicino a B
					\pgfmathsetmacro{\yAupOut} {\yB+\dOut}    % su C_2, lontano da B
					
					% Perni della traversa orizzontale (a y = yB)
					\pgfmathsetmacro{\xTin} {\xast+\dIn *\cotInt}
					\pgfmathsetmacro{\xTout}{\xast+\dOut*\cotInt}
					
					% Bielle interne (parallele, dirette in alto-destra) — corta + lunga
					\draw[gray!75, line width=0.8pt] (\xast,\yAlowIn)  -- (\xTin, \yB);
					\draw[gray!75, line width=0.8pt] (\xast,\yAlowOut) -- (\xTout,\yB);
					% Traversa orizzontale a y = yB
					\draw[gray!75, line width=0.8pt] (\xTin,\yB) -- (\xTout,\yB);
					% Bielle esterne (parallele, dirette in alto-sinistra) — corta + lunga
					\draw[gray!75, line width=0.8pt] (\xTin, \yB) -- (\xast,\yAupIn);
					\draw[gray!75, line width=0.8pt] (\xTout,\yB) -- (\xast,\yAupOut);
					% Cerniere sulla traversa
					\draw[gray!75, fill=white, line width=0.6pt] (\xTin, \yB) circle (1.5pt);
					\draw[gray!75, fill=white, line width=0.6pt] (\xTout,\yB) circle (1.5pt);
					
					% 4 anellini sulle aste AB e BC: semicerchi sporgenti a destra,
					% parte piatta appoggiata sull'asta verticale
					\draw[gray!75, fill=white, line width=0.6pt]
					(\xast,\yAlowIn -\rn) arc(-90:90:\rn) -- cycle;
					\draw[gray!75, fill=white, line width=0.6pt]
					(\xast,\yAlowOut-\rn) arc(-90:90:\rn) -- cycle;
					\draw[gray!75, fill=white, line width=0.6pt]
					(\xast,\yAupIn  -\rn) arc(-90:90:\rn) -- cycle;
					\draw[gray!75, fill=white, line width=0.6pt]
					(\xast,\yAupOut -\rn) arc(-90:90:\rn) -- cycle;
					
					%% =========================================================
					%% Aste verticali, cerniere, anellini reali (in nero, sopra)
					%% =========================================================
					\draw[line width=1.2pt] (\xast,\yA) -- (\xast,\yB);
					\draw[line width=1.2pt] (\xast,\yB) -- (\xast,\yC);
					
					\draw[thick]
					(\xast,\yA) -- (\xast-\dsuolo,\yA-\hsuolo)
					-- (\xast+\dsuolo,\yA-\hsuolo) -- cycle;
					\draw[thick] (\xast-0.28,\yA-\hsuolo) -- (\xast+0.28,\yA-\hsuolo);
					\foreach \x in {-0.20,-0.10,0,0.10,0.20}{
						\draw[thin] (\xast+\x,\yA-\hsuolo)
						-- (\xast+\x-0.09,\yA-\hsuolo-\hatch);
					}
					\node[circle, fill=white, draw=black, inner sep=0pt,
					minimum size=4pt, line width=0.7pt] at (\xast,\yA) {};
					\node[circle, fill=white, draw=black, inner sep=0pt,
					minimum size=4pt, line width=0.7pt] at (\xast,\yB) {};
					
					%% =========================================================
					%% Indicazione di q_1 (traslata insieme al bipendolo su C_1)
					%% =========================================================
					\def\qDeg{20}
					\def\Lref{0.90}
					\def\rArcQ{0.70}
					\def\Dyq{-0.30}        % traslazione verticale di q_1 (negativa = verso il basso)
					\pgfmathsetmacro{\xRot}{\xMuro+0.55}
					\pgfmathsetmacro{\yq}{\ymidUno+\Dyq}
					
					\draw[thin] (\xRot,\yq) -- (\xRot,{\yq+\Lref});
					\draw[thin] (\xRot,\yq) --
					({\xRot-\Lref*sin(\qDeg)},{\yq+\Lref*cos(\qDeg)});
					\draw[->, thin]
					(\xRot,{\yq+\rArcQ})
					arc[start angle=90, end angle={90+\qDeg}, radius=\rArcQ];
					\node[font=\scriptsize]
					at ({\xRot+(\rArcQ+0.12)*cos(90+2*\qDeg)+0.1},
					{\yq+(\rArcQ+0.12)*sin(90+\qDeg/2)-0.2}) {$q_1$};
					
					
					%% =========================================================
					%% Etichette
					%% =========================================================
					\node[font=\scriptsize, anchor=east] at (\xast-0.10,\yA) {$A$};
					\node[font=\scriptsize, anchor=east] at (\xast-0.01,\yB) {$B$};
					\node[font=\scriptsize, anchor=east] at (\xast-0.01,\yC) {$C$};
					\node[font=\scriptsize, anchor=west] at (\xast-0.10, {(\yA+\yB)/2}) {$\mathcal{C}_1$};
					\node[font=\scriptsize, anchor=west] at (\xast-0.10, {(\yB+\yC)/2}) {$\mathcal{C}_2$};
					
					%% --- Deformata in grigio: A -- B' -- C' (rotazione rigida attorno ad A) ---
					\def\Lq{0.55}
					\pgfmathsetmacro{\xBp}{\xast-\Lq}
					\pgfmathsetmacro{\yBp}{\yB}
					\pgfmathsetmacro{\xCp}{\xast-2*\Lq}
					\pgfmathsetmacro{\yCp}{\yC}
					
					% Segmento A -- B' -- C' (deformata di C_1 e C_2 allineate)
					\draw[gray!75, line width=1.0pt] (\xast,\yA) -- (\xBp,\yBp) -- (\xCp,\yCp);
					
					% Anellini grigi in B' e C'
					\draw[gray!75, fill=white, line width=0.6pt] (\xBp,\yBp) circle (1.6pt);
					\draw[gray!75, fill=white, line width=0.6pt] (\xCp,\yCp) circle (1.6pt);
					
					% Etichette B' e C'
					\node[gray!75, font=\scriptsize, anchor=east] at (\xBp-0.0,\yBp) {$B'$};
					\node[gray!75, font=\scriptsize, anchor=east] at (\xCp-0.0,\yCp) {$C'$};
					
					%% --- Archetto della rotazione q_1 (antioraria, attorno ad A) ---
					\def\qDegI{22}        % apertura grafica della rotazione
					\def\rArcI{0.85}      % raggio dell'archetto (< ell)
					
					\draw[->, gray!75, thin]
					(\xast,{\yA+\rArcI})
					arc[start angle=90, end angle={90+\qDegI}, radius=\rArcI];
					
					% Etichetta q_1 al punto medio dell'arco, sul lato esterno
					\node[gray!75, font=\scriptsize, anchor=east]	at ({\xast+\rArcI*cos(90+\qDegI/2)-0.03}, {\yA-0.1+\rArcI*sin(90+\qDegI/2)})
					{$q_1$};
					
					%% --- Etichetta del pannello ---
					\node[font=\scriptsize] at (0.8,-1.4) {(f)};
					
				\end{scope}
			\end{scope}
			
			% ============================================================
			% RIGA 3 — meccanismi q_2
			% ============================================================
			\begin{scope}[yshift=-2*\dy cm]
				% pannello (g)
				\begin{scope}[xshift=0cm]
					
					%% --- Parametri ---
					\def\ell{1.2}      % unità grafica (lunghezza dell'asta)
					\def\xast{-0.7}    % ascissa dell'asta verticale
					\def\Lpend{1.4}    % lunghezza orizzontale del pendolo
					\def\Lq{0.55}      % lunghezza della freccia di cedimento
					\def\dsuolo{0.18}
					\def\hsuolo{0.30}
					\def\hatch{0.13}
					
					%% --- Coordinate ---
					\pgfmathsetmacro{\yA}{-\ell}
					\pgfmathsetmacro{\yB}{0}
					\pgfmathsetmacro{\yC}{\ell}
					\pgfmathsetmacro{\xPend}{\xast+\Lpend}
					
					%% --- Aste verticali (in nero, due segmenti distinti) ---
					\draw[line width=1.2pt] (\xast,\yA) -- (\xast,\yB);  % C_1
					\draw[line width=1.2pt] (\xast,\yB) -- (\xast,\yC);  % C_2
					
					%% --- Pendoli orizzontali in B e in C (vincoli fittizi, grigio) ---
					\draw[gray!75, line width=0.8pt] (\xast,\yB) -- (\xPend,\yB);
					%\draw[gray!75, line width=0.8pt] (\xast,\yC) -- (\xPend,\yC);
					
					%% --- Cerniera al suolo in A (nera, pavimento sotto) ---
					\draw[thick]
					(\xast,\yA) -- (\xast-\dsuolo,\yA-\hsuolo)
					-- (\xast+\dsuolo,\yA-\hsuolo) -- cycle;
					\draw[thick] (\xast-0.28,\yA-\hsuolo) -- (\xast+0.28,\yA-\hsuolo);
					\foreach \x in {-0.20,-0.10,0,0.10,0.20}{
						\draw[thin] (\xast+\x,\yA-\hsuolo)
						-- (\xast+\x-0.09,\yA-\hsuolo-\hatch);
					}
					%% --- Anellino in A (perno della cerniera al suolo) ---
					\node[circle, fill=white, draw=black, inner sep=0pt,
					minimum size=4pt, line width=0.7pt] at (\xast,\yA) {};
					
					
					%% --- Cerniera interna in B (nera, sopra l'asta) ---
					\node[circle, fill=white, draw=black, inner sep=0pt,
					minimum size=4pt, line width=0.7pt] at (\xast,\yB) {};
					
					%% --- Anellini di estremità dei pendoli ---
					\draw[gray!75, fill=white, line width=0.6pt] (\xast,\yC)  circle (1.6pt);
					\draw[gray!75, fill=white, line width=0.6pt] (\xPend,\yB) circle (1.6pt);
					%\draw[gray!75, fill=white, line width=0.6pt] (\xPend,\yC) circle (1.6pt);
					
					%% --- Cerniera a muro per il pendolo in B (grigia, muro a destra) ---
					\draw[gray!75, thick]
					(\xPend,\yB) -- (\xPend+\hsuolo,\yB+\dsuolo)
					-- (\xPend+\hsuolo,\yB-\dsuolo) -- cycle;
					\draw[gray!75, thick] (\xPend+\hsuolo,\yB-0.28)
					-- (\xPend+\hsuolo,\yB+0.28);
					\foreach \y in {-0.20,-0.10,0,0.10,0.20}{
						\draw[gray!75, thin] (\xPend+\hsuolo,\yB+\y)
						-- (\xPend+\hsuolo+\hatch,\yB+\y-0.09);
					}
					
					%% --- Frecce di cedimento q_1 (in B) e q_2 (in C), in nero ---
					%\draw[->, thick] (\xast,\yB) -- (\xast+\Lq,\yB);
					\draw[->, thick] (\xast,\yC) -- (\xast+\Lq,\yC);
					%\node[font=\scriptsize, anchor=south]	at ({\xast+\Lq/2}, \yB+0.05) {$q_1$};
					\node[font=\scriptsize, anchor=south]	at ({\xast+\Lq/2}, \yC+0.05) {$q_2$};
					
					%% --- Etichette dei punti notevoli (a sinistra dell'asta) ---
					\node[font=\scriptsize, anchor=east] at (\xast-0.10,\yA) {$A$};
					\node[font=\scriptsize, anchor=east] at (\xast-0.10,\yB) {$B$};
					\node[font=\scriptsize, anchor=east] at (\xast-0.10,\yC) {$C$};
					
					%% --- Etichette dei corpi (a sinistra, mid-asta) ---
					\node[font=\scriptsize, anchor=east]
					at (\xast-0.10, {(\yA+\yB)/2}) {$\mathcal{C}_1$};
					\node[font=\scriptsize, anchor=east]
					at (\xast-0.10, {(\yB+\yC)/2}) {$\mathcal{C}_2$};
					
					
					%% --- Deformata in grigio: B -- C' ---
					%% C' è l'estremo della freccia q_2 in C
					\pgfmathsetmacro{\xCp}{\xast+\Lq}
					\pgfmathsetmacro{\yCp}{\yC}
					
					% Segmento B -- C' (deformata di C_2)
					\draw[gray!75, line width=1.0pt] (\xast,\yB) -- (\xCp,\yCp);
					
					% Anellino grigio in C'
					\draw[gray!75, fill=white, line width=0.6pt] (\xCp,\yCp) circle (1.6pt);
					
					% Etichetta C'
					\node[gray!75, font=\scriptsize, anchor=west] at (\xCp+0.05,\yCp) {$C'$};
					
					
					%% --- Etichetta del pannello ---
					\node[font=\scriptsize] at (0.8,-1.4) {(g)};
					
				\end{scope}
				% pannello (h)
				\begin{scope}[xshift=\dx cm]
					
					%% --- Parametri ---
					\def\ell{1.2}
					\def\xast{-0.7}
					\def\Lb{0.55}        % lunghezza delle bielle (corte)
					\def\angBip{30}      % inclinazione delle bielle interne/esterne
					\def\dAtt{0.20}      % distanza tra i due anellini sull'asta
					\def\dsuolo{0.18}
					\def\hsuolo{0.30}
					\def\hatch{0.13}
					\def\rn{0.06}        % raggio degli anellini semicircolari
					\def\rArc{0.13}      % raggio degli archi di rotazione q_i
					
					%% --- Coordinate di riferimento ---
					\pgfmathsetmacro{\yA}{-\ell}
					\pgfmathsetmacro{\yB}{0}
					\pgfmathsetmacro{\yC}{\ell}
					\pgfmathsetmacro{\ymidUno}{(\yA+\yB)/2}
					\pgfmathsetmacro{\ymidDue}{(\yB+\yC)/2}
					
					%% --- Costanti geometriche del pantografo ---
					\pgfmathsetmacro{\dxBie}{\Lb*cos(\angBip)}
					\pgfmathsetmacro{\dyBie}{\Lb*sin(\angBip)}
					\pgfmathsetmacro{\xMuro}{\xast+2*\dxBie}
					
					%				%% --- Aste verticali (in nero) ---
					%				\draw[line width=1.2pt] (\xast,\yA) -- (\xast,\yB);
					%				\draw[line width=1.2pt] (\xast,\yB) -- (\xast,\yC);
					
					%% =========================================================
					%% Bipendolo su C_1
					%% =========================================================
					\pgfmathsetmacro{\yAiU}{\ymidUno+\dAtt/2}
					\pgfmathsetmacro{\yAiD}{\ymidUno-\dAtt/2}
					\pgfmathsetmacro{\yBiU}{\yAiU-\dyBie}
					\pgfmathsetmacro{\yBiD}{\yAiD-\dyBie}
					
					% bielle interne (parallelogramma 1)
					\draw[gray!75, line width=0.8pt] (\xast,\yAiU) -- ({\xast+\dxBie},\yBiU);
					\draw[gray!75, line width=0.8pt] (\xast,\yAiD) -- ({\xast+\dxBie},\yBiD);
					% traversa
					\draw[gray!75, line width=0.8pt] ({\xast+\dxBie},\yBiU) -- ({\xast+\dxBie},\yBiD);
					% bielle esterne (parallelogramma 2)
					\draw[gray!75, line width=0.8pt] ({\xast+\dxBie},\yBiU) -- (\xMuro,\yAiU);
					\draw[gray!75, line width=0.8pt] ({\xast+\dxBie},\yBiD) -- (\xMuro,\yAiD);
					% cerniere intermedie
					\draw[gray!75, fill=white, line width=0.6pt] ({\xast+\dxBie},\yBiU) circle (1.5pt);
					\draw[gray!75, fill=white, line width=0.6pt] ({\xast+\dxBie},\yBiD) circle (1.5pt);
					% anellini sul corpo C_1 (semicerchi sporgenti a destra)
					\draw[gray!75, fill=white, line width=0.6pt]
					(\xast,\yAiU-\rn) arc(-90:90:\rn) -- cycle;
					\draw[gray!75, fill=white, line width=0.6pt]
					(\xast,\yAiD-\rn) arc(-90:90:\rn) -- cycle;
					% anellini sul muro (semicerchi sporgenti a sinistra)
					\draw[gray!75, fill=white, line width=0.6pt]
					(\xMuro,\yAiU+\rn) arc(90:270:\rn) -- cycle;
					\draw[gray!75, fill=white, line width=0.6pt]
					(\xMuro,\yAiD+\rn) arc(90:270:\rn) -- cycle;
					% muro verticale con hatching
					\pgfmathsetmacro{\yMurMin}{\yAiD-0.12}
					\pgfmathsetmacro{\yMurMax}{\yAiU+0.12}
					\draw[gray!75, thick] (\xMuro,\yMurMin) -- (\xMuro,\yMurMax);
					\foreach \k in {0,1,2,3,4,5}{
						\pgfmathsetmacro{\yh}{\yMurMin+\k*(\yMurMax-\yMurMin)/5}
						\draw[gray!75, thin] (\xMuro,\yh) -- ({\xMuro+\hatch},\yh-0.09);
					}
					
					
					%% --- Aste verticali (in nero) ---
					\draw[line width=1.2pt] (\xast,\yA) -- (\xast,\yB);
					\draw[line width=1.2pt] (\xast,\yB) -- (\xast,\yC);
					
					
					%% --- Cerniera al suolo in A (nero) ---
					\draw[thick]
					(\xast,\yA) -- (\xast-\dsuolo,\yA-\hsuolo)
					-- (\xast+\dsuolo,\yA-\hsuolo) -- cycle;
					\draw[thick] (\xast-0.28,\yA-\hsuolo) -- (\xast+0.28,\yA-\hsuolo);
					\foreach \x in {-0.20,-0.10,0,0.10,0.20}{
						\draw[thin] (\xast+\x,\yA-\hsuolo)
						-- (\xast+\x-0.09,\yA-\hsuolo-\hatch);
					}
					
					%% --- Anellini in A e B (nero) ---
					\node[circle, fill=white, draw=black, inner sep=0pt,
					minimum size=4pt, line width=0.7pt] at (\xast,\yA) {};
					\node[circle, fill=white, draw=black, inner sep=0pt,
					minimum size=4pt, line width=0.7pt] at (\xast,\yB) {};
					
					%				%% --- Frecce di rotazione q_1 e q_2 (nero), centrate sull'asta ---
					
					\def\qDeg{20}       % apertura grafica della rotazione
					\def\Lref{0.90}     % lunghezza dei due segmenti
					\def\rArcQ{0.70}    % raggio dell'archetto (minore di \Lref)
					\pgfmathsetmacro{\xRot}{\xMuro+0.55}   % a destra del muro
					
					% --- q_2: a fianco del bipendolo su C_2 ---
					\draw[thin] (\xRot,\ymidDue) -- (\xRot,{\ymidDue+\Lref});
					\draw[thin] (\xRot,\ymidDue) -- 
					({\xRot-\Lref*sin(\qDeg)},{\ymidDue+\Lref*cos(\qDeg)});
					\draw[->, thin] 
					(\xRot,{\ymidDue+\rArcQ}) 
					arc[start angle=90, end angle={90+\qDeg}, radius=\rArcQ];
					\node[font=\scriptsize] at ({\xRot+(\rArcQ+0.12)*cos(90+2*\qDeg)+0.1},{\ymidDue+(\rArcQ+0.12)*sin(90+\qDeg/2)-0.2}) {$q_2$};
					
					%% --- Etichette ---
					\node[font=\scriptsize, anchor=east] at (\xast-0.10,\yA) {$A$};
					\node[font=\scriptsize, anchor=east] at (\xast-0.10,\yB) {$B$};
					%\node[font=\scriptsize, anchor=east] at (\xast-0.10,\yC) {$C$};
					\node[font=\scriptsize, anchor=south] at (\xast+0.06,\yC) {$C$};
					\node[font=\scriptsize, anchor=east]
					at (\xast-0.10, \ymidUno) {$\mathcal{C}_1$};
					\node[font=\scriptsize, anchor=east] at (\xast-0.10, \ymidDue) {$\mathcal{C}_2$};
					
					%% --- Deformata in grigio: B -- C' (rotazione antioraria di C_2 attorno a B) ---
					%% C' è a sinistra di C di una quantità grafica pari a \Lq (= ell*q_2)
					\def\Lq{0.55}
					\pgfmathsetmacro{\xCp}{\xast-\Lq}
					\pgfmathsetmacro{\yCp}{\yC}
					
					% Segmento B -- C' (deformata di C_2)
					\draw[gray!75, line width=1.0pt] (\xast,\yB) -- (\xCp,\yCp);
					
					% Anellino grigio in C'
					\draw[gray!75, fill=white, line width=0.6pt] (\xCp,\yCp) circle (1.6pt);
					
					% Etichetta C' (a sinistra, per non sovrapporsi con C)
					\node[gray!75, font=\scriptsize, anchor=east] at (\xCp-0.05,\yCp) {$C'$};
					
					%% --- Archetto della rotazione q_2 (antioraria, attorno a B) ---
					\def\qDegII{22}        % apertura grafica della rotazione
					\def\rArcII{0.85}      % raggio dell'archetto (< ell)
					
					\draw[->, gray!75, thin]
					(\xast,{\yB+\rArcII})
					arc[start angle=90, end angle={90+\qDegII}, radius=\rArcII];
					
					% Etichetta q_2 al punto medio dell'arco, sul lato esterno
					\node[gray!75, font=\scriptsize, anchor=east]	at ({0.6*\xast+\rArcII*cos(90+\qDegII/2)}, {\yB+\rArcII*sin(90+\qDegII/2)+0.2})
					{$q_2$};
					
					%% --- Etichetta del pannello ---
					\node[font=\scriptsize] at (0.8,-1.4) {(h)};
				\end{scope}
				% pannello (i)
				\begin{scope}[xshift=2*\dx cm]
					
					%% --- Parametri comuni ---
					\def\ell{1.2}
					\def\xast{-0.7}
					\def\dsuolo{0.18}
					\def\hsuolo{0.30}
					\def\hatch{0.13}
					\def\rn{0.06}
					
					%% --- Coordinate dei punti notevoli ---
					\pgfmathsetmacro{\yA}{-\ell}
					\pgfmathsetmacro{\yB}{0}
					\pgfmathsetmacro{\yC}{\ell}
					
					%% =========================================================
					%% Bipendolo su C_1 al suolo (traslato verso il basso di Dymid)
					%% =========================================================
					\def\Lb{0.55}
					\def\angBip{30}
					\def\dAtt{0.20}
					\def\Dymid{-0.20}    % traslazione verticale rispetto al pannello (b)
					\pgfmathsetmacro{\ymidUno}{(\yA+\yB)/2 + \Dymid}
					\pgfmathsetmacro{\dxBie}{\Lb*cos(\angBip)}
					\pgfmathsetmacro{\dyBie}{\Lb*sin(\angBip)}
					\pgfmathsetmacro{\xMuro}{\xast+2*\dxBie}
					
					\pgfmathsetmacro{\yAiU}{\ymidUno+\dAtt/2}
					\pgfmathsetmacro{\yAiD}{\ymidUno-\dAtt/2}
					\pgfmathsetmacro{\yBiU}{\yAiU-\dyBie}
					\pgfmathsetmacro{\yBiD}{\yAiD-\dyBie}
					
					\draw[gray!75, line width=0.8pt] (\xast,\yAiU) -- ({\xast+\dxBie},\yBiU);
					\draw[gray!75, line width=0.8pt] (\xast,\yAiD) -- ({\xast+\dxBie},\yBiD);
					\draw[gray!75, line width=0.8pt] ({\xast+\dxBie},\yBiU) -- ({\xast+\dxBie},\yBiD);
					\draw[gray!75, line width=0.8pt] ({\xast+\dxBie},\yBiU) -- (\xMuro,\yAiU);
					\draw[gray!75, line width=0.8pt] ({\xast+\dxBie},\yBiD) -- (\xMuro,\yAiD);
					\draw[gray!75, fill=white, line width=0.6pt] ({\xast+\dxBie},\yBiU) circle (1.5pt);
					\draw[gray!75, fill=white, line width=0.6pt] ({\xast+\dxBie},\yBiD) circle (1.5pt);
					\draw[gray!75, fill=white, line width=0.6pt]
					(\xast,\yAiU-\rn) arc(-90:90:\rn) -- cycle;
					\draw[gray!75, fill=white, line width=0.6pt]
					(\xast,\yAiD-\rn) arc(-90:90:\rn) -- cycle;
					\draw[gray!75, fill=white, line width=0.6pt]
					(\xMuro,\yAiU+\rn) arc(90:270:\rn) -- cycle;
					\draw[gray!75, fill=white, line width=0.6pt]
					(\xMuro,\yAiD+\rn) arc(90:270:\rn) -- cycle;
					\pgfmathsetmacro{\yMurMin}{\yAiD-0.12}
					\pgfmathsetmacro{\yMurMax}{\yAiU+0.12}
					\draw[gray!75, thick] (\xMuro,\yMurMin) -- (\xMuro,\yMurMax);
					\foreach \k in {0,1,2,3,4,5}{
						\pgfmathsetmacro{\yh}{\yMurMin+\k*(\yMurMax-\yMurMin)/5}
						\draw[gray!75, thin] (\xMuro,\yh) -- ({\xMuro+\hatch},\yh-0.09);
					}
					
					%% =========================================================
					%% Aste verticali, cerniere, anellini reali (in nero, sopra)
					%% =========================================================
					\draw[line width=1.2pt] (\xast,\yA) -- (\xast,\yB);
					\draw[line width=1.2pt] (\xast,\yB) -- (\xast,\yC);
					
					\draw[thick]
					(\xast,\yA) -- (\xast-\dsuolo,\yA-\hsuolo)
					-- (\xast+\dsuolo,\yA-\hsuolo) -- cycle;
					\draw[thick] (\xast-0.28,\yA-\hsuolo) -- (\xast+0.28,\yA-\hsuolo);
					\foreach \x in {-0.20,-0.10,0,0.10,0.20}{
						\draw[thin] (\xast+\x,\yA-\hsuolo)
						-- (\xast+\x-0.09,\yA-\hsuolo-\hatch);
					}
					\node[circle, fill=white, draw=black, inner sep=0pt,
					minimum size=4pt, line width=0.7pt] at (\xast,\yA) {};
					\node[circle, fill=white, draw=black, inner sep=0pt,
					minimum size=4pt, line width=0.7pt] at (\xast,\yB) {};
					
					
					%% =========================================================
					%% Indicazione di q_2 (rotazione relativa C_2/C_1)
					%% =========================================================
					\def\xq{0.40}              % ascissa comune dei due semicerchi
					\def\dq{0.07}              % semi-distanza verticale fra le due frecce
					\def\rq{0.18}              % raggio dei semicerchi
					\pgfmathsetmacro{\yqup}{\yB+\dq}
					\pgfmathsetmacro{\yqdn}{\yB-\dq}
					
					% Semicerchio superiore: antiorario, convessità verso l'alto
					% (da destra, sopra il vertice, fino a sinistra — la freccia punta in basso)
					\draw[->, thin]
					({\xq+\rq},\yqup)
					arc[start angle=0, end angle=180, radius=\rq];
					
					% Semicerchio inferiore: orario, convessità verso il basso
					% (da destra, sotto il vertice, fino a sinistra — la freccia punta in alto)
					\draw[->, thin]
					({\xq+\rq},\yqdn)
					arc[start angle=0, end angle=-180, radius=\rq];
					
					% Etichetta q_2 a destra delle due frecce
					\node[font=\scriptsize] at ({\xq+\rq+0.20},\yB) {$q_2$};
					
					%% =========================================================
					%% Etichette
					%% =========================================================
					\node[font=\scriptsize, anchor=east] at (\xast-0.10,\yA) {$A$};
					\node[font=\scriptsize, anchor=east] at (\xast-0.10,\yB) {$B$};
					\node[font=\scriptsize, anchor=south] at (\xast+0.06,\yC) {$C$};
					\node[font=\scriptsize, anchor=east]
					at (\xast-0.10, {(\yA+\yB)/2}) {$\mathcal{C}_1$};
					\node[font=\scriptsize, anchor=east]
					at (\xast-0.10, {(\yB+\yC)/2}) {$\mathcal{C}_2$};
					
					%% --- Deformata in grigio: B -- C' (rotazione antioraria di C_2 attorno a B) ---
					%% C' è a sinistra di C di una quantità grafica pari a \Lq (= ell*q_2)
					\def\Lq{0.55}
					\pgfmathsetmacro{\xCp}{\xast-\Lq}
					\pgfmathsetmacro{\yCp}{\yC}
					
					% Segmento B -- C' (deformata di C_2)
					\draw[gray!75, line width=1.0pt] (\xast,\yB) -- (\xCp,\yCp);
					
					% Anellino grigio in C'
					\draw[gray!75, fill=white, line width=0.6pt] (\xCp,\yCp) circle (1.6pt);
					
					% Etichetta C' (a sinistra, per non sovrapporsi con C)
					\node[gray!75, font=\scriptsize, anchor=east] at (\xCp-0.05,\yCp) {$C'$};
					
					%% --- Archetto della rotazione q_2 (antioraria, attorno a B) ---
					\def\qDegII{22}        % apertura grafica della rotazione
					\def\rArcII{0.85}      % raggio dell'archetto (< ell)
					
					\draw[->, gray!75, thin]
					(\xast,{\yB+\rArcII})
					arc[start angle=90, end angle={90+\qDegII}, radius=\rArcII];
					
					% Etichetta q_2 al punto medio dell'arco, sul lato esterno
					\node[gray!75, font=\scriptsize, anchor=east]	at ({0.6*\xast+\rArcII*cos(90+\qDegII/2)}, {\yB+\rArcII*sin(90+\qDegII/2)+0.2})
					{$q_2$};
					
					%% --- Etichetta del pannello ---
					\node[font=\small] at (0.8,-1.4) {(i)};
					
				\end{scope}
			\end{scope}
			
		\end{tikzpicture}
		\caption{Due aste verticali collegate in serie
			(\ESEref{ex:scelte_fittizi}): tre scelte di
			vincoli fittizi e i corrispondenti meccanismi.
			\emph{Riga 1} --- sistemi principali:
			\emph{(a)}~scelta~1, pendolo orizzontale in $B$
			e pendolo orizzontale in $C$ (in grigio);
			\emph{(b)}~scelta~2, bipendolo su $\mathcal{C}_1$
			e bipendolo su $\mathcal{C}_2$ al suolo (in grigio);
			\emph{(c)}~scelta~3, bipendolo su $\mathcal{C}_1$
			al suolo e bipendolo interno tra $\mathcal{C}_2$
			e $\mathcal{C}_1$ (in grigio).
			\emph{Riga 2} --- meccanismi $q_1 \neq 0$:
			\emph{(d)}~$\mathcal{C}_1$ ruota attorno ad $A$,
			$\mathcal{C}_2$ ruota attorno a $C$;
			\emph{(e)}~$\mathcal{C}_1$ ruota attorno ad $A$,
			$\mathcal{C}_2$ trasla orizzontalmente di
			$q_1\ell$;
			\emph{(f)}~le due aste ruotano rigidamente
			attorno ad $A$.
			\emph{Riga 3} --- meccanismi $q_2 \neq 0$:
			\emph{(g)}, \emph{(h)}, \emph{(i)}~in tutti
			e tre i casi $\mathcal{C}_1$ è ferma e
			$\mathcal{C}_2$ ruota attorno a $B$.}
		\label{fig:scelte_fittizi}
	\end{figure}
	\begin{svolgimento}
		
		\esottotit{Scelta 1: pendolo orizzontale in $B$
			e pendolo orizzontale in $C$.}
		Le coordinate lagrangiane sono gli spostamenti
		orizzontali assoluti $q_1 = u(B)$ e $q_2 = u(C)$.
		
		Per $q_1 \neq 0$, $q_2 = 0$: il pendolo in $C$ è
		bloccato, quindi $C$ non si sposta orizzontalmente.
		Il vincolo in $A$ e il pendolo in $C$ fissano il
		\CRa\ di $\mathcal{C}_2$ in $C$; il pendolo in $B$
		cede di $q_1$, quindi $\mathcal{C}_1$ ruota attorno
		ad $A$ di $\vartheta_1 = q_1/\ell$ e $\mathcal{C}_2$
		ruota attorno a $C$ della stessa ampiezza.
		
		Per $q_2 \neq 0$, $q_1 = 0$: il pendolo in $B$ è
		bloccato, quindi $\mathcal{C}_1$ è ferma. Il \CRa\
		di $\mathcal{C}_2$ è $B$; il pendolo in $C$ cede di
		$q_2$, quindi $\mathcal{C}_2$ ruota attorno a $B$ di
		$\vartheta_2 = q_2/\ell$.
		
		\esottotit{Scelta 2: bipendolo su $\mathcal{C}_1$
			al suolo e bipendolo su $\mathcal{C}_2$ al suolo.}
		Le coordinate lagrangiane sono le rotazioni assolute
		$q_1 = \vartheta_1$ e $q_2 = \vartheta_2$.
		
		Per $q_1 \neq 0$, $q_2 = 0$: il bipendolo su
		$\mathcal{C}_1$ assegna a quest'ultima la rotazione
		$\vartheta_1 = q_1$; la cerniera in $A$ fissa il
		\CRa\ di $\mathcal{C}_1$ in $A$, quindi $B$ si
		sposta orizzontalmente di $q_1\ell$. Il bipendolo
		su $\mathcal{C}_2$ impone $\vartheta_2 = 0$: poiché
		$B$ si è spostato di $q_1\ell$ e $\mathcal{C}_2$
		non ruota, essa trasla orizzontalmente di $q_1\ell$.
		
		Per $q_2 \neq 0$, $q_1 = 0$: il bipendolo su
		$\mathcal{C}_1$ impone $\vartheta_1 = 0$, quindi
		$\mathcal{C}_1$ è ferma e $B$ non si sposta. Il
		bipendolo su $\mathcal{C}_2$ assegna a quest'ultima
		la rotazione $\vartheta_2 = q_2$; poiché $B$ è
		fisso, il \CRa\ di $\mathcal{C}_2$ coincide con
		$B$: $\mathcal{C}_2$ ruota attorno a $B$ di $q_2$.
		
		\esottotit{Scelta 3: bipendolo su $\mathcal{C}_1$
			al suolo e bipendolo tra $\mathcal{C}_2$ e
			$\mathcal{C}_1$.}
		Le coordinate lagrangiane sono la rotazione assoluta
		$q_1 = \vartheta_1$ di $\mathcal{C}_1$ e la rotazione
		relativa $q_2 = \vartheta_{21}$ di $\mathcal{C}_2$
		rispetto a $\mathcal{C}_1$.
		
		Per $q_1 \neq 0$, $q_2 = 0$: la rotazione relativa
		è nulla, quindi le due aste ruotano rigidamente come
		un unico corpo attorno ad $A$ di $q_1$.
		
		Per $q_2 \neq 0$, $q_1 = 0$: $\mathcal{C}_1$ è
		ferma; $\mathcal{C}_2$ ruota attorno a $B$ di $q_2$.
		
		\vspace{6pt}
		
		Nonostante le tre scelte producano meccanismi di
		forma diversa, l'insieme delle configurazioni
		ammissibili è lo stesso: qualsiasi configurazione
		descritta dalla scelta 1 può essere espressa come
		combinazione lineare dei meccanismi della scelta 2
		o della scelta 3, e viceversa. Questo riflette il
		fatto che i vincoli fittizi definiscono una base
		per $\Ker(\bm{A})$: basi diverse descrivono lo
		stesso sottospazio, e dunque lo stesso insieme di
		configurazioni. La labilità --- il numero di gradi
		di libertà residui --- è una proprietà intrinseca
		del sistema, indipendente da qualsiasi scelta di
		rappresentazione.		
	\end{svolgimento}
\end{esempio}



\subsection{Problema statico: condizione di compatibilità}
\label{sec:labile_stat}

Per un sistema labile, il problema statico $\bm{A}^{\top}\bm{r} =
-\bm{f}$ è \emph{sovradeterminato}: ha $n = 3n_c$ equazioni e $m < n$
incognite. La soluzione esiste se e solo se
$\bm{f} \in \Imm(\bm{A}^{\top})$,
ovvero (\CAPref{cap:problemi_cin_stat_SCR}):
\begin{equation}\label{eq:compat_stat}
	\bm{U}_q^{\top}\,\bm{f} = \bm{0},
	\qquad\text{ovvero}\qquad
	\bm{u}_k^{\top}\,\bm{f} = 0,
	\quad k = 1, 2, \ldots, g_\ell,
\end{equation}
ossia, per il principio dei lavori virtuali (\PARref{sec:PLV}),
le forze attive devono compiere lavoro virtuale nullo su ogni
meccanismo. Questa è la \emph{condizione di compatibilità statica}\index{compatibilita@compatibilità!statica}:
un sistema labile è in equilibrio solo se il carico è
\emph{ortogonale} allo spazio dei meccanismi.

\paragraph{Interpretazione fisica.}
La condizione~\eqref{eq:compat_stat} è la scrittura matriciale
dell'equilibrio alla rotazione del meccanismo: per ciascuna
colonna $\bm{u}_k$ di $\bm{U}_q$, l'equazione
$\bm{u}_k^{\top}\bm{f} = 0$ esprime che le forze attive non
producono momento risultante rispetto al centro di rotazione
del meccanismo.

Quando la condizione è soddisfatta, la soluzione del problema
statico è \emph{unica} (poiché $g_{\mathscr{i}} = 0$ nei sistemi
labili non degeneri) e si determina con le stesse strategie della
\PARref{sec:determinato}, applicando le equazioni di equilibrio
ai corpi liberi.

\begin{esempio}[Sistema labile di grado 1 ---
	problema statico]
	\label{ex:labile_L_stat}
	Si consideri lo stesso corpo a L
	dell'\ESEref{ex:labile_L}. Sono applicate in $D$
	una forza orizzontale $F_x$ verso destra e una
	forza verticale $F_y$ verso il basso
	(\FIGref{fig:labile_L_stat}). Verificare la
	condizione di compatibilità statica e determinare
	le reazioni vincolari quando essa è soddisfatta.
	
	\begin{figure}[ht]
		\centering
		\begin{tikzpicture}[>=Stealth, line width=0.8pt]
			
			\def\L{2.0}
			\def\F{0.5}  % lunghezza frecce forze
			
			% ============================================================
			% PANNELLO SINISTRO — dati
			% ============================================================
			\begin{scope}[xshift=0cm]
				
				% corpo a L
				
				\draw[line width=1.5pt] (0,0) -- (\L,0) -- (\L,\L);
				
				
				
				% rette tratteggiate per C
				\draw[gray!50, dashed, thin]
				(0,-0.3) -- (0,\L+0.2);
				\draw[gray!50, dashed, thin]
				(-0.3,\L) -- (\L+0.2,\L);
				
				% carrello verticale in A
				\draw[thick] (0,0) -- (-0.28,-0.48)
				-- (0.28,-0.48) -- cycle;
				\draw[thick] (-0.18,-0.55) circle (0.08);
				\draw[thick] ( 0.18,-0.55) circle (0.08);
				\draw[thick] (-0.38,-0.65) -- (0.38,-0.65);
				\foreach \x in {-0.30,-0.15,0,0.15,0.30}{
					\draw[thin] (\x,-0.65) -- (\x-0.12,-0.83);
				}
				
				% carrello orizzontale in B
				\draw[thick] (\L,\L) -- (\L+0.48,\L+0.28)
				-- (\L+0.48,\L-0.28) -- cycle;
				\draw[thick] (\L+0.55,\L+0.18) circle (0.08);
				\draw[thick] (\L+0.55,\L-0.18) circle (0.08);
				\draw[thick] (\L+0.65,\L+0.38) -- (\L+0.65,\L-0.38);
				\foreach \y in {-0.30,-0.15,0,0.15,0.30}{
					\draw[thin] (\L+0.65,\L+\y)
					-- (\L+0.83,\L+\y-0.12);
				}
				
				% punti
				\draw[fill=white, thick] (0,0) circle (3pt)
				node[above left, font=\scriptsize] {$A$};
				\draw[fill=white, thick] (\L,\L) circle (3pt)
				node[above left, font=\scriptsize] {$B$};
				\fill (\L,0) %circle (2pt)
				node[above left, font=\scriptsize] {$D$};
				\node[left, font=\scriptsize] at (0,\L) {$C$};
				
				% forza F1 orizzontale in D
				\draw[->, thick] (\L,0) -- ({\L+\F},0)
				node[right, font=\scriptsize] {$F_x$};
				
				% forza F2 verticale in D
				\draw[->, thick] (\L,0) -- (\L,-\F)
				node[right, font=\scriptsize] {$F_y$};
				
				% etichetta pannello
				\node[font=\scriptsize] at (-\L/2+0.2,-0.8) {(a)};
				
			\end{scope}
			
			% ============================================================
			% PANNELLO DESTRO — interpretazione compatibilità
			% ============================================================
			\begin{scope}[xshift=5.5cm]
				
				% corpo a L
				\draw[line width=1.5pt] (0,0) -- (\L,0) -- (\L,\L);
				
				% rette tratteggiate per C
				\draw[gray!50, dashed, thin]
				(0,-0.3) -- (0,\L+0.2);
				\draw[gray!50, dashed, thin]
				(-0.3,\L) -- (\L+0.2,\L);
				
				% carrello verticale in A
				\draw[thick] (0,0) -- (-0.28,-0.48)
				-- (0.28,-0.48) -- cycle;
				\draw[thick] (-0.18,-0.55) circle (0.08);
				\draw[thick] ( 0.18,-0.55) circle (0.08);
				\draw[thick] (-0.38,-0.65) -- (0.38,-0.65);
				\foreach \x in {-0.30,-0.15,0,0.15,0.30}{
					\draw[thin] (\x,-0.65) -- (\x-0.12,-0.83);
				}
				
				% carrello orizzontale in B
				\draw[thick] (\L,\L) -- (\L+0.48,\L+0.28)
				-- (\L+0.48,\L-0.28) -- cycle;
				\draw[thick] (\L+0.55,\L+0.18) circle (0.08);
				\draw[thick] (\L+0.55,\L-0.18) circle (0.08);
				\draw[thick] (\L+0.65,\L+0.38) -- (\L+0.65,\L-0.38);
				\foreach \y in {-0.30,-0.15,0,0.15,0.30}{
					\draw[thin] (\L+0.65,\L+\y)
					-- (\L+0.83,\L+\y-0.12);
				}
				
				% punti
				\draw[fill=white, thick] (0,0) circle (3pt)
				node[above left, font=\scriptsize] {$A$};
				\draw[fill=white, thick] (\L,\L) circle (3pt)
				node[above left, font=\scriptsize] {$B$};
				\fill (\L,0) %circle (2pt)
				node[above left, font=\scriptsize] {$D$};
				\fill (0,\L) circle (2pt)
				node[left, font=\scriptsize] {$C$};
				
				% forza F orizzontale in D
				\draw[->, thick] (\L,0) -- ({\L+\F},0)
				node[right, font=\scriptsize] {$F$};
				
				% forza F verticale in D
				\draw[->, thick] (\L,0) -- (\L,-\F)
				node[right, font=\scriptsize] {$F$};
				
				% bracci rispetto a C
				% braccio di F1 rispetto a C: distanza verticale = ell
				\draw[<->, gray!60, thin] ({\L+\F+0.6},0) -- ({\L+\F+0.6},\L) node[midway, left, font=\scriptsize] {$\ell$};
				
				% braccio di F2 rispetto a C: distanza orizzontale = ell
				\draw[<->, gray!60, thin] (0,-\F-0.2) -- (\L,-\F-0.2) node[midway, above, font=\scriptsize] {$\ell$};
				
				% momento risultante nullo in C
				\node[font=\scriptsize, align=center] at (0.6, 0.5)
				{$\mu_C = 0$};
				
				% etichetta pannello
				\node[font=\scriptsize] at (-\L/2+0.2,-0.8) {(b)};
				
			\end{scope}
			
		\end{tikzpicture}
		\caption{Corpo a L --- problema statico
			(\ESEref{ex:labile_L_stat}):
			dati \emph{(a)} e interpretazione della
			condizione di compatibilità \emph{(b)}.
			La condizione $F_x = F_y = F$ equivale
			al momento risultante nullo rispetto a $C$:
			$\mu_C = F\cdot\ell - F\cdot\ell = 0$.}
		\label{fig:labile_L_stat}
	\end{figure}
	
	\begin{svolgimento}
		
		\esottotit{Condizione di compatibilità
			statica.}
		Il meccanismo del sistema labile è
		$\bm{u}_1 = \{1,\, 0,\, 1/\ell\}^\top$
		(\ESEref{ex:labile_L}). Il vettore delle forze
		generalizzate riferito al polo $O \equiv A$ è:
		\begin{equation*}
			\bm{f} =
			\bigl\{F_x\; -F_y\; -F_y\ell
			\bigr\}^\top,
		\end{equation*}
		dove $-F_y\ell$ è il momento di $F_y$ rispetto
		ad $A$ (braccio $\ell$, verso orario). La
		condizione di compatibilità statica è:
		\begin{equation*}
			\bm{u}_1^\top\bm{f} = 0:\quad
			F_x + 0 \cdot(-F_y) +
			\frac{1}{\ell}\cdot(-F_y\ell) =
			F_x - F_y = 0
			\;\Longrightarrow\;
			F_x = F_y.
		\end{equation*}
		La condizione equivale a richiedere momento
		risultante nullo rispetto a $C = (0,\ell)$,
		centro di rotazione del meccanismo:
		\begin{equation*}
			\mu_C = F_x\cdot\ell - F_y\cdot\ell =
			(F_x - F_y)\,\ell = 0.
		\end{equation*}
		
		\esottotit{Reazioni vincolari.}
		Quando $F_x = F_y = F$, il sistema è in
		equilibrio. Le reazioni incognite sono $Y_A$
		(verticale in $A$) e $X_B$ (orizzontale in $B$).
		Dall'equilibrio globale:
		\begin{equation*}
			\sum X = 0:\quad X_B + F = 0
			\;\Longrightarrow\; X_B = -F,
		\end{equation*}
		\begin{equation*}
			\sum Y = 0:\quad Y_A - F = 0
			\;\Longrightarrow\; Y_A = F.
		\end{equation*}
		La condizione $F_x = F_y$ non riguarda le
		intensità assolute dei due carichi, ma il
		loro rapporto: qualsiasi coppia di forze
		uguali in $D$ è compatibile con l'equilibrio
		del sistema labile, indipendentemente dalla
		loro intensità comune $F$.
		
		\esottotit{Equilibrio alla rotazione.}
		L'equazione di momento rispetto ad $A$ è automaticamente
		soddisfatta:
		\begin{equation*}
			\sum \mu(A) = -F\,\ell - \ell\,X_B
			= -F\,\ell - \ell(-F) = 0.
		\end{equation*}
		Questa non è una circostanza fortuita. In un sistema labile
		di grado~$g_\ell = 1$ la condizione di compatibilità
		statica~\eqref{eq:compat_stat}, $\bm{u}_1^{\!\top}\bm{f} = 0$,
		è precisamente la relazione di dipendenza lineare tra le tre
		equazioni di equilibrio: una volta imposta, una delle tre
		equazioni diventa combinazione lineare delle altre due. Le
		equazioni scalari indipendenti sono dunque
		$m = n - g_\ell = 3 - 1 = 2$, in accordo con il rango
		$\rg(\bm{A}^{\!\top}) = m$.
	\end{svolgimento}
\end{esempio}


%----------------------------------------------------------------------------------------------------------------------------
\section{Sistemi iperstatici}\label{sec:iperstatico}
%----------------------------------------------------------------------------------------------------------------------------

Un sistema è \emph{iperstatico}\index{iperstaticita@iperstaticità} di grado $g_{\mathscr{i}}$ quando
$\dim\Ker(\Abb^{\top}) = g_{\mathscr{i}} > 0$: esistono
$g_{\mathscr{i}}$ combinazioni di reazioni vincolari linearmente
indipendenti compatibili con l'equilibrio in assenza di forze
attive, dette \emph{stati di autoequilibrio}. La condizione
necessaria è $\nu = 3n_c - m < 0$ (vincoli in numero eccessivo);
tuttavia l'iperstaticità può presentarsi anche con $\nu = 0$,
in concorrenza con una labilità di pari grado per effetto della
degenerazione geometrica: è il caso labile-iperstatico, le cui
condizioni operative sono trattate in \PARref{sec:labile_iperstatico}.

I due problemi si comportano in modo duale rispetto al caso labile
(\PARref{sec:labile}): il problema statico è ora
\emph{sottodeterminato} (infinite soluzioni), mentre quello
cinematico è \emph{sovradeterminato} (soluzione possibile solo per
certi cedimenti).


\subsection{Problema statico: vincoli ridondanti e iperstatiche}
\label{sec:iperstatico_stat}

La soluzione generale del problema statico
\eqref{eq:due_problemi}, $\bm{A}^{\top}\bm{r} = -\bm{f}$, è:
\begin{equation}\label{eq:sol_iper_stat}
	\bm{r} = \bm{r}_0 + \bm{R}_\chi\,\bm{\chi},
	\qquad \bm{\chi} \in \mathbb{R}^{g_{\mathscr{i}}},
\end{equation}
dove $\bm{r}_0$ è una soluzione particolare e le colonne di
$\bm{R}_\chi$ formano una base di $\Ker(\Abb^{\top})$. Le
componenti di $\bm{\chi}$ sono le \emph{iperstatiche}, parametri
liberi che descrivono la parte di reazione non determinata
dall'equilibrio.

La struttura della soluzione \eqref{eq:sol_iper_stat} è il duale esatto della \eqref{eq:sol_labile_cin}: il metodo operativo procede in modo speculare a quello del \PARref{sec:labile_cin}. La matrice di equilibrio
$\bm{A}^{\top}$ è rettangolare bassa ($n < m$), esattamente
come la matrice dei vincoli $\bm{A}$ è rettangolare bassa
($m < n$) nei sistemi labili: in entrambi i casi il sistema
principale si ottiene aggiungendo le $g_{\mathscr{i}}$ equazioni
mancanti. Nel problema cinematico labile le equazioni aggiunte
sono i $g_\ell$ vincoli fittizi, con i cedimenti lagrangiani
$q_k$ come parametri liberi al secondo membro; nel problema
statico iperstatico le equazioni aggiunte sono le
$g_{\mathscr{i}}$ condizioni di equilibrio che emergono
operando la \emph{sconnessione} dei vincoli ridondanti: al
punto di sconnessione la reazione $\chi_k$ del vincolo
soppresso compare come parametro libero al secondo membro
dell'equilibrio dei corpi adiacenti. In entrambi i casi il
sistema principale risultante è isostatico ($n \times n$), e
i parametri liberi sono i termini noti delle equazioni
aggiunte.


\paragraph{Procedura.}
\begin{enumerate}[label=\textbf{(\arabic*)}, leftmargin=*,
	itemsep=4pt]
	
	\item \textbf{Rimozione dei vincoli ridondanti.}\index{vincolo!ridondante|textbf}
	Si rimuovono $g_{\mathscr{i}}$ vincoli, scelti in modo che il
	sistema risultante sia isostatico ($\nu = 0$,
	$\rg(\bm{A}_{\rm rid}) = n$). Le reazioni dei vincoli rimossi
	diventano le \emph{incognite iperstatiche}\index{iperstaticita@iperstaticità!incognita iperstatica|textbf} $\chi_1, \ldots,
	\chi_{g_{\mathscr{i}}}$.
	
	\item \textbf{Soluzione particolare: iperstatiche nulle.}
	Si pone $\bm{\chi} = \bm{0}$ (reazioni dei vincoli rimossi uguali
	a zero) e si risolve il problema statico del sistema isostatico
	ridotto con le forze attive $\bm{f}$ assegnate, applicando il
	metodo dei corpi liberi (\PARref{sec:corpi_liberi}). La soluzione
	$\bm{r}_0$ così ottenuta soddisfa l'equilibrio con i vincoli
	rimasti, ma non impone alcuna condizione sui vincoli rimossi.
	
	\item \textbf{Stati di autoequilibrio: forze attive nulle,
		iperstatiche unitarie.}
	Si pone $\bm{f} = \bm{0}$ e si assegna valore unitario a
	ciascuna iperstatica alla volta: $\chi_k = 1$, $\chi_j = 0$ per
	$j \neq k$. Ogni problema così ottenuto è un problema statico
	del sistema isostatico ridotto, in cui la reazione $\chi_k = 1$
	del vincolo rimosso agisce come forza attiva nota sui corpi
	adiacenti; si risolve con il metodo dei corpi liberi
	(\PARref{sec:corpi_liberi}). La soluzione
	$\bm{r}_k$ è uno stato di autoequilibrio del sistema originale
	(soddisfa $\bm{A}^{\top}\bm{r}_k = \bm{0}$) e costituisce la
	$k$-esima colonna di $\bm{R}_\chi$.
	
	\item \textbf{Soluzione generale.}
	Per linearità, la soluzione generale si ottiene sovrapponendo
	i contributi dei passi 2 e 3:
	\begin{equation}
		\bm{r} = \bm{r}_0 +
		\sum_{k=1}^{g_{\mathscr{i}}} \chi_k\,\bm{r}_k,
	\end{equation}
	con le iperstatiche $\chi_k$ libere. 
	
\end{enumerate}

\begin{oss}[Arbitrarietà della scelta dei vincoli ridondanti]
	La scelta di quali vincoli rimuovere è arbitraria, purché il
	sistema ridotto sia isostatico. Nella formulazione algebrica, i
	vincoli ridondanti corrispondono alla scelta delle variabili
	fuori base del sistema $\bm{A}^{\top}\bm{r} = -\bm{f}$: scelte
	diverse corrispondono a basi diverse per $\Ker(\bm{A}^{\top})$,
	cioè a iperstatiche che sono combinazioni lineari l'una
	dell'altra. L'insieme di tutte le soluzioni
	$\bm{r}_0 + \Ker(\bm{A}^{\top})$ è però invariante: soluzioni
	particolari e stati di autoequilibrio cambiano forma ma
	descrivono sempre lo stesso insieme di reazioni ammissibili.
	In pratica: si rimuovono i vincoli ridondanti in modo da rendere
	i calcoli il più semplice possibile.
\end{oss}



\begin{esempio}[Trave iperstatica di grado 1 --- problema statico]
	\label{ex:iper_stat_1}
	Si consideri una trave rigida vincolata da una cerniera in
	$A$, un carrello verticale in $B$ e un carrello verticale
	in $C$, punto intermedio tra $A$ e $B$ a distanza $\ell/2$
	da ciascuno. I vincoli scalari sono $m = 4$, $n = 3$,
	$\nu = -1$: il sistema è iperstatico di grado
	$g_{\mathscr{i}} = 1$. È applicata una forza verticale $F$
	verso il basso nel punto $P$, a distanza $a$ da $A$
	(\FIGref{fig:iper_stat_1_dati}). Determinare la soluzione
	generale del problema statico.
	
	
	\begin{figure}[ht]
		\centering
		\begin{tikzpicture}[>=Stealth, line width=0.8pt]
			
			\def\L{4.0}
			\def\a{1.5}
			\def\F{0.6}
			\def\yrea{0.55} % altezza frecce reazioni
			
			% ============================================================
			% PANNELLO (a) — sistema originale con carico
			% ============================================================
			\begin{scope}[xshift=0cm]
				
				% trave
				
				\draw[line width=1.5pt] (0,0) -- (\L,0);
				
				% cerniera in A
				\draw[thick] (0,0) -- (-0.22,-0.38) -- (0.22,-0.38) -- cycle;
				\draw[thick] (-0.32,-0.38) -- (0.32,-0.38);
				\foreach \x in {-0.24,-0.12,0,0.12,0.24}{
					\draw[thin] (\x,-0.38) -- (\x-0.10,-0.52);
				}
				
				
				% carrello verticale in C
				\draw[thick] (\L/2,0) -- (\L/2-0.22,-0.38)
				-- (\L/2+0.22,-0.38) -- cycle;
				\draw[thick] (\L/2-0.15,-0.44) circle (0.07);
				\draw[thick] (\L/2+0.15,-0.44) circle (0.07);
				\draw[thick] (\L/2-0.32,-0.52) -- (\L/2+0.32,-0.52);
				\foreach \x in {-0.24,-0.12,0,0.12,0.24}{
					\draw[thin] (\L/2+\x,-0.52) -- (\L/2+\x-0.10,-0.66);
				}
				
				% carrello verticale in B
				\draw[thick] (\L,0) -- (\L-0.22,-0.38)
				-- (\L+0.22,-0.38) -- cycle;
				\draw[thick] (\L-0.15,-0.44) circle (0.07);
				\draw[thick] (\L+0.15,-0.44) circle (0.07);
				\draw[thick] (\L-0.32,-0.52) -- (\L+0.32,-0.52);
				\foreach \x in {-0.24,-0.12,0,0.12,0.24}{
					\draw[thin] (\L+\x,-0.52) -- (\L+\x-0.10,-0.66);
				}
				
				% punti notevoli
				\node[circle, fill=white, draw=black, inner sep=0pt,
				minimum size=5pt, line width=0.8pt] at (0,0) {};
				\node[above, font=\scriptsize] at (0,0) {$A$};
				
				\node[circle, fill=white, draw=black, inner sep=0pt,
				minimum size=5pt, line width=0.8pt] at (\L,0) {};
				\node[above, font=\scriptsize] at (\L,0) {$B$};
				
				% mezzaluna in C (vincolo attivo)
				\fill[white] (\L/2-0.08,0) arc(180:360:0.08) -- cycle;
				\draw[thick]  (\L/2-0.08,0) arc(180:360:0.08);
				
				
				\node[above, font=\scriptsize] at (\L/2,0) {$C$};
				\node[below, font=\scriptsize] at (\a,0)   {$P$};
				
				
				% forza F in P
				\draw[->, thick] (\a, \F) -- (\a, 0)
				node[midway, left, font=\scriptsize] {$F$};
				
				% quota a
				\draw[<->, thin] (0,-0.9) -- (\a,-0.9)
				node[midway, above, font=\scriptsize] {$a$};
				
				% quota ell
				\draw[<->, thin] (0,-1.3) -- (\L,-1.3)
				node[midway, above, font=\scriptsize] {$\ell$};
				
				\node[font=\scriptsize] at (\L/2, -1.8) {(a)};
				
			\end{scope}
			
			% ============================================================
			% PANNELLO (b) — sistema principale e stato di autoequilibrio
			% ============================================================
			\begin{scope}[xshift=6.5cm]
				
				% trave
				\draw[line width=1.5pt] (0,0) -- (\L,0);
				
				% cerniera in A
				\draw[thick] (0,0) -- (-0.22,-0.38) -- (0.22,-0.38) -- cycle;
				\draw[thick] (-0.32,-0.38) -- (0.32,-0.38);
				\foreach \x in {-0.24,-0.12,0,0.12,0.24}{
					\draw[thin] (\x,-0.38) -- (\x-0.10,-0.52);
				}
				
				% carrello verticale in C (grigio — ridondante)
				\draw[gray!50, thick] (\L/2,0) -- (\L/2-0.22,-0.38)
				-- (\L/2+0.22,-0.38) -- cycle;
				\draw[gray!50, thick] (\L/2-0.15,-0.44) circle (0.07);
				\draw[gray!50, thick] (\L/2+0.15,-0.44) circle (0.07);
				\draw[gray!50, thick] (\L/2-0.32,-0.52)
				-- (\L/2+0.32,-0.52);
				\foreach \x in {-0.24,-0.12,0,0.12,0.24}{
					\draw[gray!50, thin] (\L/2+\x,-0.52)
					-- (\L/2+\x-0.10,-0.66);
				}
				
				
				% carrello verticale in B
				\draw[thick] (\L,0) -- (\L-0.22,-0.38)
				-- (\L+0.22,-0.38) -- cycle;
				\draw[thick] (\L-0.15,-0.44) circle (0.07);
				\draw[thick] (\L+0.15,-0.44) circle (0.07);
				\draw[thick] (\L-0.32,-0.52) -- (\L+0.32,-0.52);
				\foreach \x in {-0.24,-0.12,0,0.12,0.24}{
					\draw[thin] (\L+\x,-0.52) -- (\L+\x-0.10,-0.66);
				}
				
				
				% punti notevoli
				\node[circle, fill=white, draw=black, inner sep=0pt,
				minimum size=5pt, line width=0.8pt] at (0,0) {};
				\node[above, font=\scriptsize] at (0,0) {$A$};
				
				\node[circle, fill=white, draw=black, inner sep=0pt,
				minimum size=5pt, line width=0.8pt] at (\L,0) {};
				\node[above, font=\scriptsize] at (\L,0) {$B$};
				
				%				% mezzaluna in C (vincolo ridondante, semiluna grigia)
				\fill[white]          (\L/2-0.08,0) arc(180:360:0.08) -- cycle;
				\draw[gray!50, thick] (\L/2-0.08,0) arc(180:360:0.08);
				
				
				\node[above, font=\scriptsize] at (\L/2,0) {$C$};
				
				
				% reazioni stato di autoequilibrio
				% Y_A = -1/2 (verso il basso): freccia traslata sopra la trave
				\draw[->, thick] (0, 2*\yrea) -- (0, \yrea)
				node[midway, left, font=\scriptsize] {$\tfrac{1}{2}$};
				% Y_C = 1 (verso l'alto): freccia traslata sopra la trave
				\draw[->, thick] (\L/2, \yrea) -- (\L/2, 2*\yrea)
				node[midway, right, font=\scriptsize] {$1$};
				% Y_B = -1/2 (verso il basso): freccia traslata sopra la trave
				\draw[->, thick] (\L, 2*\yrea) -- (\L, \yrea)
				node[midway, right, font=\scriptsize] {$\tfrac{1}{2}$};
				
				\node[font=\scriptsize] at (\L/2, -1.8) {(b)};
			\end{scope}
			
		\end{tikzpicture}
		\caption{Trave iperstatica di grado 1
			(\ESEref{ex:iper_stat_1}):
			\emph{(a)}~sistema originale con carico $F$ in $P$;
			\emph{(b)}~sistema principale (carrello in $C$ in grigio)
			e stato di autoequilibrio $\bm{r}_1$: reazione unitaria
			verso l'alto in $C$, reazioni $1/2$ verso il basso in $A$
			e $B$.}
		\label{fig:iper_stat_1_dati}
	\end{figure}
	
	\begin{svolgimento}
		
		\esottotit{Sistema principale.}
		Si rimuove il carrello in $C$ (vincolo ridondante): il sistema
		principale è la trave con cerniera in $A$ e carrello in $B$,
		isostatico. La reazione del carrello rimosso diventa
		l'incognita iperstatica $\chi_1 \equiv Y_C$.
		
		\esottotit{Soluzione particolare} ($\chi_1 = 0$,
		carico $F$ assegnato).
		Si risolve l'equilibrio del sistema principale. Dal momento
		rispetto ad $A$:
		\[
		Y_{B_0} \cdot \ell - F \cdot a = 0
		\;\Longrightarrow\;
		Y_{B_0} = \frac{Fa}{\ell}
		\quad\text{(verso l'alto)}.
		\]
		Dall'equilibrio verticale:
		\[
		Y_{A_0}+ Y_{B_0} - F = 0
		\;\Longrightarrow\;
		Y_{A_0} = F\,\frac{\ell - a}{\ell}
		\quad\text{(verso l'alto)}.
		\]
		Dall'equilibrio orizzontale: $X_{A_0} = 0$.
		
		\esottotit{Stato di autoequilibrio} ($\bm{f} = \bm{0}$,
		$\chi_1 = 1$).
		Si applica al sistema principale la sola reazione unitaria
		$Y_C = 1$ verso l'alto in $C$. Dal momento rispetto ad $A$:
		\[
		Y_{B_1} \cdot \ell + 1 \cdot \frac{\ell}{2} = 0
		\;\Longrightarrow\;
		Y_{B_1} = -\frac{1}{2}
		\quad\text{(verso il basso)}.
		\]
		Dall'equilibrio verticale:
		\[
		Y_{A_1} + Y_{B_1} + 1 = 0
		\;\Longrightarrow\;
		Y_{A_1} = -\frac{1}{2}
		\quad\text{(verso il basso)}.
		\]
		Dall'equilibrio orizzontale: $X_{A_1}= 0$. Lo stato di
		autoequilibrio è (\FIGref{fig:iper_stat_1_dati}\,b):
		\[
		\bm{r}_1 =
		\Bigl\{0\,\;-\tfrac{1}{2}\,\;1\,\;-\tfrac{1}{2}\Bigr\}^{\!\top}
		\quad
		\bigl(X_A,\;Y_A,\;Y_C,\;Y_B\bigr).
		\]
		
		\esottotit{Soluzione generale.}
		Per sovrapposizione:
		\begin{align*}
			X_A &= 0, \\
			Y_A &= F\,\frac{\ell-a}{\ell}
			- \frac{\chi_1}{2}
			\quad\text{(verso l'alto per $\chi_1$ piccolo)},\\
			Y_C &= \chi_1, \\
			Y_B &= \frac{Fa}{\ell} - \frac{\chi_1}{2},
		\end{align*}
		con $\chi_1 \in \mathbb{R}$ libera. Lo stato reattivo
		equilibrato è dunque una famiglia mono-parametrica di
		reazioni: al variare di $\chi_1$ si ottengono infinite
		distribuzioni di reazioni che soddisfano tutte le
		equazioni di equilibrio.
		
		\begin{oss}[Arbitrarietà della scelta del vincolo ridondante]
			La scelta del vincolo ridondante non è unica. Un sistema
			principale è isostatico quando le reazioni rimaste sono
			esattamente tre e indipendenti: tante quante le equazioni
			di equilibrio disponibili (somma delle forze orizzontali,
			somma delle forze verticali, somma dei momenti rispetto a
			un polo). Per questo sistema si possono adottare almeno
			tre scelte:
			\begin{itemize}[noitemsep]
				\item rimozione del carrello in $C$ (scelta sviluppata
				sopra): restano $X_A$, $Y_A$, $Y_B$ --- le tre equazioni
				di equilibrio le determinano univocamente;
				\item rimozione del carrello in $B$: restano $X_A$,
				$Y_A$, $Y_C$ --- stessa struttura, isostatico;
				\item rilascio della componente verticale della cerniera
				in $A$: restano $X_A$, $Y_C$, $Y_B$ --- la somma delle
				forze orizzontali dà $X_A$, quella verticale e i momenti
				determinano $Y_C$ e $Y_B$.
			\end{itemize}
			In ciascun caso la soluzione particolare $\bm{r}_0$ e lo
			stato di autoequilibrio $\bm{r}_1$ cambiano forma, ma la
			famiglia di soluzioni $\bm{r}_0 + \chi_1\,\bm{r}_1$ descrive
			lo stesso insieme di stati reattivi equilibrati.
		\end{oss}
	\end{svolgimento}
\end{esempio}

\begin{esempio}[Trave iperstatica di grado 3 ---
	tre sistemi principali]
	\label{ex:iper_stat_3}
	Si consideri una trave rigida con incastro in $A$,
	cerniera in $B$ e carrello verticale in $C$, punto
	medio di $AB$. I vincoli scalari sono $m = 6$
	($3$ dall'incastro, $2$ dalla cerniera, $1$ dal
	carrello), $n = 3$, $\nu = -3$: il sistema è
	iperstatico di grado $g_{\mathscr{i}} = 3$.
	Individuare tre sistemi principali isostatici e
	descrivere i corrispondenti stati di autoequilibrio
	(\FIGref{fig:iper_stat_3}).
	
	\begin{figure}[ht]
		\centering
		\begin{tikzpicture}[>=Stealth, line width=0.8pt]
			
			\def\L{3.6}
			\def\dy{1.50}
			
			% ============================================================
			% SP1 — cerniera in A + carrello in C
			%        (rimossi: momento in A, cerniera in B)
			% ============================================================
			
			\begin{scope}[yshift=0cm]
				
				
				\draw[line width=1.5pt] (0,0) -- (\L,0);
				
				
				% incastro in A rimosso (grigio) — disegnato per primo
				\draw[gray!50, thick] (0,-0.38) -- (0, 0.38);
				\foreach \y in {-0.28,-0.14,0,0.14,0.28}{
					\draw[gray!50, thin] (0,\y) -- (-0.20,\y-0.13);
				}
				
				% cerniera in A (attiva) — triangolo
				\draw[thick] (0,0) -- (-0.20,-0.34)
				-- (0.20,-0.34) -- cycle;
				\draw[thick] (-0.30,-0.34) -- (0.30,-0.34);
				\foreach \x in {-0.22,-0.11,0,0.11,0.22}{
					\draw[thin] (\x,-0.34) -- (\x-0.10,-0.50);
				}
				
				% carrello verticale in C (attivo)
				\draw[thick] (\L/2,0) -- (\L/2-0.20,-0.34)
				-- (\L/2+0.20,-0.34) -- cycle;
				\draw[thick] (\L/2-0.14,-0.40) circle (0.07);
				\draw[thick] (\L/2+0.14,-0.40) circle (0.07);
				\draw[thick] (\L/2-0.30,-0.50) -- (\L/2+0.30,-0.50);
				\foreach \x in {-0.22,-0.11,0,0.11,0.22}{
					\draw[thin] (\L/2+\x,-0.50)
					-- (\L/2+\x-0.10,-0.64);
				}
				% mezzaluna sopra il vertice del triangolo in C
				\fill[white]      (\L/2-0.08,0) arc(180:360:0.08) -- cycle;
				\draw[thick]      (\L/2-0.08,0) arc(180:360:0.08);
				
				% cerniera in B rimossa (grigio)
				\draw[gray!50, thick] (\L,0) -- (\L-0.20,-0.34)
				-- (\L+0.20,-0.34) -- cycle;
				\draw[gray!50, thick] (\L-0.30,-0.34)
				-- (\L+0.30,-0.34);
				\foreach \x in {-0.22,-0.11,0,0.11,0.22}{
					\draw[gray!50, thin] (\L+\x,-0.34)
					-- (\L+\x-0.10,-0.50);
				}
				
				% pallini A e B — disegnati per ultimi, sopra i vincoli
				\draw[fill=white, thick]       (0,0)  circle (2pt);
				\draw[gray!50, fill=white, thick] (\L,0) circle (2pt);
				
				% etichette dei punti
				\node[above right,  font=\scriptsize]            at (0,0)    {$A$};
				\node[above,        font=\scriptsize]            at (\L/2,0) {$C$};
				\node[above, font=\scriptsize, gray!60]    at (\L,0)   {$B$};
				
				% etichetta del sistema principale e nomenclatura iperstatiche
				\node[font=\scriptsize, anchor=east] at (-0.5,0) {SP1};
				\node[font=\scriptsize, anchor=west] at (\L+0.3,0)
				{$\chi_1\!\equiv\!\mu_A,\;
					\chi_2\!\equiv\!X_B,\;
					\chi_3\!\equiv\!Y_B$};
				
			\end{scope}
			
			% ============================================================
			% SP2 — mensola (incastro in A)
			%        (rimossi: carrello in C, cerniera in B)
			% ============================================================
			\begin{scope}[yshift=-\dy cm]
				
				\draw[line width=1.5pt] (0,0) -- (\L,0);
				
				% incastro in A (attivo) — niente perno: l'incastro non ha pallino
				\draw[thick] (0,-0.38) -- (0, 0.38);
				\foreach \y in {-0.28,-0.14,0,0.14,0.28}{
					\draw[thin] (0,\y) -- (-0.20,\y-0.13);
				}
				
				% carrello verticale in C rimosso (grigio)
				\draw[gray!50, thick] (\L/2,0) -- (\L/2-0.20,-0.34)
				-- (\L/2+0.20,-0.34) -- cycle;
				\draw[gray!50, thick] (\L/2-0.14,-0.40) circle (0.07);
				\draw[gray!50, thick] (\L/2+0.14,-0.40) circle (0.07);
				\draw[gray!50, thick] (\L/2-0.30,-0.50)
				-- (\L/2+0.30,-0.50);
				\foreach \x in {-0.22,-0.11,0,0.11,0.22}{
					\draw[gray!50, thin] (\L/2+\x,-0.50)
					-- (\L/2+\x-0.10,-0.64);
				}
				% mezzaluna grigia sopra il vertice del triangolo in C (vincolo rimosso)
				\fill[white]            (\L/2-0.08,0) arc(180:360:0.08) -- cycle;
				\draw[gray!50, thick]   (\L/2-0.08,0) arc(180:360:0.08);
				
				% cerniera in B rimossa (grigio)
				\draw[gray!50, thick] (\L,0) -- (\L-0.20,-0.34)
				-- (\L+0.20,-0.34) -- cycle;
				\draw[gray!50, thick] (\L-0.30,-0.34)
				-- (\L+0.30,-0.34);
				\foreach \x in {-0.22,-0.11,0,0.11,0.22}{
					\draw[gray!50, thin] (\L+\x,-0.34)
					-- (\L+\x-0.10,-0.50);
				}
				
				% pallino in B — disegnato per ultimo, sopra il triangolo grigio
				\draw[gray!50, fill=white, thick] (\L,0) circle (2pt);
				
				% etichette dei punti
				\node[above right,  font=\scriptsize]              at (0,0)    {$A$};
				\node[above,        font=\scriptsize, gray!60]    at (\L/2,0) {$C$};
				\node[above, font=\scriptsize, gray!60]     at (\L,0)   {$B$};
				
				% etichetta del sistema principale e nomenclatura iperstatiche
				\node[font=\scriptsize, anchor=east] at (-0.5,0) {SP2};
				\node[font=\scriptsize, anchor=west] at (\L+0.3,0)
				{$\chi_1\!\equiv\!X_B,\;
					\chi_2\!\equiv\!Y_B,\;
					\chi_3\!\equiv\!Y_C$};
				
			\end{scope}
			
			% ============================================================
			% SP3 — cerniera in B + carrello in C
			%        (rimosso: incastro in A)
			% ============================================================
			\begin{scope}[yshift=-2*\dy cm]
				
				\draw[line width=1.5pt] (0,0) -- (\L,0);
				
				% incastro in A rimosso (grigio)
				\draw[gray!50, thick] (0,-0.38) -- (0, 0.38);
				\foreach \y in {-0.28,-0.14,0,0.14,0.28}{
					\draw[gray!50, thin] (0,\y) -- (-0.20,\y-0.13);
				}
				
				% carrello verticale in C (attivo)
				\draw[thick] (\L/2,0) -- (\L/2-0.20,-0.34)
				-- (\L/2+0.20,-0.34) -- cycle;
				\draw[thick] (\L/2-0.14,-0.40) circle (0.07);
				\draw[thick] (\L/2+0.14,-0.40) circle (0.07);
				\draw[thick] (\L/2-0.30,-0.50) -- (\L/2+0.30,-0.50);
				\foreach \x in {-0.22,-0.11,0,0.11,0.22}{
					\draw[thin] (\L/2+\x,-0.50)
					-- (\L/2+\x-0.10,-0.64);
				}
				% mezzaluna nera sopra il vertice del triangolo in C (vincolo attivo)
				\fill[white]      (\L/2-0.08,0) arc(180:360:0.08) -- cycle;
				\draw[thick]      (\L/2-0.08,0) arc(180:360:0.08);
				
				% cerniera in B (attiva)
				\draw[thick] (\L,0) -- (\L-0.20,-0.34)
				-- (\L+0.20,-0.34) -- cycle;
				\draw[thick] (\L-0.30,-0.34) -- (\L+0.30,-0.34);
				\foreach \x in {-0.22,-0.11,0,0.11,0.22}{
					\draw[thin] (\L+\x,-0.34) -- (\L+\x-0.10,-0.50);
				}
				
				% pallino in B — disegnato per ultimo, sopra il triangolo nero
				\draw[fill=white, thick] (\L,0) circle (2pt);
				
				% etichette dei punti
				\node[above right,  font=\scriptsize, gray!60]   at (0,0)    {$A$};
				\node[above,        font=\scriptsize]            at (\L/2,0) {$C$};
				\node[above, font=\scriptsize]             at (\L,0)   {$B$};
				
				% etichetta del sistema principale e nomenclatura iperstatiche
				\node[font=\scriptsize, anchor=east] at (-0.5,0) {SP3};
				\node[font=\scriptsize, anchor=west] at (\L+0.3,0)
				{$\chi_1\!\equiv\!X_A,\;
					\chi_2\!\equiv\!Y_A,\;
					\chi_3\!\equiv\!\mu_A$};
				
			\end{scope}
			
		\end{tikzpicture}
		\caption{Trave iperstatica di grado 3
			(\ESEref{ex:iper_stat_3}): tre sistemi principali
			isostatici. I vincoli rimossi sono in grigio;
			le corrispondenti reazioni diventano le incognite
			iperstatiche $\chi_1$, $\chi_2$, $\chi_3$.}
		\label{fig:iper_stat_3}
	\end{figure}
	
	
	\begin{svolgimento}
		
		Per ciascun sistema principale si verifica l'isostaticità
		e si descrivono i corrispondenti stati di autoequilibrio.
		
		\esottotit{SP1: cerniera in $A$ + carrello in $C$.}
		Si rilascia il momento dell'incastro ($\chi_1 \equiv \mu_A$)
		e si rimuove la cerniera in $B$
		($\chi_2 \equiv X_B$, $\chi_3 \equiv Y_B$): restano
		due reazioni in $A$ ($X_A$, $Y_A$) e una in $C$ ($Y_C$),
		tre reazioni indipendenti per tre equazioni di equilibrio.
		I tre stati di autoequilibrio si ottengono applicando
		al sistema principale le iperstatiche unitarie una
		alla volta:
		\begin{itemize}[noitemsep]
			\item $\chi_1 = 1$ ($\mu_A$ unitaria, $X_B = Y_B = 0$):
			coppia in $A$ equilibrata da una forza verticale in $C$
			e da una forza verticale in $A$, senza componente
			orizzontale;
			\item $\chi_2 = 1$ ($X_B$ unitaria, $\mu_A = Y_B = 0$):
			forza orizzontale in $B$ equilibrata dalla sola reazione
			orizzontale in $A$, senza reazioni verticali;
			\item $\chi_3 = 1$ ($Y_B$ unitaria, $\mu_A = X_B = 0$):
			forza verticale in $B$ equilibrata da forze verticali
			in $A$ e $C$.
		\end{itemize}
		
		\esottotit{SP2: mensola con incastro in $A$.}
		Si rimuovono la cerniera in $B$
		($\chi_1 \equiv X_B$, $\chi_2 \equiv Y_B$) e il carrello
		in $C$ ($\chi_3 \equiv Y_C$): restano le tre reazioni
		dell'incastro ($X_A$, $Y_A$, $\mu_A$), determinate
		dall'equilibrio. I tre stati di autoequilibrio:
		\begin{itemize}[noitemsep]
			\item $\chi_1 = 1$ ($X_B$ unitaria): forza orizzontale
			in $B$ equilibrata da $X_A$ uguale e opposta, senza
			effetti verticali né momento;
			\item $\chi_2 = 1$ ($Y_B$ unitaria): forza verticale
			in $B$ equilibrata da $Y_A$ e da un momento in $A$
			proporzionale alla distanza $AB$;
			\item $\chi_3 = 1$ ($Y_C$ unitaria): forza verticale
			in $C$ equilibrata da $Y_A$ e da un momento in $A$
			proporzionale alla distanza $AC = \ell/2$.
		\end{itemize}
		
		\esottotit{SP3: cerniera in $B$ + carrello in $C$.}
		Si rimuove l'intero incastro in $A$
		($\chi_1 \equiv X_A$, $\chi_2 \equiv Y_A$,
		$\chi_3 \equiv \mu(A)$): restano $X_B$, $Y_B$ e $Y_C$,
		tre reazioni indipendenti. I tre stati di autoequilibrio:
		\begin{itemize}[noitemsep]
			\item $\chi_1 = 1$ ($X_A$ unitaria): forza orizzontale
			in $A$ equilibrata da $X_B$ uguale e opposta;
			\item $\chi_2 = 1$ ($Y_A$ unitaria): forza verticale
			in $A$ equilibrata da $Y_B$ e $Y_C$ secondo le
			distanze;
			\item $\chi_3 = 1$ ($\mu_A$ unitaria): coppia in $A$
			equilibrata da forze verticali in $B$ e $C$ uguali
			e opposte, con bracci $\ell$ e $\ell/2$.
		\end{itemize}
		
		\smallskip
		Nonostante i tre sistemi principali producano stati di
		autoequilibrio di forma diversa, l'insieme di tutte le
		soluzioni generali descrive lo stesso spazio
		$\Ker(\bm{A}^{\top})$: basi diverse per lo stesso
		sottospazio di dimensione tre.
		
	\end{svolgimento}
\end{esempio}




\subsection{Problema cinematico: condizione di compatibilità}
\label{sec:iperstatico_cin}


Per un sistema iperstatico\index{compatibilita@compatibilità!cinematica}, il problema cinematico
$\bm{A}\bm{u} = \bm{s}$ è \emph{sovradeterminato}: ha $m > n$
equazioni e $n$ incognite. La soluzione esiste se e solo se
$\bm{s} \in \Imm(\Abb)$, ovvero (\CAPref{cap:problemi_cin_stat_SCR}):
\begin{equation}\label{eq:compat_cin}
	\bm{R}_\chi^{\top}\,\bm{s} = \bm{0},
	\qquad\text{ovvero}\qquad
	\bm{r}_k^{\top}\,\bm{s} = 0,
	\quad k = 1, 2, \ldots, g_{\mathscr{i}},
\end{equation}
ossia, per il principio dei lavori virtuali (\PARref{sec:PLV}),
i cedimenti devono compiere lavoro virtuale nullo su ogni
stato di autoequilibrio. Questa è la \emph{condizione di
	compatibilità cinematica}: un sistema iperstatico ammette una
configurazione compatibile con i cedimenti assegnati solo se
questi sono \emph{ortogonali} allo spazio degli stati di
autoequilibrio.

\paragraph{Interpretazione fisica.}
La condizione~\eqref{eq:compat_cin} è la scrittura matriciale
della compatibilità dei cedimenti con gli stati di autoequilibrio:
per ciascuna colonna $\bm{r}_k$ di $\bm{R}_\chi$, l'equazione
$\bm{r}_k^{\top}\bm{s} = 0$ esprime che i cedimenti assegnati
non compiono lavoro virtuale nelle corrispondenti reazioni di
autoequilibrio.

Quando la condizione è soddisfatta, la soluzione del problema
cinematico è \emph{unica} (poiché $g_\ell = 0$ nei sistemi
iperstatici non degeneri) e si determina con le stesse strategie
della \PARref{sec:determinato}, applicando il metodo
grafo-analitico al sistema isostatico ridotto con i cedimenti
assegnati.

\begin{esempio}[Trave iperstatica di grado 1 --- 
	compatibilità cinematica]
	\label{ex:iper_cin}
	Si riprende il sistema dell'\ESEref{ex:iper_stat_1}.
	Sono assegnati i cedimenti: spostamento orizzontale
	$\bar{u}_A$ alla cerniera in $A$, cedimento verticale
	$\bar{v}_A$ alla cerniera in $A$, cedimento verticale
	$\bar{v}_C$ al carrello in $C$, cedimento verticale
	$\bar{v}_B$ al carrello in $B$. Determinare la condizione
	di compatibilità cinematica.
	
	
	\begin{svolgimento}
		
		Il vettore dei cedimenti, ordinato come le componenti
		di $\bm{r}_1$, è:
		\[
		\bm{s} =
		\bigl(\bar{u}_A,\;\bar{v}_A,\;\bar{v}_C,\;
		\bar{v}_B\bigr)^{\!\top}
		\quad
		(X_A,\;Y_A,\;Y_C,\;Y_B).
		\]
		Lo stato di autoequilibrio costruito
		nell'\ESEref{ex:iper_stat_1} è:
		\[
		\bm{r}_1 =
		\Bigl(0,\;-\tfrac{1}{2},\;1,\;
		-\tfrac{1}{2}\Bigr)^{\!\top}.
		\]
		La condizione di compatibilità cinematica
		$\bm{r}_1^{\top}\bm{s} = 0$ diventa:
		\[
		0 \cdot \bar{u}_A
		- \tfrac{1}{2}\,\bar{v}_A
		+ 1 \cdot \bar{v}_C
		- \tfrac{1}{2}\,\bar{v}_B = 0,
		\]
		ovvero:
		\[
		\bar{v}_C = \frac{\bar{v}_A + \bar{v}_B}{2}.
		\]
		Il cedimento orizzontale $\bar{u}_A$ è libero: non compare
		nello stato di autoequilibrio perché la cerniera in $A$
		non trasmette reazioni orizzontali nel sistema principale.
		
		\esottotit{Interpretazione geometrica.}
		La condizione appena ottenuta esprime la linearità della
		deformata di un corpo rigido: il cedimento verticale del
		punto intermedio $C$ deve essere la media dei cedimenti
		verticali degli estremi $A$ e $B$. Un corpo rigido non
		può deformarsi, quindi i cedimenti dei suoi vincoli non
		possono essere assegnati arbitrariamente: devono essere
		compatibili con uno spostamento rigido, e questa è
		precisamente tale condizione.
		
	\end{svolgimento}
\end{esempio}

%----------------------------------------------------------------------------------------------------------------------------
\section{Il caso labile-iperstatico: condizioni operative di compatibilità}
\label{sec:labile_iperstatico}
%----------------------------------------------------------------------------------------------------------------------------

Quando la geometria dei vincoli è degenere il sistema può essere
simultaneamente labile ($g_\ell > 0$) e iperstatico
($g_{\mathscr{i}} > 0$), con $g_\ell = g_{\mathscr{i}}$ poiché
$g_\ell - g_{\mathscr{i}} = \nu$. Le cause geometriche di questa
condizione --- parallelismo o concorrenza degli assi di vincolo
per un singolo corpo, allineamento dei centri di rotazione per
sistemi a più corpi --- e i criteri per riconoscerla sono trattati
nel \CAPref{cap:classificazione_SCR}. Qui si assume che il tipo
sia già noto e si descrivono le condizioni operative di compatibilità
che entrambi i problemi devono soddisfare.

Entrambe le condizioni di compatibilità --- la condizione statica
\eqref{eq:compat_stat} sui carichi e quella cinematica
\eqref{eq:compat_cin} sui cedimenti --- devono essere verificate
simultaneamente: $\bm{u}_k^{\top}\bm{f} = 0$ per ogni meccanismo
($k = 1,\ldots,g_\ell$) e $\bm{r}_k^{\top}\bm{s} = 0$ per ogni stato
di autoequilibrio ($k = 1,\ldots,g_{\mathscr{i}}$). Se entrambe sono soddisfatte,
la soluzione del problema cinematico è non unica (a meno del
meccanismo) e quella del problema statico è non unica (a meno
dello stato di autoequilibrio): i parametri liberi rimangono
indeterminati nell'ambito dei corpi rigidi.

\begin{oss}[Natura distinta delle due indeterminazioni]
	Nel caso labile-iperstatico le due indeterminazioni non si
	compensano e hanno natura distinta. I $g_{\mathscr{i}}$
	parametri di autoequilibrio potranno essere determinati
	nel modello deformabile, dove la compatibilità delle
	deformazioni fornisce le equazioni aggiuntive mancanti.
	I $g_\ell$ parametri di meccanismo rimangono invece
	indeterminati anche nel modello deformabile: il sistema
	è cinematicamente mobile e nessuna ipotesi sulla
	deformabilità degli elementi può bloccarne il moto.
\end{oss}

\begin{esempio}[Trave con tre carrelli paralleli ---
	caso labile-iperstatico]
	\label{ex:labile_iperstatico_3carrelli}
	
	Si consideri una trave rigida di lunghezza $\ell$
	vincolata da tre carrelli verticali (asse $\ab_y$) in
	$A$, $C$, $B$, con $C$ punto medio di $AB$. Il sistema
	ha $n_c = 1$, $n = 3$, $m = 3$ e $\nu = 0$: la
	differenza fra il numero di vincoli e i gradi di libertà
	è nulla, condizione necessaria per l'isostaticità. La
	configurazione geometrica è tuttavia degenere: gli assi
	dei tre vincoli sono \emph{paralleli}, dunque
	$\rg(\bm{A}) = 2 < n$. Il sistema è simultaneamente
	labile e iperstatico, con
	$g_\ell = g_{\mathscr{i}} = n - \rg(\bm{A}) = 1$
	(\FIGref{fig:labile_iperstatico_3carrelli}).
	
	Sono applicate una forza $F_x$ orizzontale e una forza
	$F_y$ verticale (entrambe verso destra/basso) nel
	punto $C$. Sono inoltre assegnati cedimenti verticali
	$\bar{v}_A$, $\bar{v}_C$, $\bar{v}_B$ ai tre vincoli.
	Verificare le due condizioni di compatibilità e
	determinare, per quanto possibile, la soluzione dei due
	problemi.
	
	% ============================================================
	% FIGURA: Trave con tre carrelli paralleli — caso degenere
	% ============================================================
	\begin{figure}[ht]
		\centering
		\begin{tikzpicture}[>=Stealth, line width=0.8pt]
			
			
			\def\L{4.0}
			\def\F{0.9}     % ← era 0.6 (forze allungate)
			\def\sv{0.30}
			\def\rea{0.55}
			\def\ced{0.55}  % ← nuovo: lunghezza frecce cedimenti
			
		
			% ============================================================
			% PANNELLO (a) — dati: vincoli, cedimenti, carichi
			% ============================================================
			\begin{scope}[xshift=0cm]
				
				% trave
				\draw[line width=1.5pt] (0,0) -- (\L,0);
				
				% --- carrello verticale in A ---
				\draw[thick] (0,0) -- (-0.22,-0.38) -- (0.22,-0.38) -- cycle;
				\draw[thick] (-0.15,-0.44) circle (0.07);
				\draw[thick] ( 0.15,-0.44) circle (0.07);
				\draw[thick] (-0.32,-0.52) -- (0.32,-0.52);
				\foreach \x in {-0.24,-0.12,0,0.12,0.24}{\draw[thin] (\x,-0.52) -- (\x-0.10,-0.66);
				}
				
				% --- carrello verticale in C ---
				\draw[thick] (\L/2,0) -- (\L/2-0.22,-0.38)	-- (\L/2+0.22,-0.38) -- cycle;
				\draw[thick] (\L/2-0.15,-0.44) circle (0.07);
				\draw[thick] (\L/2+0.15,-0.44) circle (0.07);
				\draw[thick] (\L/2-0.32,-0.52) -- (\L/2+0.32,-0.52);
				\foreach \x in {-0.24,-0.12,0,0.12,0.24}{\draw[thin] (\L/2+\x,-0.52) -- (\L/2+\x-0.10,-0.66);
				}
				
				% --- carrello verticale in B ---
				\draw[thick] (\L,0) -- (\L-0.22,-0.38)	-- (\L+0.22,-0.38) -- cycle;
				\draw[thick] (\L-0.15,-0.44) circle (0.07);
				\draw[thick] (\L+0.15,-0.44) circle (0.07);
				\draw[thick] (\L-0.32,-0.52) -- (\L+0.32,-0.52);
				\foreach \x in {-0.24,-0.12,0,0.12,0.24}{\draw[thin] (\L+\x,-0.52) -- (\L+\x-0.10,-0.66);
				}
				
				% --- perni/mezzaluna ed etichette dei punti ---
				% anellino in A (estremità)
				\node[circle, fill=white, draw=black, inner sep=0pt,
				minimum size=5pt, line width=0.8pt] at (0,0) {};
				\node[above left, font=\scriptsize] at (0,0) {$A$};
				
				% mezzaluna in C (carrello a mezzeria, non interrompe la continuità rotazionale)
				\fill[white] (\L/2-0.08,0) arc(180:360:0.08) -- cycle;
				\draw[thick] (\L/2-0.08,0) arc(180:360:0.08);
				\node[above left, font=\scriptsize] at (\L/2,0.10) {$C$};
				
				% anellino in B (estremità)
				\node[circle, fill=white, draw=black, inner sep=0pt,
				minimum size=5pt, line width=0.8pt] at (\L,0) {};
				\node[above right, font=\scriptsize] at (\L,0) {$B$};
				
				% --- cedimenti: frecce in asse ai vincoli, verso l'alto ---
				\draw[->, thick] (0,0)    -- (0,\ced)	node[right, font=\scriptsize] {$\bar{v}_A$};
				\draw[->, thick] (\L/2,0) -- (\L/2,\ced)	node[right, font=\scriptsize] {$\bar{v}_C$};
				\draw[->, thick] (\L,0)   -- (\L,\ced)	node[right, font=\scriptsize] {$\bar{v}_B$};
				
				% --- forze attive in C ---
				\draw[->, line width=1.2pt] (\L/2,0) -- ({\L/2+\F},0)	node[above, font=\scriptsize] {$F_x$};
				\draw[->, line width=1.2pt] (\L/2,0) -- (\L/2,-\F)
				node[right, font=\scriptsize] {$F_y$};
				
				% --- quote sotto i carrelli (immutate) ---
				\draw[<->, thin] (0,-1.05) -- (\L/2,-1.05)	node[midway, above=1pt, font=\scriptsize] {$\ell/2$};
				\draw[<->, thin] (\L/2,-1.05) -- (\L,-1.05)	node[midway, above=1pt, font=\scriptsize] {$\ell/2$};
				
				% etichetta pannello
				\node[font=\scriptsize] at (\L/2,-1.55) {(a)};
				
			\end{scope}
			
			% ============================================================
			% PANNELLO (b) — meccanismo u_1 e autoequilibrio r_1
			% ============================================================
			\begin{scope}[xshift={\L+6.0cm}]
				
				% trave traslata orizzontalmente (meccanismo)
				\draw[gray!50, line width=1.0pt] (0,0) -- (\L,0);
				\draw[line width=1.5pt] (\sv,0) -- ({\L+\sv},0);
				
				% carrelli (grigi, in posizione originale)
				\foreach \xp in {0,\L/2,\L}{
					\draw[gray!50, thick] (\xp,0) -- (\xp-0.22,-0.38)
					-- (\xp+0.22,-0.38) -- cycle;
					\draw[gray!50, thick] (\xp-0.15,-0.44) circle (0.07);
					\draw[gray!50, thick] (\xp+0.15,-0.44) circle (0.07);
					\draw[gray!50, thick] (\xp-0.32,-0.52)
					-- (\xp+0.32,-0.52);
					\foreach \x in {-0.24,-0.12,0,0.12,0.24}{
						\draw[gray!50, thin] (\xp+\x,-0.52)
						-- (\xp+\x-0.10,-0.66);
					}
				}
				
				% perni vuoti grigi in A e B (estremità) e mezzaluna grigia in C (mezzeria),
				% sulla configurazione di riferimento
				\draw[gray!60, fill=white, thick] (0,0)  circle (2.5pt);
				\draw[gray!60, fill=white, thick] (\L,0) circle (2.5pt);
				
				\fill[white]          (\L/2-0.08,0) arc(180:360:0.08) -- cycle;
				\draw[gray!60, thick] (\L/2-0.08,0) arc(180:360:0.08);
				
				% etichette
				\node[gray!70, above,  font=\scriptsize] at (0,0)        {$A$};
				\node[gray!70, above left,       font=\scriptsize] at (\L/2,0.0)  {$C$};
				\node[gray!70, above, font=\scriptsize] at (\L,0)       {$B$};
				
				
				% anellini e mezzaluna sulla configurazione traslata
				\node[circle, fill=white, draw=black, inner sep=0pt,
				minimum size=5pt, line width=0.8pt] at (\sv,0) {};
				\fill[white] ({\L/2+\sv-0.08},0) arc(180:360:0.08) -- cycle;
				\draw[thick] ({\L/2+\sv-0.08},0) arc(180:360:0.08);
				\node[circle, fill=white, draw=black, inner sep=0pt,
				minimum size=5pt, line width=0.8pt] at ({\L+\sv},0) {};
				
				% freccia traslazione (meccanismo u_1): coda in A, punta in A'
				\draw[->, thick] (0,-1.05) -- (\sv,-1.05)
				node[right, font=\scriptsize] {$\bm{u}_1$: traslazione $\ab_x$};
				
				
				
				% reazioni di autoequilibrio r_1: -1/2 in A, +1 in C, -1/2 in B
				\draw[<-, thick] (0,-\rea) -- (0,0)
				node[below=12pt, font=\scriptsize] {$\tfrac{1}{2}$};
				\draw[<-, thick] (\L/2,\rea) -- (\L/2,0)
				node[above=12pt, font=\scriptsize] {$1$};
				\draw[<-, thick] (\L,-\rea) -- (\L,0)
				node[below=12pt, font=\scriptsize] {$\tfrac{1}{2}$};
				
				% etichette punti
				\fill (\sv+0.1,0)        node[below,  font=\scriptsize] {$A'$};
				\fill ({\L/2+\sv+0.1},0) node[below,       font=\scriptsize] {$C'$};
				\fill ({\L+\sv+0.1},0)   node[below, font=\scriptsize] {$B'$};
				
				\node[font=\scriptsize] at (\L/2,-1.55) {(b)};
				
			\end{scope}
			
		\end{tikzpicture}
		\caption{Trave con tre carrelli verticali paralleli: {(a)}
			dati del problema ; {(b)}
			meccanismo $\bm{u}_1$ (traslazione orizzontale) e
			stato di autoequilibrio $\bm{r}_1$
			(reazioni verticali a risultante e momento nulli)				.}
		\label{fig:labile_iperstatico_3carrelli}
	\end{figure}
	
	\begin{svolgimento}
		
		\esottotit{Meccanismo.}
		Tutti gli assi di vincolo sono paralleli ad $\ab_y$:
		ogni vincolo annulla la componente verticale dello
		spostamento del proprio punto di applicazione, ma
		lascia libera la componente orizzontale. L'unica
		configurazione cinematica ammessa è quindi una
		traslazione rigida orizzontale, con $\vartheta = 0$ e
		$\ub(P) = \alpha\,\ab_x$ per ogni punto $P$ del corpo,
		$\alpha \in \mathbb{R}$. Posto $\alpha = 1$ il
		meccanismo, riferito al polo $O \equiv A$, è:
		\[
		\bm{u}_1 =
		\bigl(1\,\;0\,\;0\bigr)^{\!\top}.
		\]
		Equivalentemente, il \CRa\ è all'infinito nella
		direzione $\ab_y$, ortogonale alla direzione di
		traslazione, come previsto dal caso limite del
		\PARref{sec:grafo_analitico}.
		
		\esottotit{Stato di autoequilibrio.}
		Si rimuove il carrello in $C$ e si assume la sua
		reazione $Y_C$ come iperstatica
		($\chi = Y_C$). Nel sistema ridotto (trave appoggiata
		in $A$ e $B$, isostatica) si pone $\chi = 1$ e si
		richiede equilibrio in assenza di forze attive.
		Dall'equilibrio alla rotazione attorno ad $A$ si
		ottiene
		$Y_B = -\tfrac{1}{2}$;
		dall'equilibrio alla traslazione verticale,
		$Y_A = -\tfrac{1}{2}$. Lo stato di autoequilibrio è:
		\[
		\bm{r}_1 =
		\bigl(0\,\;-\tfrac{1}{2}\,\;1\,\;-\tfrac{1}{2}\bigr)^{\!\top}
		\quad (X_A,\;Y_A,\;Y_C,\;Y_B),
		\]
		dove la prima componente è nulla perché nessun
		vincolo trasmette reazioni orizzontali. È lo stesso
		stato di autoequilibrio di
		\ESEref{ex:iper_stat_1}/\ESEref{ex:iper_cin}: la
		degenerazione geometrica non modifica $\Ker(\bm{A}^{\!\top})$,
		ma fa emergere in più un $\Ker(\bm{A})$ non banale.
		
		\esottotit{Condizione di compatibilità statica.}
		Il vettore delle forze attive riferito al polo
		$O \equiv A$ è
		$\bm{f} = \bigl(F_x\,\;-F_y\,\;-F_y\,\ell/2\bigr)^{\!\top}$.
		La condizione~\eqref{eq:compat_stat} per il meccanismo
		$\bm{u}_1$ è:
		\[
		\bm{u}_1^{\!\top}\bm{f}
		= 1\cdot F_x + 0\cdot(-F_y) + 0\cdot(-F_y\,\ell/2)
		= F_x = 0.
		\]
		L'equilibrio è possibile solo se la risultante
		orizzontale delle forze attive è nulla: la
		componente $F_x$ non può essere sopportata dai
		vincoli, che non trasmettono alcuna reazione
		orizzontale.
		
		\esottotit{Condizione di compatibilità cinematica.}
		Il vettore dei cedimenti, ordinato come $\bm{r}_1$,
		è $\bm{s} = \bigl(0\,\;\bar{v}_A\,\;\bar{v}_C\,\;\bar{v}_B\bigr)^{\!\top}$
		(la cerniera in $A$ è qui sostituita da un carrello,
		quindi non c'è cedimento orizzontale). La
		condizione~\eqref{eq:compat_cin} è:
		\[
		\bm{r}_1^{\!\top}\bm{s}
		= -\tfrac{1}{2}\,\bar{v}_A + \bar{v}_C
		- \tfrac{1}{2}\,\bar{v}_B = 0
		\;\Longrightarrow\;
		\bar{v}_C = \frac{\bar{v}_A + \bar{v}_B}{2}.
		\]
		Stessa condizione di \ESEref{ex:iper_cin}: i
		cedimenti devono essere allineati su una retta,
		coerentemente con la linearità della deformata di
		un corpo rigido.
		
		\esottotit{Soluzione dei due problemi.}
		Se entrambe le condizioni sono soddisfatte
		($F_x = 0$ e $\bar{v}_C = (\bar{v}_A+\bar{v}_B)/2$),
		i due problemi ammettono soluzione, ma in entrambi
		i casi essa è \emph{non unica}:
		\begin{itemize}[noitemsep]
			\item la configurazione cinematica è determinata
			a meno della traslazione orizzontale $\alpha\,\bm{u}_1$,
			con $\alpha \in \mathbb{R}$ libero;
			\item le reazioni vincolari sono determinate a
			meno dello stato di autoequilibrio $\chi\,\bm{r}_1$,
			con $\chi \in \mathbb{R}$ libero.
		\end{itemize}
		Le due indeterminazioni hanno natura distinta:
		$\chi$ verrà fissato nel modello deformabile imponendo
		la compatibilità delle deformazioni; $\alpha$
		rimane invece indeterminato anche nel modello
		deformabile, perché la trave è cinematicamente
		mobile in direzione orizzontale.
		
	\end{svolgimento}
\end{esempio}


%----------------------------------------------------------------------------------------------------------------------------
\section{Gerarchia strutturale}
\label{sec:gerarchia}
%----------------------------------------------------------------------------------------------------------------------------

Molti sistemi di corpi rigidi di interesse pratico non sono
semplicemente una collezione disordinata di corpi e vincoli, ma
presentano una \emph{struttura gerarchica}\index{gerarchia strutturale|textbf}: esiste un ordine
naturale in cui i corpi si appoggiano l'uno sull'altro, così che
il moto (o l'equilibrio) di ciascun corpo dipende da quello dei
corpi che lo precedono nella gerarchia, ma non viceversa.
Riconoscere questa struttura consente di risolvere i due problemi
in cascata, un corpo alla volta, evitando di affrontare il sistema
lineare completo.

% ──────────────────────────────────────────────────────────────
\subsection{Trascinante e trascinato nel problema cinematico}
\label{sec:trascinante}
% ──────────────────────────────────────────────────────────────

Nel problema cinematico si dice che il sottosistema
$\mathcal{S}_j$ è \emph{trascinato}\index{gerarchia strutturale!trascinato}\index{trascinato|see{gerarchia strutturale}} dal sottosistema
$\mathcal{S}_i$ quando la configurazione di $\mathcal{S}_j$ è
completamente determinata da quella di $\mathcal{S}_i$ tramite
i vincoli interni che li collegano, insieme agli eventuali
vincoli esterni di $\mathcal{S}_j$. Il sottosistema
$\mathcal{S}_i$ è il \emph{trascinante}\index{gerarchia strutturale!trascinante}\index{trascinante|see{gerarchia strutturale}}: il suo moto è noto
(o determinabile indipendentemente) e impone il moto di
$\mathcal{S}_j$. Ciascun sottosistema può essere un singolo
corpo oppure un insieme di corpi che formano internamente una
struttura determinata --- come un arco a tre cerniere --- e
che si comporta come un unico elemento cinematico nei confronti
del resto del sistema.

\paragraph{Risoluzione in cascata.}
In un sistema con gerarchia $\mathcal{S}_1 \to \mathcal{S}_2
\to \cdots \to \mathcal{S}_k$ (dove la freccia indica il
trascinamento), il problema cinematico si risolve nell'ordine:
\begin{enumerate}[label=\textbf{(\arabic*)}, leftmargin=*,
	itemsep=2pt]
	\item Si determina la configurazione di $\mathcal{S}_1$ dai
	suoi soli vincoli esterni.
	\item Si determina la configurazione di $\mathcal{S}_2$ dai
	vincoli interni con $\mathcal{S}_1$ (già nota) e dai vincoli
	esterni di $\mathcal{S}_2$.
	\item Si procede così fino a $\mathcal{S}_k$.
\end{enumerate}
Ad ogni passo si risolve un sistema cinematico relativo al solo
sottosistema corrente, invece del sistema completo.

\begin{oss}[Struttura a blocchi triangolare di $\bm{A}$]
	La gerarchia trascinante/trascinato si riflette direttamente
	nella struttura di $\bm{A}$: se i sottosistemi sono ordinati
	secondo la gerarchia, la matrice $\bm{A}$ assume una struttura
	a \emph{blocchi triangolare inferiore}\index{matrice dei vincoli!struttura a blocchi triangolare}, dove ciascun blocco
	corrisponde a un sottosistema (\FIGref{fig:blocchi_triangolari}). Le righe relative ai vincoli
	di $\mathcal{S}_j$ coinvolgono solo le componenti di
	configurazione di $\mathcal{S}_1, \ldots, \mathcal{S}_j$,
	mai quelle di $\mathcal{S}_{j+1}, \ldots, \mathcal{S}_k$.
	Il sistema $\bm{A}\bm{u} = \bm{s}$ si risolve quindi per
	sostituzione in avanti, un blocco alla volta.
	
	
	\begin{figure}[ht]
		\centering
		\begin{tikzpicture}[>=Stealth, line width=0.8pt,
			font=\small]
			
			% dimensioni blocchi
			\def\bone{1.0}
			\def\btwo{1.0}
			\def\bthree{1.0}
			% valori derivati precalcolati
			\pgfmathsetmacro{\bonebtwo}{\bone+\btwo}
			\pgfmathsetmacro{\bonebtwobthree}{\bone+\btwo+\bthree}
			\pgfmathsetmacro{\halftot}{(\bone+\btwo+\bthree)/2}
			\pgfmathsetmacro{\halfone}{\bone/2}
			\pgfmathsetmacro{\halftwo}{\btwo/2}
			\pgfmathsetmacro{\halfthree}{\bthree/2}
			
			\colorlet{diag}{gray!30}
			\colorlet{subdiag}{gray!15}
			\colorlet{zero}{white}
			
			% ============================================================
			% MATRICE A (triangolare inferiore)
			% ============================================================
			\begin{scope}[xshift=0cm]
				
				% blocchi diagonali
				\fill[diag] (0,-\bone) rectangle (\bone,0);
				\fill[diag] (\bone,-\bonebtwo)
				rectangle (\bonebtwo,-\bone);
				\fill[diag] (\bonebtwo,-\bonebtwobthree)
				rectangle (\bonebtwobthree,-\bonebtwo);
				
				% blocchi sotto-diagonali
				\fill[subdiag] (0,-\bonebtwo)
				rectangle (\bone,-\bone);
				\fill[subdiag] (0,-\bonebtwobthree)
				rectangle (\bone,-\bonebtwo);
				\fill[subdiag] (\bone,-\bonebtwobthree)
				rectangle (\bonebtwo,-\bonebtwo);
				
				% blocchi nulli
				\fill[zero] (\bone,-\bone)
				rectangle (\bonebtwobthree,0);
				\fill[zero] (\bonebtwo,-\bonebtwo)
				rectangle (\bonebtwobthree,-\bone);
				
				% griglia
				\draw[thick] (0,0)
				rectangle (\bonebtwobthree,-\bonebtwobthree);
				\draw[thick] (\bone,0) -- (\bone,-\bonebtwobthree);
				\draw[thick] (\bonebtwo,0)
				-- (\bonebtwo,-\bonebtwobthree);
				\draw[thick] (0,-\bone)
				-- (\bonebtwobthree,-\bone);
				\draw[thick] (0,-\bonebtwo)
				-- (\bonebtwobthree,-\bonebtwo);
				
				% etichette blocchi
				\node at (\halfone,-\halfone) {$\bm{A}_{11}$};
				\node at (\bone+\halftwo,-\bone-\halftwo)
				{$\bm{A}_{22}$};
				\node at (\bonebtwo+\halfthree,-\bonebtwo-\halfthree)
				{$\bm{A}_{33}$};
				\node at (\halfone,-\bone-\halftwo)
				{$\bm{A}_{21}$};
				\node at (\halfone,-\bonebtwo-\halfthree)
				{$\bm{A}_{31}$};
				\node at (\bone+\halftwo,-\bonebtwo-\halfthree)
				{$\bm{A}_{32}$};
				\node at (\bone+\halftwo,-\halfone) {$\bm{0}$};
				\node at (\bonebtwo+\halfthree,-\halfone) {$\bm{0}$};
				\node at (\bonebtwo+\halfthree,-\bone-\halftwo)
				{$\bm{0}$};
				
				% annotazioni colonne
				\draw[decorate,
				decoration={brace,amplitude=5pt,raise=4pt}]
				(0,0) -- (\bone,0)
				node[midway,above=9pt]
				{$\bm{u}_{\mathcal{S}_1}$};
				\draw[decorate,
				decoration={brace,amplitude=5pt,raise=4pt}]
				(\bone,0) -- (\bonebtwo,0)
				node[midway,above=9pt]
				{$\bm{u}_{\mathcal{S}_2}$};
				\draw[decorate,
				decoration={brace,amplitude=5pt,raise=4pt}]
				(\bonebtwo,0) -- (\bonebtwobthree,0)
				node[midway,above=9pt]
				{$\bm{u}_{\mathcal{S}_3}$};
				
				% annotazioni righe
				\draw[decorate,
				decoration={brace,amplitude=5pt,
					raise=4pt,mirror}]
				(0,0) -- (0,-\bone)
				node[midway,left=10pt,align=right,
				font=\scriptsize]
				{vincoli\\di $\mathcal{S}_1$};
				\draw[decorate,
				decoration={brace,amplitude=5pt,
					raise=4pt,mirror}]
				(0,-\bone) -- (0,-\bonebtwo)
				node[midway,left=10pt,align=right,
				font=\scriptsize]
				{vincoli int.\\
					$\mathcal{S}_1$-$\mathcal{S}_2$,\\
					est. di $\mathcal{S}_2$};
				\draw[decorate,
				decoration={brace,amplitude=5pt,
					raise=4pt,mirror}]
				(0,-\bonebtwo) -- (0,-\bonebtwobthree)
				node[midway,left=10pt,align=right,
				font=\scriptsize]
				{vincoli int.\\
					$\mathcal{S}_2$-$\mathcal{S}_3$,\\
					est. di $\mathcal{S}_3$};
				
				\node[left=60pt] at (0,-\halftot)
				{$\bm{A}=$};
				
			\end{scope}
			
			% ============================================================
			% MATRICE A^T (triangolare superiore)
			% ============================================================
			\begin{scope}[xshift=5.2cm]
				
				% blocchi diagonali
				\fill[diag] (0,-\bone) rectangle (\bone,0);
				\fill[diag] (\bone,-\bonebtwo)
				rectangle (\bonebtwo,-\bone);
				\fill[diag] (\bonebtwo,-\bonebtwobthree)
				rectangle (\bonebtwobthree,-\bonebtwo);
				
				% blocchi sopra-diagonali
				\fill[subdiag] (\bone,-\bone)
				rectangle (\bonebtwo,0);
				\fill[subdiag] (\bonebtwo,-\bone)
				rectangle (\bonebtwobthree,0);
				\fill[subdiag] (\bonebtwo,-\bonebtwo)
				rectangle (\bonebtwobthree,-\bone);
				
				% blocchi nulli
				\fill[zero] (0,-\bonebtwo)
				rectangle (\bone,-\bone);
				\fill[zero] (0,-\bonebtwobthree)
				rectangle (\bonebtwo,-\bonebtwo);
				
				% griglia
				\draw[thick] (0,0)
				rectangle (\bonebtwobthree,-\bonebtwobthree);
				\draw[thick] (\bone,0) -- (\bone,-\bonebtwobthree);
				\draw[thick] (\bonebtwo,0)
				-- (\bonebtwo,-\bonebtwobthree);
				\draw[thick] (0,-\bone)
				-- (\bonebtwobthree,-\bone);
				\draw[thick] (0,-\bonebtwo)
				-- (\bonebtwobthree,-\bonebtwo);
				
				% etichette blocchi
				\node at (\halfone,-\halfone)
				{$\bm{A}_{11}^{\top}$};
				\node at (\bone+\halftwo,-\bone-\halftwo)
				{$\bm{A}_{22}^{\top}$};
				\node at (\bonebtwo+\halfthree,-\bonebtwo-\halfthree)
				{$\bm{A}_{33}^{\top}$};
				\node at (\bone+\halftwo,-\halfone)
				{$\bm{A}_{21}^{\top}$};
				\node at (\bonebtwo+\halfthree,-\halfone)
				{$\bm{A}_{31}^{\top}$};
				\node at (\bonebtwo+\halfthree,-\bone-\halftwo)
				{$\bm{A}_{32}^{\top}$};
				\node at (\halfone,-\bone-\halftwo) {$\bm{0}$};
				\node at (\halfone,-\bonebtwo-\halfthree) {$\bm{0}$};
				\node at (\bone+\halftwo,-\bonebtwo-\halfthree)
				{$\bm{0}$};
				
				% annotazioni colonne
				\draw[decorate,
				decoration={brace,amplitude=5pt,raise=4pt}]
				(0,0) -- (\bone,0)
				node[midway,above=9pt]
				{$\bm{r}_{\mathcal{S}_1}$};
				\draw[decorate,
				decoration={brace,amplitude=5pt,raise=4pt}]
				(\bone,0) -- (\bonebtwo,0)
				node[midway,above=9pt]
				{$\bm{r}_{\mathcal{S}_2}$};
				\draw[decorate,
				decoration={brace,amplitude=5pt,raise=4pt}]
				(\bonebtwo,0) -- (\bonebtwobthree,0)
				node[midway,above=9pt]
				{$\bm{r}_{\mathcal{S}_3}$};
				
				% annotazioni righe
				\draw[decorate,
				decoration={brace,amplitude=5pt,raise=4pt}]
				(\bonebtwobthree,0) -- (\bonebtwobthree,-\bone)
				node[midway,right=10pt,align=left,
				font=\scriptsize]
				{equilibrio\\di $\mathcal{S}_1$};
				\draw[decorate,
				decoration={brace,amplitude=5pt,raise=4pt}]
				(\bonebtwobthree,-\bone)
				-- (\bonebtwobthree,-\bonebtwo)
				node[midway,right=10pt,align=left,
				font=\scriptsize]
				{equilibrio\\di $\mathcal{S}_2$};
				\draw[decorate,
				decoration={brace,amplitude=5pt,raise=4pt}]
				(\bonebtwobthree,-\bonebtwo)
				-- (\bonebtwobthree,-\bonebtwobthree)
				node[midway,right=10pt,align=left,
				font=\scriptsize]
				{equilibrio\\di $\mathcal{S}_3$};
				
				\node[left=10pt] at (0,-\halftot)
				{$\bm{A}^{\top}=$};
				
			\end{scope}
			
		\end{tikzpicture}
		\caption{Struttura a blocchi triangolare di $\bm{A}$
			(sinistra) e $\bm{A}^{\top}$ (destra) per un
			sistema con gerarchia
			$\mathcal{S}_1 \to \mathcal{S}_2 \to \mathcal{S}_3$.
			Blocchi diagonali (grigio scuro): vincoli propri
			di ciascun sottosistema. Blocchi fuori diagonale
			(grigio chiaro): vincoli di connessione.
			Il problema cinematico $\bm{A}\bm{u}=\bm{s}$ si
			risolve per sostituzione in avanti (da
			$\mathcal{S}_1$ a $\mathcal{S}_3$); quello statico
			$\bm{A}^{\top}\bm{r}=-\bm{f}$ per sostituzione
			all'indietro (da $\mathcal{S}_3$ a $\mathcal{S}_1$).}
		\label{fig:blocchi_triangolari}
	\end{figure}
	
	I blocchi hanno significati meccanici distinti.
	I \emph{blocchi diagonali} $\bm{A}_{jj}$ raccolgono
	le equazioni dei vincoli che, per $\mathcal{S}_j$,
	si comportano come vincoli con il suolo: sia i
	vincoli fisici con il suolo, sia i vincoli interni
	con i sottosistemi trascinanti, la cui configurazione
	è già nota al passo $j$ e si trasferisce nel termine
	noto $\bm{s}_j$ come cedimento assegnato. I
	\emph{blocchi fuori diagonale} $\bm{A}_{ji}$
	(con $i < j$) contengono i coefficienti con cui
	la configurazione di $\mathcal{S}_i$ compare nelle
	equazioni di vincolo di $\mathcal{S}_j$: essi
	quantificano il cedimento trasmesso da
	$\mathcal{S}_i$ a $\mathcal{S}_j$ attraverso i
	vincoli interni di connessione. I vincoli tra
	$\mathcal{S}_j$ e i sottosistemi che esso trascina
	compaiono invece nelle righe dei blocchi
	$\bm{A}_{ij}$ con $i > j$.
	
	La sostituzione in avanti rispecchia questa
	struttura: al passo $j$ si risolve il sistema
	\begin{equation}
		\bm{A}_{jj}\,\bm{u}_{\mathcal{S}_j} =
		\bm{s}_j -
		\sum_{i=1}^{j-1}
		\bm{A}_{ji}\,\bm{u}_{\mathcal{S}_i},
	\end{equation}
	dove $\bm{s}_j$ contiene tutti i cedimenti
	assegnati ai vincoli le cui equazioni compaiono
	nel blocco $j$: i cedimenti ai vincoli con il
	suolo fisico e i cedimenti relativi ai vincoli
	interni con i sottosistemi trascinanti. Il termine
	$\bm{A}_{ji}\,\bm{u}_{\mathcal{S}_i}$ rappresenta
	invece il contributo cinematico trasmesso dai
	sottosistemi trascinanti attraverso i vincoli
	interni: non è un dato del problema ma si
	determina man mano che si risolvono i blocchi
	precedenti.
\end{oss}

\subsection{Portante e portato nel problema statico}
\label{sec:portante}


Nel problema statico la gerarchia si inverte: il sottosistema
$\mathcal{S}_j$ è \emph{portato}\index{gerarchia strutturale!portato}\index{portato|see{gerarchia strutturale}} da $\mathcal{S}_i$ quando le
reazioni che $\mathcal{S}_j$ trasmette a $\mathcal{S}_i$
attraverso i vincoli interni sono determinabili dall'equilibrio
di $\mathcal{S}_j$ prima di affrontare l'equilibrio di
$\mathcal{S}_i$. Il sottosistema $\mathcal{S}_i$ è il
\emph{portante}\index{gerarchia strutturale!portante}\index{portante|see{gerarchia strutturale}}: riceve i carichi trasmessi da $\mathcal{S}_j$
e li scarica sui propri vincoli esterni. Ciascun sottosistema
può essere un singolo corpo oppure un insieme di corpi che
formano internamente una struttura determinata e che si
comporta come un unico elemento statico nei confronti del
resto del sistema.

\paragraph{Risoluzione in cascata.}
In un sistema con gerarchia $\mathcal{S}_1 \leftarrow
\mathcal{S}_2 \leftarrow \cdots \leftarrow \mathcal{S}_k$
(dove la freccia indica il portamento, dal portato al
portante), il problema statico si risolve nell'ordine
\emph{inverso} rispetto alla cinematica:
\begin{enumerate}[label=\textbf{(\arabic*)}, leftmargin=*,
	itemsep=2pt]
	\item Si determina l'equilibrio di $\mathcal{S}_k$
	dai suoi soli vincoli esterni e dai carichi applicati.
	Le reazioni interne trasmesse a $\mathcal{S}_{k-1}$
	diventano dati noti.
	\item Si determina l'equilibrio di $\mathcal{S}_{k-1}$,
	trattando le reazioni interne appena calcolate come
	forze attive note.
	\item Si procede così fino a $\mathcal{S}_1$.
\end{enumerate}
Ad ogni passo si risolve un problema statico relativo al
solo sottosistema corrente, invece del sistema completo.

\begin{oss}[Struttura a blocchi triangolare
	di $\bm{A}^{\top}$]
	L'inversione dell'ordine di risoluzione è la
	manifestazione del dualismo $\bm{B} = \bm{A}^{\top}$:
	se $\bm{A}$ è triangolare inferiore a blocchi,
	$\bm{A}^{\top}$ è triangolare \emph{superiore} a
	blocchi (\FIGref{fig:blocchi_triangolari}).
	
	I blocchi hanno significati meccanici distinti,
	duali di quelli del problema cinematico. I
	\emph{blocchi diagonali} $\bm{A}_{jj}^{\top}$
	raccolgono le equazioni di equilibrio che
	coinvolgono esclusivamente le reazioni dei vincoli
	di $\mathcal{S}_j$ con il suo ``suolo'': sia i
	vincoli fisici con il suolo, sia i vincoli interni
	con i sottosistemi portati, le cui reazioni sono
	già note al passo $j$ e si trasferiscono nel
	termine noto $-\bm{f}_j$ come forze attive
	assegnate. I \emph{blocchi fuori diagonale}
	$\bm{A}_{ji}^{\top}$ (con $i < j$, cioè sopra la
	diagonale in $\bm{A}^\top$) contengono i
	coefficienti con cui le reazioni di $\mathcal{S}_i$
	compaiono nelle equazioni di equilibrio di
	$\mathcal{S}_j$: essi quantificano la forza
	trasmessa dai sottosistemi portati a $\mathcal{S}_j$
	attraverso i vincoli interni di connessione.
	
	La sostituzione all'indietro rispecchia questa
	struttura: al passo $j$ (procedendo da $k$ a $1$)
	si risolve il sistema
	\begin{equation}
		\bm{A}_{jj}^{\top}\,\bm{r}_{\mathcal{S}_j} =
		-\bm{f}_j -
		\sum_{i=j+1}^{k}
		\bm{A}_{ij}^{\top}\,\bm{r}_{\mathcal{S}_i},
	\end{equation}
	dove $-\bm{f}_j$ contiene tutti i termini noti
	delle equazioni di equilibrio di $\mathcal{S}_j$:
	le forze attive applicate direttamente su
	$\mathcal{S}_j$ e le reazioni trasmesse dai
	vincoli interni con il suolo fisico. Il termine
	$\bm{A}_{ij}^{\top}\,\bm{r}_{\mathcal{S}_i}$
	rappresenta invece le reazioni trasmesse dai
	sottosistemi portati, già determinate nei passi
	precedenti.
	
	Cinematica e statica si risolvono lungo la stessa
	gerarchia, ma in direzioni opposte: in avanti per
	la cinematica (dal trascinante al trascinato), 
	all'indietro per la statica (dal portato al
	portante).
\end{oss}





\begin{esempio}[Trave di Gerber --- gerarchia strutturale]
	\label{ex:gerber_gerarchia}
	Si consideri la trave di Gerber\index{trave!di Gerber}: $\mathcal{C}_1$
	con cerniera in $A$ e carrello verticale in $B$;
	$\mathcal{C}_2$ con cerniera interna in $C$ e
	carrello verticale in $D$. I poli sono
	$O_1 \equiv A$ e $O_2 \equiv C$.
	
	\smallskip
	\noindent\textit{Gerarchia cinematica.}
	$\mathcal{C}_1$ è il trascinante: la cerniera
	in $A$ e il carrello in $B$ forniscono tre
	vincoli scalari che determinano univocamente
	la configurazione di $\mathcal{C}_1$
	indipendentemente dal resto del sistema.
	$\mathcal{C}_2$ è il trascinato: una volta nota
	la configurazione di $\mathcal{C}_1$, la
	cerniera interna in $C$ trasmette lo spostamento
	di $C$ come dato noto, e il carrello in $D$
	fornisce l'equazione mancante per determinare
	$\vartheta_2$. La cascata procede da
	$\mathcal{C}_1$ a $\mathcal{C}_2$.
	
	\smallskip
	\noindent\textit{Gerarchia statica.}
	$\mathcal{C}_2$ è il portato: ha il solo
	vincolo esterno in $D$ e la reazione interna
	in $C$. L'equilibrio di $\mathcal{C}_2$ ---
	con polo in $C$ per eliminare le reazioni
	interne --- determina $R_D$; le componenti
	della reazione interna in $C$ diventano
	forze attive note per $\mathcal{C}_1$.
	$\mathcal{C}_1$ è il portante: riceve le
	reazioni interne da $\mathcal{C}_2$ e le
	scarica sulla cerniera in $A$ e sul carrello
	in $B$. La cascata procede da $\mathcal{C}_2$
	a $\mathcal{C}_1$, in ordine inverso rispetto
	alla cinematica.
	
	Le due direzioni di cascata sono illustrate
	nella \FIGref{fig:gerarchia}.
	
	\begin{figure}[ht]
		\centering
		\begin{tikzpicture}[>=Stealth, line width=0.8pt]
			
			% ============================================================
			% PANNELLO (a) — gerarchia cinematica
			% ============================================================
			\begin{scope}[xshift=0cm]
				
				
				% linee d'asse (più spesse)
				\draw[line width=1.5pt] (0,0) -- (3,0);
				\draw[line width=1.5pt] (3,0) -- (5,0);
				
				% etichette corpi
				\node[above, font=\scriptsize] at (1.,0) {$\mathcal{C}_1$};
				\node[above, font=\scriptsize] at (4.0,0) {$\mathcal{C}_2$};
				
				% etichette punti notevoli
				\node[font=\scriptsize, above] at (0,0) {$A$};
				\node[font=\scriptsize, above] at (2,0) {$B$};
				\node[font=\scriptsize, above] at (3,0) {$C$};
				\node[font=\scriptsize, above] at (5,0) {$D$};
				
				% cerniera in A
				\draw[thick] (0,0) -- (-0.22,-0.38) -- (0.22,-0.38) -- cycle;
				\draw[thick] (-0.32,-0.38) -- (0.32,-0.38);
				\foreach \x in {-0.24,-0.12,0,0.12,0.24}{
					\draw[thin] (\x,-0.38) -- (\x-0.10,-0.52);
				}
				% perno bianco in A (sopra il vertice del triangolo)
				\draw[fill=white, thick] (0,0) circle (2.5pt);
				
				% carrello in B (intermedio)
				\draw[thick] (2,0) -- (2-0.22,-0.38) -- (2+0.22,-0.38) -- cycle;
				\draw[thick] (2-0.32,-0.38) -- (2+0.32,-0.38);
				\draw[thick] (2-0.15,-0.44) circle (0.07);
				\draw[thick] (2+0.15,-0.44) circle (0.07);
				\draw[thick] (2-0.32,-0.54) -- (2+0.32,-0.54);
				\foreach \x in {-0.24,-0.12,0,0.12,0.24}{
					\draw[thin] (2+\x,-0.54) -- (2+\x-0.10,-0.68);
				}
				% mezzaluna in B (carrello intermedio: la trave attraversa il vincolo)
				\fill[white] (2-0.08,0) arc(180:360:0.08) -- cycle;
				\draw[thick] (2-0.08,0) arc(180:360:0.08);
				
				% cerniera interna in C (cerchio bianco di taglia maggiore del perno)
				\draw[fill=white, thick] (3,0) circle (2.5pt);
				
				% carrello in D (terminale)
				\draw[thick] (5,0) -- (5-0.22,-0.38) -- (5+0.22,-0.38) -- cycle;
				\draw[thick] (5-0.32,-0.38) -- (5+0.32,-0.38);
				\draw[thick] (5-0.15,-0.44) circle (0.07);
				\draw[thick] (5+0.15,-0.44) circle (0.07);
				\draw[thick] (5-0.32,-0.54) -- (5+0.32,-0.54);
				\foreach \x in {-0.24,-0.12,0,0.12,0.24}{
					\draw[thin] (5+\x,-0.54) -- (5+\x-0.10,-0.68);
				}
				% perno bianco in D (carrello terminale)
				\draw[fill=white, thick] (5,0) circle (2.5pt);
				
				
				% freccia cascata cinematica (sinistra -> destra)
				\draw[->, very thick, gray!60]
				(0.2, 0.75) -- (4.8, 0.75);
				\node[font=\scriptsize, gray!70] at (2.5, 1.05)
				{trascinante $\to$ trascinato};
				
				
				\node[font=\scriptsize] at (3.0,-0.80) {(a)};
				
			\end{scope}
			
			% ============================================================
			% PANNELLO (b) — gerarchia statica
			% ============================================================
			\begin{scope}[xshift=7.0cm]
				
				% linee d'asse (più spesse)
				\draw[line width=1.5pt] (0,0) -- (3,0);
				\draw[line width=1.5pt] (3,0) -- (5,0);
				
				% etichette corpi
				\node[above, font=\scriptsize] at (1.0,0) {$\mathcal{C}_1$};
				\node[above, font=\scriptsize] at (4.0,0) {$\mathcal{C}_2$};
				
				% etichette punti notevoli
				\node[font=\scriptsize, above] at (0,0) {$A$};
				\node[font=\scriptsize, above] at (2,0) {$B$};
				\node[font=\scriptsize, above] at (3,0) {$C$};
				\node[font=\scriptsize, above] at (5,0) {$D$};
				
				% cerniera in A
				\draw[thick] (0,0) -- (-0.22,-0.38) -- (0.22,-0.38) -- cycle;
				\draw[thick] (-0.32,-0.38) -- (0.32,-0.38);
				\foreach \x in {-0.24,-0.12,0,0.12,0.24}{
					\draw[thin] (\x,-0.38) -- (\x-0.10,-0.52);
				}
				% perno bianco in A (sopra il vertice del triangolo)
				\draw[fill=white, thick] (0,0) circle (2.5pt);
				
				% carrello in B (intermedio)
				\draw[thick] (2,0) -- (2-0.22,-0.38) -- (2+0.22,-0.38) -- cycle;
				\draw[thick] (2-0.32,-0.38) -- (2+0.32,-0.38);
				\draw[thick] (2-0.15,-0.44) circle (0.07);
				\draw[thick] (2+0.15,-0.44) circle (0.07);
				\draw[thick] (2-0.32,-0.54) -- (2+0.32,-0.54);
				\foreach \x in {-0.24,-0.12,0,0.12,0.24}{
					\draw[thin] (2+\x,-0.54) -- (2+\x-0.10,-0.68);
				}
				% mezzaluna in B (carrello intermedio: la trave attraversa il vincolo)
				\fill[white] (2-0.08,0) arc(180:360:0.08) -- cycle;
				\draw[thick] (2-0.08,0) arc(180:360:0.08);
				
				% cerniera interna in C (cerchio bianco di taglia maggiore del perno)
				\draw[fill=white, thick] (3,0) circle (2.5pt);
				
				% carrello in D (terminale)
				\draw[thick] (5,0) -- (5-0.22,-0.38) -- (5+0.22,-0.38) -- cycle;
				\draw[thick] (5-0.32,-0.38) -- (5+0.32,-0.38);
				\draw[thick] (5-0.15,-0.44) circle (0.07);
				\draw[thick] (5+0.15,-0.44) circle (0.07);
				\draw[thick] (5-0.32,-0.54) -- (5+0.32,-0.54);
				\foreach \x in {-0.24,-0.12,0,0.12,0.24}{
					\draw[thin] (5+\x,-0.54) -- (5+\x-0.10,-0.68);
				}
				% perno bianco in D (carrello terminale)
				\draw[fill=white, thick] (5,0) circle (2.5pt);
				
				% freccia cascata statica (destra -> sinistra)
				\draw[->, very thick, gray!60]
				(4.8, 0.75) -- (0.2, 0.75);
				\node[font=\scriptsize, gray!70] at (2.5, 1.05)
				{portato $\to$ portante};
				
				\node[font=\scriptsize] at (3.0,-0.80) {(b)};
				
			\end{scope}
			
		\end{tikzpicture}
		\caption{Gerarchia strutturale della trave di Gerber:
			\emph{(a)}~cascata cinematica, dal trascinante
			$\mathcal{C}_1$ al trascinato $\mathcal{C}_2$
			(freccia da sinistra a destra);
			\emph{(b)}~cascata statica, dal portato
			$\mathcal{C}_2$ al portante $\mathcal{C}_1$
			(freccia da destra a sinistra). Le due cascate
			percorrono la stessa gerarchia in direzioni
			opposte.}
		\label{fig:gerarchia}
	\end{figure}
\end{esempio}


\subsection{Riconoscimento della gerarchia per ispezione}
\label{sec:riconoscimento_gerarchia}


Per riconoscere la gerarchia strutturale di un sistema per ispezione
diretta, senza costruire $\bm{A}$, si può applicare il seguente
criterio:

\begin{itemize}[noitemsep, itemsep=4pt]
	\item Si individua il sottosistema (corpo singolo o insieme di corpi)
	che ha vincoli esterni sufficienti a determinarne la configurazione
	indipendentemente dal resto: questo è il sottosistema
	\emph{trascinante} (nella cinematica) o \emph{portante} (nella
	statica).
	\item Si rimuove concettualmente quel sottosistema dal sistema
	globale, sostituendo i vincoli interni di connessione con cedimenti
	noti (cinematica) o con forze note (statica).
	\item Si ripete il procedimento sul sistema ridotto, individuando
	il nuovo sottosistema trascinante/portante.
	\item Si continua fino a esaurire tutti i sottosistemi.
\end{itemize}

Se il procedimento ha successo, la sequenza dei sottosistemi
individuati costituisce la gerarchia strutturale del sistema. Il
sottosistema trascinante/portante può essere un singolo corpo oppure
un insieme di corpi che formano internamente una struttura
determinata --- come un arco a tre cerniere o un anello a tre
cerniere --- e che si comporta come un unico elemento nei confronti
del resto del sistema. Se a un certo passo nessun sottosistema è
determinabile indipendentemente dagli altri, il sistema non ha una
struttura gerarchica semplice e occorre risolvere il problema
completo.



\begin{esempio}[Arco a tutto sesto con incastro ---	cinematica e statica]
	\label{ex:arco_incastro}\index{arco!a tutto sesto}
	Si consideri un arco a profilo circolare di raggio $R$,
	composto da due semiarchi $\mathcal{C}_1$ (sinistro) e
	$\mathcal{C}_2$ (destro). $\mathcal{C}_1$ è vincolato
	al suolo da un incastro in $A$; $\mathcal{C}_2$ è
	vincolato al suolo da un carrello verticale in $B$.
	I due semiarchi sono collegati da una cerniera interna
	in chiave $C$. Si scelgono come poli $O_1 \equiv A$
	e $O_2 \equiv B$ (\FIGref{fig:arco_incastro_dati}).
	
	\smallskip
	\noindent\textit{Problema cinematico.}
	Sono assegnati un cedimento verticale $\bar{v}_A$ verso
	il basso all'incastro in $A$ e un cedimento verticale
	$\bar{v}_B$ verso l'alto al carrello in $B$. Determinare
	la configurazione del sistema.
	
	\smallskip
	\noindent\textit{Problema statico.}
	Sono applicate una coppia $\mu$ antioraria su
	$\mathcal{C}_1$ e una forza verticale $F$ verso il basso
	su $\mathcal{C}_2$, a distanza orizzontale $R/2$ da $B$.
	Determinare le reazioni vincolari.
	
	% ============================================================
	% FIGURA: Arco a tutto sesto con incastro - dati cin/stat
	% ============================================================
	\begin{figure}[ht]
		\centering
		\begin{tikzpicture}[>=Stealth, line width=0.8pt]
			
			\def\R{2.0}
			\def\arr{0.55}
			
			% ============================================================
			% PANNELLO (a) — Dato cinematico
			% ============================================================
			\begin{scope}[xshift=0cm]
				
				% semiarchi
				\draw[line width=1.5pt] (-\R,0)
				arc[start angle=180, end angle=90, radius=\R];
				\draw[line width=1.5pt] (0,\R)
				arc[start angle=90, end angle=0, radius=\R];
				
				
				% incastro in A (suolo orizzontale + tratteggio obliquo verso il basso)
				\draw[thick] (-\R-0.38, 0) -- (-\R+0.38, 0);
				\foreach \x in {-0.28,-0.14,0,0.14,0.28}{
					\draw[thin] (-\R+\x, 0) -- (-\R+\x-0.13, -0.20);
				}
				
				
				% carrello verticale in B (triangolo + 2 ruote + suolo + tratteggio)
				\draw[thick] (\R,0) -- (\R-0.28,-0.48)
				-- (\R+0.28,-0.48) -- cycle;
				\draw[thick] (\R-0.18,-0.55) circle (0.08);
				\draw[thick] (\R+0.18,-0.55) circle (0.08);
				\draw[thick] (\R-0.38,-0.65) -- (\R+0.38,-0.65);
				\foreach \x in {-0.30,-0.15,0,0.15,0.30}{
					\draw[thin] (\R+\x,-0.65) -- (\R+\x-0.12,-0.83);
				}
				
				% anellino in B (perno tra corpo e carrello)
				\node[circle, fill=white, draw=black, inner sep=0pt,
				minimum size=5pt, line width=0.8pt] at (\R,0) {};
				
				% cerniera interna in C — perno unico
				\node[circle, fill=white, draw=black, inner sep=0pt,
				minimum size=5pt, line width=0.8pt] at (0,\R) {};
				
				% etichette punti
				\node[left, font=\scriptsize] at (-\R-0.22,+0.15)
				{$A \!\equiv\! O_1$};
				\node[left, font=\scriptsize] at (\R,+0.20)
				{$B \!\equiv\! O_2$};
				\node[above, font=\scriptsize] at (0,\R+0.10) {$C$};
				
				% etichette corpi (esterne, parte bassa per non collidere)
				\node[font=\small] at ({(\R+0.35)*cos(150)},{(\R+0.35)*sin(150)})
				{$\mathcal{C}_1$};
				\node[font=\small] at ({(\R+0.35)*cos(30)},{(\R+0.35)*sin(30)})
				{$\mathcal{C}_2$};
				
				
				% retta ausiliaria da A all'origine di bar{v}_A
				\draw[gray!50, thin] (-\R,0) -- (-\R+0.50,0);
				% cedimento bar{v}_A verso il basso (decentrato a destra dell'incastro)
				\draw[->, thick] (-\R+0.50,0) -- (-\R+0.50,-\arr)
				node[right, font=\scriptsize] {$\bar{v}_A$};
				
				% retta ausiliaria da B all'origine di bar{v}_B
				\draw[gray!50, thin] (\R,0) -- (\R+0.55,0);
				% cedimento bar{v}_B verso l'alto (decentrato a destra del carrello)
				\draw[->, thick] (\R+0.55,0) -- (\R+0.55,\arr)
				node[left, font=\scriptsize] {$\bar{v}_B$};
				
				% etichetta pannello
				\node[font=\scriptsize] at (0,-0.50) {(a)};
				
			\end{scope}
			
			% ============================================================
			% PANNELLO (b) — Dato statico
			% ============================================================
			\begin{scope}[xshift=7.0cm]
				
				% semiarchi
				\draw[line width=1.5pt] (-\R,0)
				arc[start angle=180, end angle=90, radius=\R];
				\draw[line width=1.5pt] (0,\R)
				arc[start angle=90, end angle=0, radius=\R];
				
				% incastro in A (suolo orizzontale + tratteggio obliquo verso il basso)
				\draw[thick] (-\R-0.38, 0) -- (-\R+0.38, 0);
				\foreach \x in {-0.28,-0.14,0,0.14,0.28}{
					\draw[thin] (-\R+\x, 0) -- (-\R+\x-0.13, -0.20);
				}
				
				% carrello verticale in B
				\draw[thick] (\R,0) -- (\R-0.28,-0.48)
				-- (\R+0.28,-0.48) -- cycle;
				\draw[thick] (\R-0.18,-0.55) circle (0.08);
				\draw[thick] (\R+0.18,-0.55) circle (0.08);
				\draw[thick] (\R-0.38,-0.65) -- (\R+0.38,-0.65);
				\foreach \x in {-0.30,-0.15,0,0.15,0.30}{
					\draw[thin] (\R+\x,-0.65) -- (\R+\x-0.12,-0.83);
				}
				
				% anellino in B (perno tra corpo e carrello)
				\node[circle, fill=white, draw=black, inner sep=0pt,
				minimum size=5pt, line width=0.8pt] at (\R,0) {};
				
				
				% cerniera interna in C
				\node[circle, fill=white, draw=black, inner sep=0pt,
				minimum size=5pt, line width=0.8pt] at (0,\R) {};
				
				% etichette punti
				\node[left, font=\scriptsize] at (-\R-0.22,+0.15)
				{$A \!\equiv\! O_1$};
				\node[left, font=\scriptsize] at (\R,+0.20)
				{$B \!\equiv\! O_2$};
				\node[above, font=\scriptsize] at (0,\R+0.10) {$C$};
				
				% etichette corpi
				\node[font=\small] at ({(\R+0.35)*cos(150)},{(\R+0.35)*sin(150)}) {$\mathcal{C}_1$};
				\node[font=\small] at ({(\R+0.35)*cos(30)},{(\R+0.35)*sin(30)})	{$\mathcal{C}_2$};
				
				% coppia mu antioraria su C_1 (centro asse C_1, a 135°)
				\node[circle, fill=black, inner sep=0pt, minimum size=2pt]
				at ({\R*cos(135)},{\R*sin(135)}) {};
				\draw[->, thick]
				({\R*cos(135)+0.30},{\R*sin(135)+0.05})
				arc[start angle=10, end angle=310, radius=0.30];
				\node[font=\scriptsize]
				at ({\R*cos(135)+0.55},{\R*sin(135)+0.10}) {$\mu$};
				
				% forza F su C_2 a distanza orizzontale R/2 da B
				% P = (R - R/2, R*sin(60)) = (R/2, R*sqrt(3)/2)
				\node[circle, fill=black, inner sep=0pt, minimum size=2pt]
				at ({\R/2},{\R*0.866}) {};
				\draw[->, thick]
				({\R/2},{\R*0.866+\arr}) -- ({\R/2},{\R*0.866});
				\node[font=\scriptsize]
				at ({\R/2-0.20},{\R*0.866+\arr-0.10}) {$F$};
				\node[font=\scriptsize]
				at ({\R/2},{\R*0.866-0.25}) {$P$};
				
				% linea tratteggiata grigia da B a P (per evidenziare la distanza R/2)
				\draw[gray!50,  thin] (\R,{\R*0.866}) -- ({\R/2},{\R*0.866});
				\draw[gray!50,  thin] (\R,0) -- (\R,{\R*0.866});
				\node[gray!70, font=\scriptsize] at ({(\R+\R/2)/2},{\R*0.866+0.18})
				{$R/2$};
				
				% etichetta pannello
				\node[font=\scriptsize] at (0,-0.50) {(b)};
				
			\end{scope}
			
		\end{tikzpicture}
		\caption{Arco a tutto sesto con incastro: poli
			$O_1 \equiv A$ e $O_2 \equiv B$.
			\emph{(a)}~Dato cinematico: cedimento
			verticale $\bar{v}_A$ verso il basso
			all'incastro in $A$ e cedimento verticale
			$\bar{v}_B$ verso l'alto al carrello in $B$.
			\emph{(b)}~Dato statico: coppia $\mu$
			antioraria su $\mathcal{C}_1$ e forza
			verticale $F$ verso il basso su
			$\mathcal{C}_2$, a distanza orizzontale
			$R/2$ da $B$.}
		\label{fig:arco_incastro_dati}
	\end{figure}
	
	\begin{svolgimento}
		
		\esottotit{Classificazione e gerarchia.}
		Il sistema ha $n_c = 2$ corpi, $n = 6$ gradi di libertà,
		$m = 6$ vincoli scalari (3 dall'incastro in $A$, 2 dalla
		cerniera in $C$, 1 dal carrello in $B$): $\nu = 0$,
		sistema isostatico.
		
		La struttura gerarchica è immediata: $\mathcal{C}_1$ ha
		l'incastro in $A$ che ne fissa completamente la
		configurazione (a meno dei cedimenti assegnati) ed è
		quindi il sottosistema \emph{trascinante} nella
		cinematica e \emph{portante} nella statica.
		$\mathcal{C}_2$ è \emph{trascinato}/\emph{portato}:
		la sua configurazione è determinata dalla cerniera in
		$C$ (che trasmette il moto di $\mathcal{C}_1$) e dal
		carrello in $B$.
		
		\esottotit{Formulazione matriciale.}
		Il vettore di configurazione è:
		\[
		\bm{u} =
		\bigl\{
		u(O_1)\,\; v(O_1)\,\; \vartheta_1\,\;
		u(O_2)\,\; v(O_2)\,\; \vartheta_2
		\bigr\}^\top.
		\]
		Le equazioni di vincolo, con $O_1\equiv A$ e
		$O_2\equiv B$, sono:
		
		\noindent\textit{Incastro in $A$} (3 equazioni):
		\begin{equation*}
			u(O_1) = 0, \quad
			v(O_1) = -\bar{v}_A, \quad
			\vartheta_1 = 0.
		\end{equation*}
		
		\noindent\textit{Cerniera interna in $C$}
		(2 equazioni, spostamento relativo nullo):
		\begin{align*}
			u(O_2) - \vartheta_2\cdot(-R) -
			\bigl[u(O_1) - \vartheta_1\cdot R\bigr] &= 0, \\
			v(O_2) + \vartheta_2 (-R) -
			\bigl[v(O_1) + \vartheta_1R\bigr] &= 0.
		\end{align*}
		Con $O_1\equiv A$: le coordinate di $C$ rispetto ad
		$O_1$ sono $(R, R)$; rispetto ad $O_2\equiv B$ sono
		$(-R, R)$.
		
		\noindent\textit{Carrello verticale in $B$}
		(1 equazione):
		\[
		v(O_2) = \bar{v}_B.
		\]
		
		La matrice dei vincoli $\bm{A}$ assume la struttura
		a blocchi triangolare inferiore della
		\FIGref{fig:blocchi_triangolari}:
		\begin{equation*}
			\bm{A} =
			\begin{bmatrix}
				\bm{A}_{11} & \bm{0} \\[4pt]
				\bm{A}_{21} & \bm{A}_{22}
			\end{bmatrix}
			=
			\left[
			\begin{array}{ccc|ccc}
				1 & 0 &  0 & 0 & 0 &  0 \\
				0 & 1 &  0 & 0 & 0 &  0 \\
				0 & 0 &  1 & 0 & 0 &  0 \\
				\hline
				-1 & 0 &  R & 1 & 0 &  R \\
				0 &-1 & -R & 0 & 1 & -R \\
				0 & 0 &  0 & 0 & 1 &  0
			\end{array}
			\right],
		\end{equation*}
		\esottotit{Soluzione del problema cinematico.}
		Per la struttura triangolare inferiore il sistema
		$\bm{A}\bm{u}=\bm{s}$ si risolve per sostituzione
		in avanti.
		
		\noindent\textit{Blocco $\mathcal{C}_1$}
		(prime tre equazioni):
		\[
		u(O_1) = 0, \qquad
		v(O_1) = -\bar{v}_A, \qquad
		\vartheta_1 = 0.
		\]
		$\mathcal{C}_1$ trasla verticalmente verso il basso
		di $\bar{v}_A$ senza ruotare.
		
		\noindent\textit{Blocco $\mathcal{C}_2$}
		(ultime tre equazioni), sostituiti i valori di
		$\mathcal{C}_1$:
		\begin{equation*}
			u(O_2) + \vartheta_2\,R = 0,\quad
			\bar{v}_A + v(O_2) - \vartheta_2\,R = 0,\quad
			v(O_2) = \bar{v}_B.
		\end{equation*}
		Dalla terza: $v(O_2) = \bar{v}_B$. Sostituendo
		nella seconda:
		\[
		\vartheta_2 = \frac{\bar{v}_A + \bar{v}_B}{R}.
		\]
		Dalla prima:
		\[
		u(O_2) = \vartheta_2\,R = \bar{v}_A + \bar{v}_B.
		\]
		La configurazione del sistema è univocamente
		determinata:
		\begin{equation*}
			\bm{u} =
			\Bigl\{
			0\;\;-\bar{v}_A\,\;0\;\;
			\bar{v}_A+\bar{v}_B\;\;\bar{v}_B\,\;
			\dfrac{\bar{v}_A+\bar{v}_B}{R}
			\Bigr\}^\top.
		\end{equation*}
		\esottotit{Interpretazione cinematica:
			metodo grafo-analitico.}
		Si applicano i due cedimenti separatamente per
		sovrapposizione. Per ciascun contributo si
		individuano i centri di rotazione assoluti $C_1$,
		$C_2$ e relativo $C_{12}$, si tracciano le rette
		ausiliarie $r_H$ e $r_V$ e si propagano le
		ampiezze con il metodo delle rette di proiezione
		(\FIGref{fig:arco_incastro_grafo}).
		
		\smallskip
		\noindent\textit{Contributo $\bar{v}_A$
			($\bar{v}_B = 0$).}
		Si sopprime la componente verticale dell'incastro
		in $A$. I vincoli rimasti su $\mathcal{C}_1$ sono
		la componente orizzontale dell'incastro e il
		bipendolo: il corpo può solo traslare verticalmente,
		con rotazione nulla. Il suo \CRa\ $C_1$ è dunque
		il punto improprio della direzione orizzontale.
		
		Il \CRr\ $C_{12}$ coincide con la cerniera in $C$
		(cedimento relativo nullo). Per il teorema di
		allineamento, il \CRa\ di $\mathcal{C}_2$ giace
		sulla retta per $C_1$ e $C_{12}$: poiché $C_1$ è
		il punto improprio della direzione orizzontale,
		questa retta è l'orizzontale per $C$. Il \CRa\
		$C_2$ giace inoltre sulla verticale per $B$
		(carrello verticale con cedimento nullo): la loro
		intersezione individua il centro $C_2^{(1)}$,
		la cui posizione è indicata in
		\FIGref{fig:arco_incastro_grafo}\,a.
		
		Si usa la retta $r_H$. Le tracce dei centri su
		$r_H$ sono: $C_1$ all'infinito (traslazione pura
		di $\mathcal{C}_1$); $C_{12}\equiv C$ e
		$C_2^{(1)}$ nelle posizioni indicate dalla figura.
		Poiché $\mathcal{C}_1$ trasla di $\bar{v}_A$ verso
		il basso, lo spostamento verticale di $C$ trasmesso
		a $\mathcal{C}_2$ è $-\bar{v}_A$. La rotazione di
		$\mathcal{C}_2$ si ricava dal braccio orizzontale
		tra $C_2^{(1)}$ e $C$:
		\[
		\vartheta_2^{(1)} =
		-\frac{\bar{v}_A}{x_C - x_{C_2^{(1)}}},
		\]
		dove $x_C - x_{C_2^{(1)}}$ è la distanza
		orizzontale orientata da $C_2^{(1)}$ a $C$,
		leggibile dalla figura (\FIGref{fig:arco_incastro_grafo}\,a):
		\[
		\vartheta_2^{(1)} = \frac{\bar{v}_A}{R}.
		\]
		
		\smallskip
		\noindent\textit{Contributo $\bar{v}_B$
			($\bar{v}_A = 0$).}
		Si sopprime il carrello in $B$. L'incastro in $A$
		è perfetto: $\mathcal{C}_1$ è fisso e la cerniera
		$C$ rimane ferma. Il \CRa\ di $\mathcal{C}_2$
		coincide con la cerniera $C$ (unico punto fisso di
		$\mathcal{C}_2$): $C_2^{(2)} \equiv C$.
		
		Si usa la retta $r_H$. La traccia di
		$C_2^{(2)}\equiv C$ su $r_H$ è nella posizione
		indicata in \FIGref{fig:arco_incastro_grafo}\,b.
		La rotazione di $\mathcal{C}_2$ si ricava dal
		braccio orizzontale tra $C$ e $B$:
		\[
		\vartheta_2^{(2)} =
		\frac{\bar{v}_B}{x_B - x_C},
		\]
		dove $x_B - x_C$ è la distanza orizzontale
		orientata da $C$ a $B$, leggibile dalla figura:
		\[
		\vartheta_2^{(2)} = \frac{\bar{v}_B}{R}.
		\]
		
		\noindent\textit{Sovrapposizione:}
		\[
		\vartheta_2 =
		\vartheta_2^{(1)} + \vartheta_2^{(2)} =
		\frac{\bar{v}_A + \bar{v}_B}{R},
		\]
		in accordo con la soluzione analitica.
		
		
		% ============================================================
		\begin{figure}[!ht]
			\centering
			\resizebox{\textwidth}{!}{%
				\begin{tikzpicture}[>=Stealth, line width=0.8pt]
					
					\def\R{2.0}
					\def\vA{0.40}                        % cedimento amplificato per chiarezza
					\pgfmathsetmacro\thetaB{atan(\vA/\R)} % theta_2 in gradi (~11.3°)
					\def\arr{0.55}
					\def\xV{-3.2}                        % ascissa di r_V
					\def\yH{-1.4}                        % ordinata di r_H
					\def\vB{0.40}
					\pgfmathsetmacro\thetaC{atan(\vB/\R)}
					
					% ============================================================
					% PANNELLO (a) — contributo v_A
					% ============================================================
					\begin{scope}[xshift=0cm]
						
						% ----- Sistema originale (grigio chiaro, riferimento) -----
						\draw[gray!50, thick] (-\R,0)
						arc[start angle=180, end angle=90, radius=\R];
						\draw[gray!50, thick] (0,\R)
						arc[start angle=90, end angle=0, radius=\R];
						
						% incastro in A (suolo orizzontale + tratteggio obliquo verso il basso)
						\draw[thick, gray!80] (-\R-0.38,0) -- (-\R+0.38,0);
						\foreach \x in {-0.28,-0.14,0,0.14,0.28}{\draw[thin, gray!80] (-\R+\x,0) -- (-\R+\x-0.13,-0.20);
						}
						
						% carrello verticale in B (configurazione indeformata, in grigio)
						\draw[gray!50, thick] (\R,0) -- (\R-0.28,-0.48)
						-- (\R+0.28,-0.48) -- cycle;
						\draw[gray!50, thick] (\R-0.18,-0.55) circle (0.08);
						\draw[gray!50, thick] (\R+0.18,-0.55) circle (0.08);
						\draw[gray!50, thick] (\R-0.38,-0.65) -- (\R+0.38,-0.65);
						\foreach \x in {-0.30,-0.15,0,0.15,0.30}{
							\draw[gray!50, thin] (\R+\x,-0.65) -- (\R+\x-0.12,-0.83);
						}
						\node[circle, fill=white, draw=gray!50, inner sep=0pt,
						minimum size=5pt, line width=0.8pt] at (\R,0) {};
						
						% cerniera interna in C (sistema originale - perno grigio)
						\node[circle, fill=white, draw=gray!50, inner sep=0pt,
						minimum size=5pt, line width=0.8pt] at (0,\R) {};
						
						% ----- Sistema deformato (nero spesso) -----
						% C_1 deformato: tutto traslato verso il basso di v_A
						\begin{scope}[yshift=-\vA cm]
							\draw[line width=1.5pt] (-\R,0)
							arc[start angle=180, end angle=90, radius=\R];
							% incastro deformato in A'
							\draw[thick] (-\R-0.38,0) -- (-\R+0.38,0);
							\foreach \x in {-0.28,-0.14,0,0.14,0.28}{
								\draw[thin] (-\R+\x,0) -- (-\R+\x-0.13,-0.20);
							}
						\end{scope}
						
						% perno B' (rotazione di B attorno a C_2^(1))
						\node[circle, fill=white, draw=black, inner sep=0pt,
						minimum size=5pt, line width=0.8pt] at (\R+\vA,0) {};
						
						% C_2 deformato: ruotato di theta_2 attorno a C_2^(1) = (R, R)
						\begin{scope}[rotate around={\thetaB:(\R,\R)}]
							\draw[line width=1.5pt] (0,\R)
							arc[start angle=90, end angle=0, radius=\R];
						\end{scope}
						
						% perno C' (nuova posizione di C, comune a C_1 e C_2 deformati)
						\node[circle, fill=white, draw=black, inner sep=0pt,
						minimum size=5pt, line width=0.8pt] at (0,\R-\vA) {};
						
						% ----- Centri di rotazione -----
						% C_2^(1) = (R, R) — centro assoluto di C_2
						\fill (\R,\R) circle (1.5pt);
						\node[above right=-1pt and 1pt, font=\small] at (\R,\R)
						{$C_2^{(1)}$};

						\node[above=4pt, font=\small] at (0,\R)	{$C_{12}\!\equiv\!C$};
						
						% C_1 = punto improprio di x: freccia che esce dal disegno
						\draw[->, thick] (-\R-0.85,\R+0.55) -- (-\R-0.15,\R+0.55);
						\node[font=\small] at (-\R-0.30,\R+0.85) {$C_1\!\to\!\infty$};
						
						% retta orizzontale per C (passante per C_1 improprio e C_2^(1))
						\draw[gray!60, dashed, thin] (-\R-0.30,\R) -- (\R+0.40,\R);
						% retta verticale per B (passante per C_2^(1))
						\draw[gray!60, dashed, thin] (\R,\R+0.40) -- (\R,-0.20);
						
						% ----- Etichette punti -----
						\node[gray!80, left, font=\small] at (-\R,+0.20) {$A$};
						\node[gray!80, left, font=\small] at (\R,0) {$B$};
						\node[gray!80, font=\small] at (-0.20,\R-0.20) {$C$};
						\node[left=2pt, font=\small] at (-\R,-\vA) {$A'$};
						\node[below right=-2pt and 2pt, font=\small] at (0,\R-\vA) {$C'$};
						\node[above right=-2pt and 2pt, font=\small] at (\R+\vA,0) {$B'$};
						
						% ----- Etichette corpi -----
						\node[gray!80, font=\small] at ({(\R+0.45)*cos(150)},{(\R+0.45)*sin(150)}) {$\mathcal{C}_1$};
						\node[gray!80, font=\small] at ({(\R-0.30)*cos(30)},{(\R-0.30)*sin(30)})	{$\mathcal{C}_2$};
						
						% ----- Etichette corpi -----
						\node[font=\small] at ({(\R-0.6)*cos(150)},{(\R-0.6)*sin(150)}) {$\mathcal{C}_1^{\prime}$};
						\node[font=\small] at ({(\R+0.45)*cos(30)},{(\R+0.45)*sin(30)})	{$\mathcal{C}_2^{\prime}$};
						
						% ============================================================
						% Retta r_V (campo u) — verticale a sinistra
						% ============================================================
						\draw[gray!40, thin] (\xV,-0.5) -- (\xV,\R+0.6)
						node[above, font=\small] {$r_V$};
						\node[gray!60, left=2pt, font=\small] at (\xV,\R*0.5) {$u_1^{(1)}$};
						
						% skyline del triangolo (verso destra: B si sposta di +v_A)
						\draw[gray!60, thin] ({\xV+\vA},0) -- (\xV,\R);
						
						% frecce discrete in 5 punti (k=0..4)
						\foreach \k in {0,1,2,3,4}{
							\pgfmathsetmacro{\yy}{\k*\R/4}
							\pgfmathsetmacro{\uu}{\vA*(1-\k/4)}
							\fill[gray!60] (\xV,\yy) circle (1pt);
							\ifnum\k<4
							\draw[->, gray!70, thin]
							(\xV,\yy) -- ({\xV+\uu},\yy);
							\fi
						}
						% freccia thick + etichetta a y=0 (punto B)
						\draw[->, gray!70, thick] (\xV,0) -- ({\xV+\vA},0)
						node[below, font=\small] {$\bar{v}_A$};
						
						% linee di proiezione tratteggiate dal sistema a r_V
						\draw[gray!40, dashed, thin] (0,\R) -- (\xV,\R);    % C, C_2^(1)
						\draw[gray!40, dashed, thin] (\R,0) -- (\xV,0);     % B
						
						% traccia di C_2^(1) e C su r_V (a y=R)
						\node[gray!70, left=4pt, font=\scriptsize] at (\xV,\R)
						{$C_2^{(1)}\!,\,C$};
						
						% ============================================================
						% Retta r_H (campo v) — orizzontale in basso
						% ============================================================
						\draw[gray!40, thin] (-\R-0.5,\yH) -- (\R+0.5,\yH)
						node[right, font=\small] {$r_H$};
						\node[gray!60, above=3pt, font=\small] at (-\R*0.5,\yH) {$v_1^{(1)}$};
						\node[gray!60, above=3pt, font=\small] at (\R*0.5,\yH) {$v_2^{(1)}$};
						
						% skyline: v_1 costante (-v_A) tra A e C, poi v_2 lineare a 0 in B
						\draw[gray!60, thin] (-\R,\yH-\vA) -- (0,\yH-\vA) -- (\R,\yH);
						
						% frecce discrete v_1 (5 punti tra A e C, costante = -v_A)
						\foreach \k in {0,1,2,3,4}{
							\pgfmathsetmacro{\xx}{-\R + \k*\R/4}
							\fill[gray!60] (\xx,\yH) circle (1pt);
							\draw[->, gray!70, thin] (\xx,\yH) -- (\xx,\yH-\vA);
						}
						% frecce discrete v_2 (3 punti tra C e B; salto k=0 già in v_1, k=4 ha v=0)
						\foreach \k in {1,2,3}{
							\pgfmathsetmacro{\xx}{\k*\R/4}
							\pgfmathsetmacro{\vv}{\vA*(1-\k/4)}
							\fill[gray!60] (\xx,\yH) circle (1pt);
							\draw[->, gray!70, thin] (\xx,\yH) -- (\xx,\yH-\vv);
						}
						\fill[gray!60] (\R,\yH) circle (1pt);   % B / C_2^(1)
						
						% freccia thick + etichetta a x=0 (punto C, v=-v_A)
						\draw[->, gray!70, thick] (-\R,\yH) -- (-\R,\yH-\vA)
						node[left, font=\small] {$\bar{v}_A$};
						
						% proiezioni tratteggiate sui punti notevoli
						\draw[gray!40, dashed, thin] (-\R,-0.20) -- (-\R,\yH);
						\draw[gray!40, dashed, thin] (0,\R-\vA-0.10) -- (0,\yH);
						\draw[gray!40, dashed, thin] (\R,-0.20) -- (\R,\yH);
						
						% tracce dei centri sulla r_H
						\node[gray!70, above, font=\scriptsize] at (0,\yH)	{$C_{12}\!\equiv\!C$};
						\node[gray!70, above, font=\scriptsize] at (\R,\yH)	{$C_2^{(1)}$};
						
						% etichetta pannello
						\node[font=\small] at (-\R,\yH-1.0) {(a)};
						
					\end{scope}
					
					% ============================================================
					% PANNELLO (b) — contributo v_B (da costruire)
					% ============================================================
					
					\begin{scope}[xshift=8.0cm]
						
						% ----- Sistema originale: solo C_2 in grigio (sarà ruotato) -----
						\draw[gray!50, thick] (0,\R)
						arc[start angle=90, end angle=0, radius=\R];
						
						% incastro in A (spettatore, GRIGIO)
						\draw[gray!50, thick] (-\R-0.38,0) -- (-\R+0.38,0);
						\foreach \x in {-0.28,-0.14,0,0.14,0.28}{
							\draw[gray!50, thin] (-\R+\x,0) -- (-\R+\x-0.13,-0.20);
						}
						
						% carrello verticale in B (configurazione indeformata, in grigio)
						\draw[gray!50, thick] (\R,0) -- (\R-0.28,-0.48)
						-- (\R+0.28,-0.48) -- cycle;
						\draw[gray!50, thick] (\R-0.18,-0.55) circle (0.08);
						\draw[gray!50, thick] (\R+0.18,-0.55) circle (0.08);
						\draw[gray!50, thick] (\R-0.38,-0.65) -- (\R+0.38,-0.65);
						\foreach \x in {-0.30,-0.15,0,0.15,0.30}{
							\draw[gray!50, thin] (\R+\x,-0.65) -- (\R+\x-0.12,-0.83);
						}
						\node[circle, fill=white, draw=gray!50, inner sep=0pt,
						minimum size=5pt, line width=0.8pt] at (\R,0) {};
						
						% ----- C_1 fisso (nero spesso, una sola volta) -----
						\draw[line width=1.5pt] (-\R,0)
						arc[start angle=180, end angle=90, radius=\R];
						
						% ----- C_2 deformato (ruotato di theta attorno a C) -----
						\begin{scope}[rotate around={\thetaC:(0,\R)}]
							\draw[line width=1.5pt] (0,\R)
							arc[start angle=90, end angle=0, radius=\R];
						\end{scope}
						
						% perno C (fisso, comune a C_1 e C_2 deformato)
						\node[circle, fill=white, draw=black, inner sep=0pt, minimum size=5pt, line width=0.8pt] at (0,\R) {};
						
						% perno B' = (R + v_B, v_B) — nuova posizione di B
						\node[circle, fill=white, draw=black, inner sep=0pt,
						minimum size=5pt, line width=0.8pt] at (\R+\vB,\vB) {};
						
						% ----- Centro di rotazione C_2^(2) ≡ C -----
						%\fill (0,\R) circle (1.5pt);
						\node[above=4pt, font=\small] at (0,\R) {$C_2^{(2)}\!\equiv\!C$};
						
						% ----- Etichette punti -----
						\node[gray!80, left, font=\small] at (-\R,0.20) {$A$};
						\node[gray!80,left, font=\small] at (\R,0.20) {$B$};
						\node[above right=-2pt and 2pt, font=\small] at (\R+\vB,\vB) {$B'$};
						
						% ----- Etichette corpi -----
						\node[gray!80, font=\small] at ({(\R+0.45)*cos(150)},{(\R+0.45)*sin(150)}) {$\mathcal{C}_1$};
						\node[gray!80, font=\small] at ({(\R-0.30)*cos(30)},{(\R-0.30)*sin(30)})	{$\mathcal{C}_2$};
						\node[font=\small] at ({(\R+0.65)*cos(30)},{(\R+0.65)*sin(30)})	{$\mathcal{C}_2^{\prime}$};
						
						% ============================================================
						% Retta r_V (campo u) — verticale a sinistra
						% ============================================================
						\draw[gray!40, thin] (\xV,-0.5) -- (\xV,\R+0.6)
						node[above, font=\small] {$r_V$};
						\node[gray!60, left=2pt, font=\small] at (\xV,\R*0.5) {$u_2^{(2)}$};
						
						% skyline del triangolo (verso destra: B si sposta di +v_B)
						\draw[gray!60, thin] ({\xV+\vB},0) -- (\xV,\R);
						
						% frecce discrete in 5 punti
						\foreach \k in {0,1,2,3,4}{
							\pgfmathsetmacro{\yy}{\k*\R/4}
							\pgfmathsetmacro{\uu}{\vB*(1-\k/4)}
							\fill[gray!60] (\xV,\yy) circle (1pt);
							\ifnum\k<4
							\draw[->, gray!70, thin]
							(\xV,\yy) -- ({\xV+\uu},\yy);
							\fi
						}
						% freccia thick + etichetta a y=0 (punto B)
						\draw[->, gray!70, thick] (\xV,0) -- ({\xV+\vB},0)
						node[below, font=\small] {$\bar{v}_B$};
						
						% linee di proiezione tratteggiate
						\draw[gray!40, dashed, thin] (0,\R) -- (\xV,\R);    % C / C_2^(2)
						\draw[gray!40, dashed, thin] (\R,0) -- (\xV,0);     % B
						
						% traccia di C_2^(2) ≡ C su r_V (a y=R)
						\node[gray!70, left=4pt, font=\scriptsize] at (\xV,\R)
						{$C_2^{(2)}\!\equiv\!C$};
						
						% ============================================================
						% Retta r_H (campo v) — orizzontale in basso
						% ============================================================
						\draw[gray!40, thin] (-\R-0.5,\yH) -- (\R+0.5,\yH)
						node[right, font=\small] {$r_H$};
						\node[gray!60, below=3pt, font=\small] at (-\R*0.5,\yH) {$v_1^{(2)}$};
						\node[gray!60, below=3pt, font=\small] at (\R*0.5,\yH) {$v_2^{(2)}$};
						
						% skyline: v_2 lineare da 0 in x=0 a v_B in x=R (frecce SOPRA r_H)
						\draw[gray!60, thin] (0,\yH) -- (\R,\yH+\vB);
						
						% frecce discrete v_2 (verso ALTO)
						\foreach \k in {0,1,2,3,4}{
							\pgfmathsetmacro{\xx}{\k*\R/4}
							\pgfmathsetmacro{\vv}{\vB*\k/4}
							\fill[gray!60] (\xx,\yH) circle (1pt);
							\ifnum\k>0
							\draw[->, gray!70, thin] (\xx,\yH) -- (\xx,\yH+\vv);
							\fi
						}
						% freccia thick + etichetta a x=R (punto B, dove v=v_B)
						\draw[->, gray!70, thick] (\R,\yH) -- (\R,\yH+\vB)
						node[right, font=\small] {$\bar{v}_B$};
						
						% proiezioni tratteggiate
						\draw[gray!40, dashed, thin] (0,\R-0.10) -- (0,\yH);
						\draw[gray!40, dashed, thin] (\R,-0.20) -- (\R,\yH);
						
						% traccia di C_2^(2) ≡ C su r_H (a x=0)
						\node[gray!70, above=4pt, font=\scriptsize] at (0,\yH)
						{$C_2^{(2)}\!\equiv\!C$};
						
						% etichetta pannello
						\node[font=\small] at (-\R,\yH-1.0) {(b)};
						
					\end{scope}
				\end{tikzpicture}
			}
			\caption{Arco a tutto sesto con incastro: contributi separati
				dei due cedimenti al metodo grafo-analitico.
				\emph{(a)}~Contributo $\bar{v}_A$: $\mathcal{C}_1$
				trasla verso il basso ($C_1$ punto improprio di $x$);
				$\mathcal{C}_2$ ruota attorno a $C_2^{(1)}$
				(intersezione dell'orizzontale per $C$ con la
				verticale per $B$); su $r_H$ sono riportate le
				tracce di $C_{12}\equiv C$ e $C_2^{(1)}$ e i
				campi $u$ e $v$ di $\mathcal{C}_1$ e $\mathcal{C}_2$.
				\emph{(b)}~Contributo $\bar{v}_B$: $\mathcal{C}_1$
				fisso, $C$ fisso; $\mathcal{C}_2$ ruota attorno a
				$C_2^{(2)}\equiv C$; su $r_H$ sono riportate la
				traccia di $C_2^{(2)}\equiv C$ e i campi $u$ e $v$
				di $\mathcal{C}_2$.}
			\label{fig:arco_incastro_grafo}
		\end{figure}
		
		
		
		\esottotit{Soluzione del problema statico.}
		Il vettore delle forze generalizzate, con polo
		$O_1\equiv A$ per $\mathcal{C}_1$ e polo
		$O_2\equiv B$ per $\mathcal{C}_2$, è:
		\[
		\bm{f} =
		\Bigl\{
		0\,\;0\,\;\mu\,\;
		0\,\;-F\,\;\tfrac{FR}{2}
		\Bigr\}^\top.
		\]
		Il momento di $F$ rispetto a $B$ è $FR/2$
		antiorario: la forza è verticale verso il basso
		a distanza orizzontale $R/2$ a sinistra di $B$,
		quindi $\mu(B) = F\cdot R/2 > 0$.
		
		Il sistema $\bm{A}^\top\bm{r}=-\bm{f}$ ha
		struttura triangolare superiore a blocchi e si
		risolve per sostituzione all'indietro.
		
		\esottotit{Interpretazione statica:	corpi liberi in cascata.}
		
		\noindent\textit{Corpo libero di $\mathcal{C}_2$
			(portato).}
		Le forze su $\mathcal{C}_2$ sono: reazioni interne
		$X_C$, $Y_C$ in $C$, reazione $Y_B$ al carrello
		in $B$, forza $F$ verso il basso a distanza
		orizzontale $R/2$ da $B$.
		
		Momento rispetto a $C$ (elimina $X_C$ e $Y_C$):
		\[
		\sum_{(2)}\mu_C = 0:\quad
		Y_B\cdot R - F\cdot\frac{R}{2} = 0
		\;\Longrightarrow\;
		Y_B = \frac{F}{2}
		\quad\text{(verso l'alto)}.
		\]
		Equilibrio delle forze:
		\[
		\sum_{(2)} X = 0:\quad X_C = 0,
		\]
		\[
		\sum_{(2)} Y = 0:\quad
		Y_C = F - Y_B = \frac{F}{2}
		\quad\text{(verso l'alto su $\mathcal{C}_2$)}.
		\]
		
		\noindent\textit{Corpo libero di $\mathcal{C}_1$
			(portante).}
		Le forze su $\mathcal{C}_1$ sono: reazioni
		$X_A$, $Y_A$, $\mu(A)$ all'incastro in $A$,
		reazioni interne trasmesse da $\mathcal{C}_2$
		in $C$ uguali e opposte ($X_C=0$, $-Y_C=-F/2$),
		coppia $\mu$ antioraria.
		
		Equilibrio delle forze:
		\[
		\sum_{(1)} X = 0:\quad X_A = 0,
		\]
		\[
		\sum_{(1)} Y = 0:\quad
		Y_A - \frac{F}{2} = 0
		\;\Longrightarrow\;
		Y_A = \frac{F}{2}
		\quad\text{(verso l'alto)}.
		\]
		Momento rispetto a $A$ (le coordinate di $C$
		rispetto ad $A$ sono $(R,R)$):
		\[
		\sum_{(1)}\mu(A) = 0:\quad
		\mu(A) + \mu
		- \frac{F}{2}\cdot R = 0
		\;\Longrightarrow\;
		\mu(A) = \frac{FR}{2} - \mu.
		\]
		
		
		% ============================================================
		\begin{figure}[ht]
			\centering
			\begin{tikzpicture}[>=Stealth, line width=0.8pt]
				
				\def\R{2.0}
				\def\arr{0.55}
				
				% ============================================================
				% PANNELLO (a) — corpo libero di C_2
				% ============================================================
				\begin{scope}[xshift=0cm]
					
					% semiarco C_2
					\draw[line width=1.5pt] (0,\R)
					arc[start angle=90, end angle=0, radius=\R];
					
					% carrello verticale in B (grigio: vincolo rotto)
					\draw[gray!50, thick] (\R,0) -- (\R-0.28,-0.48)
					-- (\R+0.28,-0.48) -- cycle;
					\draw[gray!50, thick] (\R-0.18,-0.55) circle (0.08);
					\draw[gray!50, thick] (\R+0.18,-0.55) circle (0.08);
					\draw[gray!50, thick] (\R-0.38,-0.65) -- (\R+0.38,-0.65);
					\foreach \x in {-0.30,-0.15,0,0.15,0.30}{
						\draw[gray!50, thin] (\R+\x,-0.65) -- (\R+\x-0.12,-0.83);
					}
					
					% perni B e C
					\node[circle, fill=white, draw=black, inner sep=0pt,
					minimum size=5pt, line width=0.8pt] at (\R,0) {};
					\node[circle, fill=white, draw=black, inner sep=0pt,
					minimum size=5pt, line width=0.8pt] at (0,\R) {};
					
					% etichette punti
					\node[left=2pt, font=\scriptsize] at (\R+0.05,0) {$B$};
					\node[left=2pt, font=\scriptsize] at (0,\R) {$C$};
					
					% etichetta corpo
					\node[font=\small] at ({(\R-0.45)*cos(35)},{(\R-0.45)*sin(35)})	{$\mathcal{C}_2$};
					
					% ----- Reazioni in B -----
					\draw[->, thick] (\R+0.05,0) -- (\R+0.05,\arr)	node[right, font=\scriptsize] {$Y_B$};
					
					% ----- Reazioni in C (interne) -----
					% X_C = 0: annotazione in scriptsize
					\node[below=1pt, font=\scriptsize] at (0,\R)
					{$X_C\!=\!0$};
					% Y_C verticale verso l'alto, decentrata a sinistra
					\draw[->, thick] (-0.0,\R) -- (-0.0,\R+\arr)
					node[left, yshift=-3pt, font=\scriptsize] {$Y_C$};
					
					% ----- Forza F in P = (R/2, R*sqrt(3)/2) -----
					\draw[->, thick]
					({\R/2},{\R*0.866+\arr}) -- ({\R/2},{\R*0.866});
					\node[font=\scriptsize]
					at ({\R/2+0.20},{\R*0.866+\arr-0.10}) {$F$};
					
					% linea tratteggiata grigia per R/2
					\draw[gray!50, dashed, thin] (\R,{\R*0.866}) -- ({\R/2},{\R*0.866});
					\draw[gray!50, dashed, thin] (\R,0) -- (\R,{\R*0.866});
					\node[gray!70, font=\scriptsize] at ({(\R+\R/2)/2},{\R*0.866+0.18})
					{$R/2$};
					
					% etichetta pannello
					\node[font=\scriptsize] at (\R/2-0.50,-0.80) {(a)};
					
				\end{scope}
				
				% ============================================================
				% PANNELLO (b) — corpo libero di C_1
				% ============================================================
				\begin{scope}[xshift=7.0cm]
					
					% semiarco C_1
					\draw[line width=1.5pt] (-\R,0)	arc[start angle=180, end angle=90, radius=\R];
					
					% incastro in A (grigio: vincolo rotto)
					\draw[gray!50, thick] (-\R-0.38,0) -- (-\R+0.38,0);
					\foreach \x in {-0.28,-0.14,0,0.14,0.28}{
						\draw[gray!50, thin] (-\R+\x,0) -- (-\R+\x-0.13,-0.20);
					}
					
					% perno  C
					
					\node[circle, fill=white, draw=black, inner sep=0pt,
					minimum size=5pt, line width=0.8pt] at (0,\R) {};
					
					% etichette punti
					\node[left=2pt, font=\scriptsize] at (-\R,0.2) {$A$};
					\node[right=2pt, font=\scriptsize] at (0,\R) {$C$};
					
					% etichetta corpo
					\node[font=\small] at ({(\R+0.45)*cos(145)},{(\R+0.45)*sin(145)}) {$\mathcal{C}_1$};
					
					% ----- Reazioni in A -----
					% X_A = 0: annotazione in scriptsize
					\node[above right=4pt and 4pt, font=\scriptsize] at (-\R,0)
					{$X_A\!=\!0$};
					% Y_A verticale verso l'alto, allineata con A
					\draw[->, thick] (-\R-0.05,0) -- (-\R-0.05,\arr)	node[left, font=\scriptsize] {$Y_A$};
					% mu(A) semicerchio antiorario con convessità verso il basso (180°)
					\draw[<-, thick]
					({-\R+0.30},0) arc[start angle=0, end angle=-180, radius=0.30];
					\node[font=\scriptsize] at (-\R,-0.55) {$\mu_A$};
					
					
					% ----- Reazioni interne in C (uguali e opposte al pannello a) -----
					% X_C = 0: annotazione in scriptsize
					\node[above left=2pt and 4pt, font=\scriptsize] at (0,\R)
					{$X_C\!=\!0$};
					% Y_C verticale verso il basso (opposta al pannello a)
					\draw[->, thick] (0.05,\R) -- (0.05,\R-\arr)
					node[right, yshift=3pt, font=\scriptsize] {$Y_C$};
					
					% ----- Coppia esterna applicata mu in P_1 -----
					\draw[->, thick]
					({\R*cos(135)+0.30},{\R*sin(135)+0.05})
					arc[start angle=10, end angle=310, radius=0.30];
					\node[font=\scriptsize]
					at ({\R*cos(135)+0.55},{\R*sin(135)+0.10}) {$\mu$};
					
					% etichetta pannello
					\node[font=\scriptsize] at (\R/2-0.80,-0.80) {(b)};
					
				\end{scope}
				
			\end{tikzpicture}
			\caption{Arco a tutto sesto con incastro: schemi dei corpi	liberi di $\mathcal{C}_2$ \emph{(a)} e		$\mathcal{C}_1$ \emph{(b)}. La cascata statica	procede da $\mathcal{C}_2$ a $\mathcal{C}_1$,	in ordine inverso rispetto alla cinematica;	polo di momento in $C$ per $\mathcal{C}_2$,	in $A$ per $\mathcal{C}_1$.}
			\label{fig:arco_incastro_cl}
		\end{figure}
		
	\end{svolgimento}
\end{esempio}

